\documentclass[11pt]{article}
\usepackage[T1]{fontenc}
\usepackage{lmodern}
\usepackage[margin=1in]{geometry}
\usepackage{amsmath,amssymb,amsthm,mathtools,mathrsfs}
\usepackage{microtype}
\usepackage{graphicx,xcolor}
\definecolor{linkblue}{RGB}{30,55,100}
\usepackage[colorlinks=true,linkcolor=linkblue,citecolor=linkblue,urlcolor=linkblue]{hyperref}
\hypersetup{pdftitle={The rational Hodge conjecture for products of K3 surfaces},pdfauthor={OpenAI},bookmarksdepth=2}
\setlength{\emergencystretch}{2em}
\numberwithin{equation}{section}
\newtheorem{theorem}{Theorem}[section]
\newtheorem{proposition}[theorem]{Proposition}
\newtheorem{lemma}[theorem]{Lemma}
\newtheorem{corollary}[theorem]{Corollary}
\theoremstyle{definition}
\newtheorem{definition}[theorem]{Definition}
\theoremstyle{remark}
\newtheorem{remark}[theorem]{Remark}
\newcommand{\Q}{\mathbb Q}
\newcommand{\R}{\mathbb R}
\newcommand{\C}{\mathbb C}
\newcommand{\Z}{\mathbb Z}
\newcommand{\cO}{\mathcal O}
\DeclareMathOperator{\HH}{HH}
\DeclareMathOperator{\Ext}{Ext}
\DeclareMathOperator{\Hom}{Hom}
\DeclareMathOperator{\End}{End}
\DeclareMathOperator{\Perf}{Perf}
\DeclareMathOperator{\PD}{PD}
\DeclareMathOperator{\Spin}{Spin}
\DeclareMathOperator{\SU}{SU}
\DeclareMathOperator{\SO}{SO}
\DeclareMathOperator{\GL}{GL}
\DeclareMathOperator{\SL}{SL}
\DeclareMathOperator{\NS}{NS}
\DeclareMathOperator{\Hg}{Hg}
\DeclareMathOperator{\CH}{CH}
\DeclareMathOperator{\cl}{cl}
\DeclareMathOperator{\Res}{Res}
\DeclareMathOperator{\Tr}{Tr}
\DeclareMathOperator{\str}{str}
\DeclareMathOperator{\Sym}{Sym}
\DeclareMathOperator{\id}{id}
\DeclareMathOperator{\rk}{rk}
\DeclareMathOperator{\td}{td}
\DeclareMathOperator{\ch}{ch}
\DeclareMathOperator{\Pic}{Pic}
\DeclareMathOperator{\im}{im}
\DeclareMathOperator{\Spec}{Spec}
\DeclareMathOperator{\HF}{HF}
\newcommand{\Ainf}{A_\infty}
\newcommand{\KS}{\mathrm{KS}}
\newcommand{\Hdg}{\mathrm{Hdg}}
\title{The rational Hodge conjecture for products of K3 surfaces}
\author{OpenAI}
\date{October 4, 2026}

\begin{document}
\maketitle
\begin{abstract}
We prove the rational Hodge conjecture for every finite product of projective
complex K3 surfaces: every rational Hodge class is algebraic.  The factors
may be distinct or repeated, with no restrictions on their Picard numbers,
periods, or endomorphism fields.
\end{abstract}
\tableofcontents

\section{Introduction}\label{sec:introduction}

For a smooth projective complex variety \(X\), an algebraic subvariety
of codimension \(p\) has a rational cohomology class of Hodge type
\((p,p)\).  The rational Hodge conjecture asks whether these classes
span every rational class of that type.  Write
\[
 \Hdg^p(X)=H^{2p}(X,\Q)\cap H^{p,p}(X,\C),
\]
where the intersection is taken inside \(H^{2p}(X,\C)\), and write
\(\CH^p(X)_\Q\) for codimension-\(p\) cycles modulo rational
equivalence with rational coefficients.

\begin{theorem}\label{thm:main}
Let \(m\ge1\), and let \(S_1,\ldots,S_m\) be any projective complex
K3 surfaces.  For \(X=S_1\times\cdots\times S_m\) and every
\(0\le p\le2m\), the cycle-class map
\[
 \cl:\CH^p(X)_\Q\longrightarrow\Hdg^p(X)
\]
is surjective.
\end{theorem}

The factors may repeat, and no relation between their periods or
endomorphism fields is assumed.  The distinction matters: a product
can contain Hodge tensors coupling different transcendental spaces,
while a power of one surface can contain tensors beyond those generated
by pairwise contractions.  The proof treats both phenomena.

For a projective K3 surface \(S\), its transcendental space is
\[
 T(S)=\NS(S)_\Q^\perp\subset H^2(S,\Q),
\]
with the intersection form.  

\subsection{History and the remaining tensor problem}
The Hodge conjecture, formulated in Hodge's 1950 Congress address
\cite{Hodge1952}, asks which rational cohomology classes come from algebraic
cycles.  For a projective K3 surface itself, the Lefschetz $(1,1)$ theorem,
together with the classes of the surface and of a point, answers the question.
Products introduce a different problem: their Hodge classes include tensors
in the transcendental cohomologies of the factors.  Already for two factors,
the nontrivial part concerns Hodge homomorphisms between these cohomologies;
on higher powers, further invariant tensors occur.  We use the customary
background on K3 surfaces, their periods, and their transcendental Hodge
structures as presented in \cite{HuybrechtsK3}.

Kuga and Satake associated an abelian variety to a polarized Hodge structure
of K3 type \cite{KugaSatake}.  Its first cohomology carries a Clifford-algebra
representation, and the original weight-two Hodge structure embeds into a
tensor square of this weight-one structure.  Deligne developed this
construction in families and used it in his proof of the Weil conjecture
for K3 surfaces \cite[Sections~5--6]{DeligneK3}.  Andr\'e subsequently proved
that the Kuga--Satake correspondence is motivated
\cite[Section~7.1]{AndreMotives}.  These results provide the Hodge-theoretic
and motivic relation to abelian varieties; obtaining an algebraic cycle
representing the correspondence is the additional requirement in the
Kuga--Satake Hodge conjecture.  Van Geemen's account explains the Clifford
construction and its relation to exceptional Hodge classes on abelian
varieties \cite{VanGeemenKS}.

A second line of work constructs correspondences between K3 surfaces using
moduli spaces of sheaves.  Mukai's work supplied the fundamental examples
\cite{MukaiBundles}, and Nikulin extended the isometry result using
quadratic-lattice methods \cite{NikulinCorrespondences}.  Buskin proved that every rational Hodge isometry between
the second cohomologies of two projective K3 surfaces is algebraic
\cite[Theorem~1.1]{Buskin}; Huybrechts gave a proof using twisted derived
categories and established an isomorphism of their Chow motives
\cite[Theorems~0.1--0.2]{HuybrechtsIsogenous}.  The isometry hypothesis is
significant: arbitrary Hodge homomorphisms need not preserve the quadratic
forms.  Varesco proved that algebraic Kuga--Satake correspondences also make
Hodge similarities algebraic \cite[Theorem~4.5]{VarescoSimilarities}.

Zarhin determined the Hodge groups of transcendental K3 cohomology in terms
of its endomorphism field, which is either totally real or CM \cite{Zarhin}.
This description explains why a result about squares need not settle the
problem for powers.  In the totally real case, orthogonal invariant theory
produces determinant tensors in addition to the quadratic contractions.
Varesco identified the corresponding relative determinant classes and
proved a precise generation statement
\cite[Theorem~2.16]{VarescoPowers}.  When the endomorphism field is CM, the
invariants are generated in degree two; the work of Ram\'on Mar\'i, together
with algebraicity of Hodge isometries, gives the Hodge conjecture for all
powers \cite{RamonMari,Buskin}; see also
\cite[Theorem~2.2 and Corollary~2.3]{VarescoPowers}.  Here a CM endomorphism
field is the relevant condition; it need not make the whole transcendental
Hodge structure of CM type.

Algebraic Kuga--Satake correspondences have also been constructed for
important special families.  Paranjape treated the K3 surfaces arising
from double covers of the plane branched along six lines \cite{Paranjape},
and Schlickewei used their Kuga--Satake varieties to establish the Hodge
conjecture for their squares \cite[Theorem~2]{Schlickewei}.  Varesco proved an
all-powers theorem for four-dimensional families of generic Picard number
16 with a specified rational transcendental quadratic space, assuming
algebraicity of the Kuga--Satake correspondences throughout the family
\cite[Theorem~0.2]{VarescoPowers}.  The quadratic-space and family
hypotheses in that statement are essential to its scope.

Recent geometric constructions considerably enlarge these cases.
Floccari proved algebraicity of the Kuga--Satake correspondence and the
Hodge conjecture for all powers of the K3 surfaces associated with varieties
of generalized Kummer type \cite[Theorem~1.2]{FloccariKummer}.
Floccari and Fu obtained related results through singular varieties of
O'Grady's six-dimensional deformation type; their construction also gives
products of powers of the abelian, K3, and hyper-K\"ahler varieties
associated with one such singular variety
\cite[Theorem~5.7 and Corollary~5.8]{FloccariFu}.  These theorems illustrate
how geometric relations with abelian varieties control exceptional tensors
on particular families.  For a general list of K3 surfaces, a further
issue is to control the joint Hodge group: algebraicity of the tensors on
each individual factor does not by itself determine the mixed invariants.



\subsection{The proof and its principal constructions}

The Kuga--Satake realization of \(T(S)\) uses the
weight-one Hodge structure on the full even Clifford algebra
\(W_S=C^+(T(S))\), the \(H^1\) of an abelian variety up to isogeny.
For a rational vector \(w\in T(S)\) of nonzero square, the standard
correspondence is the specified Hodge map
\begin{equation}\label{intro:ks-map}
 T(S)(1)\longrightarrow\End(W_S),\qquad
 v\longmapsto\bigl(x\longmapsto vxw\bigr).
\end{equation}
Here \(\Q(1)\) has weight \(-2\).  A polarization converts this into
a degree-two tensor on the square of the abelian variety.
The precise map in \eqref{intro:ks-map} is essential to the argument.

The bridge from geometry to this exact tensor is
Proposition~\ref{prop:realization}.  It starts with a rational
middle-dimensional class killed by the symplectic form on a quartic
times a large standard torus, or on the torus alone.  Under a condition
on its degree-two annihilator, it produces a perfect complex on the
ordinary mirror.
The conclusion controls two different features of the complex:
its Euler pairings with specified test objects and the dimension of
its second self-extension space.  Part~\ref{part:algebraic} first
uses this proposition to prove Theorem~\ref{thm:main};
Part~\ref{part:realization} then proves the proposition.

The force of the realization statement lies in combining its two
conclusions.  In the K3 construction, the Euler pairings initially
determine only the algebraic component of the complex's Mukai vector.
We choose that component
using Clifford multiplication and separated powers of invariant
alternating tensors.  The Hochschild action on the whole vector
factors through \(\Ext^2\), so the small extension space forces
many cohomological operators to annihilate it.  A Clifford commutant
calculation shows that so large a kernel is possible only for the
prescribed transcendental component, with its exact normalization.
Thus an upper bound on self-extensions determines a tensor that the
Euler pairings cannot see.

The same calculation attains equality in the rank bound.  Equality
makes the semiregularity trace injective, so the complex formally
deforms along its Hodge locus.  This uses the deformation
principle developed by
Buchweitz--Flenner and Pridham
\cite{BuchweitzFlennerSemiregularity,PridhamSemiregularity};
Markman's construction of Weil classes by deforming semiregular
sheaves gives a related application
\cite[Sections~7--9]{MarkmanSecant}.  We apply this principle along
the coupled K3 and Kuga--Satake period family.  Algebraization and
specialization give \eqref{intro:ks-map} on every projective K3
surface, as proved in
Theorem~\ref{thm:exact-ks}.

This use of an extension-space bound is useful beyond identifying
one tensor: it obtains the required Hochschild rank on the algebraic
mirror itself.  It requires no identification of the Hochschild
action under the mirror equivalence.  The same method, applied to a
torus alone, constructs the degree-four spin tensors of
Proposition~\ref{prop:auxiliary-spin}.

Two companion results enter the final algebraic argument.  The exact
full-even-Clifford criterion of~\cite[Theorem~1.1]{AKS} gives all
Hodge classes on self-powers once \eqref{intro:ks-map} is algebraic.
The rational Hodge theorem for CM abelian varieties
\cite[Theorem~1.1]{CM} controls the torus constituents of auxiliary
weight-one Hodge structures.  These are theorem inputs with the stated
scope; the proof below supplies the exact KS correspondences and the
additional tensors needed to combine different surfaces.

For the combination, compare the group preserving algebraic
correspondences with the joint Hodge group of the transcendental
spaces.  The self-power results identify the projections of these
groups to each factor; the remaining question is whether the
relations between different factors agree.  Algebraic morphisms
between weight-one abelian realizations
identify the shared simple factors projectively.  Central signs can
still differ, especially among the low-rank orthogonal and unitary
factors sharing \(\mathfrak{sl}_2\).  The degree-four spin tensors
give perfect pairings between local Hom spaces and transport those
signs across a finite grid of representation choices.  A graph
calculation then shows that the remaining signs come from the center
of the joint derived group.  This is the mixed-product step in
Section~\ref{sec:mixed}; self-power algebraicity alone does not
replace it.

It remains to prove the realization proposition.  We represent the
input class, after multiplication, by a spin graded Lagrangian
immersion whose second cohomology is controlled by the class's
degree-two annihilator.  Its self-intersection generators have odd
integer degrees.  A disk-moduli Stokes identity makes the immersed
object unobstructed, and its degree-two endomorphism space is then
bounded by that second cohomology.  Only the embedded test objects
are passed through ordinary mirror symmetry.  A directed comparison
and a categorical projection produce the perfect complex while
retaining the Euler pairings and the extension bound.
The geometry adapts the controlled-immersion construction
of~\cite{A8}.  Sections~\ref{sec:topology}--\ref{sec:realization}
provide the additional arguments for a K3 factor, arbitrary large
torus dimension, integer gradings, the Stokes identity, and the
categorical comparisons.

\subsection{Conventions and organization}

All cycle classes and correspondences have rational coefficients.
An algebraic tensor means a cohomology tensor represented by a rational
algebraic correspondence; no assertion about relations modulo rational
equivalence is implicit.  Hodge groups use the special, or circle,
convention and fix Tate structures.  In calculations with total
cohomology, Tate twists equalizing weights are left implicit only
when the circle action makes their meaning unambiguous.

The symplectic constructions use spin branes, zero relative-spin
background, and integer Maslov gradings.  Novikov fields and their
extensions are embedded in \(\C\) as abstract fields; these embeddings
are not analytic evaluation of a Novikov parameter.  The relevant
mirror objects are fixed before the rational input class is chosen.

In Part~\ref{part:algebraic}, Section~\ref{sec:setup} states the
realization proposition and analyzes its mirror tests.
Section~\ref{sec:cohomology} develops cap-action and deformation tools.
Sections~\ref{sec:exactks} and~\ref{sec:auxiliary} construct algebraic
spin tensors; Section~\ref{sec:mixed} uses them to prove
Theorem~\ref{thm:main}.  Part~\ref{part:realization} proves the proposition:
Sections~\ref{sec:topology} and~\ref{sec:curvature} construct the
immersion and remove its curvature; Sections~\ref{sec:comparison}
and~\ref{sec:realization} compare the categories.


\part{Algebraic tensors and the Hodge conjecture}
\label{part:algebraic}
\section{A realization statement and its mirror test objects}
\label{sec:setup}

The first step is to realize a rational middle-dimensional class by a
perfect complex whose second self-extension group is small.  The
realization will preserve its pairings with a fixed collection of
Lagrangian test objects.  We specify those objects here and then use the
realization statement to construct the algebraic tensors.  Its proof occupies
Sections~\ref{sec:topology}--\ref{sec:realization} in the second part
of the paper.

\subsection{The symplectic spaces and the realization statement}

Let $Y$ be either a smooth quartic K3 surface with the integral
symplectic form $\omega_Y$ representing its hyperplane class, or a point.
In the latter case all $Y$-factors below are omitted.  Set
\begin{equation}
 T_g=\R^{2g}/\Z^{2g},\qquad
 \omega_T=\sum_{j=1}^g q_j\wedge p_j,\qquad
 D=Y\times T_g,\qquad 2n=\dim_{\R}D.
 \label{setup:spaces}
\end{equation}
Here $q_j,p_j$ denote the coordinate one-forms; we use the same letters
for coordinates when describing graphs.  Write $\omega$ for the product
symplectic form.  Throughout this construction $g$ is even and $n\ge12$.
All branes have integer Maslov gradings, spin structures with zero
background class, and the orientations determined by their gradings.

The ordinary mirrors are the mirror quartic $B$ of $Y$ and the product
$E^g$, where $E$ is the Tate elliptic curve; set $B$ equal to a point
when $Y$ is a point.  These varieties are initially
defined over Novikov fields.  The change of parameter in quartic
homological mirror symmetry is included in the definition of $B$.
We allow further field extensions and then an abstract field embedding
into $\C$.  Thus passage to a complex variety means extension of scalars
along a field embedding, not evaluation of a Novikov series at a complex
number.  All fields used below have cardinality at most the continuum.

Choose a finite collection of graded Lagrangian spheres $\Sigma_\ell\subset Y$
which split-generates the quartic Fukaya category and whose classes span
\begin{equation}
 P_Y=\omega_Y^\perp\subset H^2(Y,\Q).
 \label{setup:primitive-space}
\end{equation}
Their existence is recalled below.  For the point use one test object,
with the unit convention.  In $T_g$ choose finitely many graphs
\begin{equation}
 G_M=\{p=Mq\},\qquad M\in\operatorname{Sym}_g(\Z),
 \label{setup:graphs}
\end{equation}
and make small Hamiltonian perturbations when transversality is required.
The graph gradings vary continuously with the symmetric graph planes;
the spin structures come from the base frame.  The distinct slopes are
totally ordered by positive-definite difference.  We include a generating
set of scalar slopes and may enlarge the list as described below.
This order on slopes is separate from any directed ordering of their
Hamiltonian copies.

Denote the mirror images of $\Sigma_\ell$ and $G_M$ by $Q_\ell$ and $N_M$.
The graph images are line bundles after one common shift.  Put
$K_{\ell,M}=\Sigma_\ell\times G_M$.  For perfect complexes we use the Euler
pairing
$\chi(A,F)=\sum_j(-1)^j\dim\Ext^j(A,F)$.

\begin{proposition}[Realization with a self-extension bound]
\label{prop:realization}
Fix the spaces and test systems above.  Let $a\in H^n(D,\Q)$ satisfy
$\omega\smile a=0$, and define
\[
 \mathcal Z=\{x\in H^2(D,\Q):x\smile a=0\}.
\]
Suppose that cup product induces an injective map
\begin{equation}
 H^1(D,\Q)\otimes\mathcal Z\longrightarrow H^3(D,\Q).
 \label{setup:wedge-injective}
\end{equation}
After a field extension and an abstract embedding into $\C$, there is
$\mathcal E\in\Perf(B\times E^g)$ such that
\begin{equation}
 \dim\Ext^2(\mathcal E,\mathcal E)
 \le b_2(D)-\dim\mathcal Z.
 \label{setup:ext-bound}
\end{equation}
There is a common nonzero $\lambda\in\Q$, independent of $\ell,M$, with
\begin{equation}
 \chi(Q_\ell\boxtimes N_M,\mathcal E)
   =\lambda\,\langle a,[K_{\ell,M}]\rangle.
 \label{setup:euler}
\end{equation}
The mirror varieties, their enhanced categories, and the images of the
fixed test systems can be chosen before $a$ is specified.
\end{proposition}

The last assertion lets us choose divisor classes and ample classes on
the geometric mirror before choosing $a$.  One may use a preliminary
complex embedding to describe their lattices and cones: the geometric
N\'eron--Severi lattice, intersection pairing, and ampleness of its
classes are unchanged by algebraically closed field extension.  Later
embeddings therefore preserve the algebraic data used in the
construction.

The proof of Proposition~\ref{prop:realization} first constructs a graded
Lagrangian immersion representing a multiple of $a$.  Its degree-two
cohomology gives the right-hand side of \eqref{setup:ext-bound}.
Unobstructedness places it in a proper cyclic category.  We then identify
the subcategory generated by the embedded test objects with
$\Perf(B\times E^g)$ and project to that subcategory.  The categorical
completion is given in Section~\ref{sec:realization}.

\subsection{Sphere tests and the algebraic Mukai space}

In this subsection assume that $Y$ is a quartic K3 surface.
Seidel's quartic theorem gives the enhanced equivalence on a generating
category and split-generation by vanishing spheres
\cite[Theorem~1.3, Corollary~9.6, and Lemma~11.4]{SeidelQuartic}.
Rescaling the Novikov parameter accommodates the normalization of
$\omega_Y$.  Sphere branes are rational: no one-cycle holonomy condition
has to be imposed.  Compatibly graded isotopic spheres determine
quasi-isomorphic objects by \cite[Lemma~8.14]{SeidelQuartic}.  These facts
apply in particular to all the Hamiltonian copies used later.

We may add vanishing spheres from a generic pencil of quartics in
$\mathbf P^3$.  Symplectic transport identifies its smooth fibers with
$Y$, and the vanishing cycles span the kernel of ambient inclusion.
Their Poincar\'e duals therefore span $P_Y$.  Equivalently one can use
a hyperplane pencil in the corresponding Veronese embedding.

For any finite or countable set of the embedded sphere copies, there is
a common compatible almost complex structure with no nonconstant
holomorphic spheres and no nonconstant holomorphic disks on any
individual copy.  For somewhere injective curves the unmarked expected
dimensions are respectively $-2$ and $-1$.  Genericity, sphere
factorization, and the disk decomposition results of Kwon--Oh and
Lazzarini \cite{KwonOh,Lazzarini} give the assertion, as used in
\cite[Section~8]{SeidelQuartic}.  The torus graphs have vanishing relative
second homotopy, so a product almost complex structure has the same
exclusion property for the embedded product tests.

Write
\begin{equation}
 P_B=H^0(B,\Q)\oplus\NS(B)_{\Q}\oplus H^4(B,\Q),
 \label{setup:algebraic-mukai}
\end{equation}
with the Mukai pairing, in which $H^0$ and $H^4$ pair with a minus sign.
Before choosing a complex embedding, this denotes the geometric divisor
space together with its degree-zero and degree-four generators.
For an object $Q$ its Mukai vector is
$v(Q)=\operatorname{ch}(Q)\sqrt{\operatorname{td}(B)}$.

\begin{lemma}
\label{setup:rho-lemma}
The space $P_B$ has dimension $21$.  Sending the Poincar\'e dual of a
sphere test to the Mukai vector of its mirror image induces a rational
isometry
\begin{equation}
 \rho:P_Y\xrightarrow{\sim}P_B.
 \label{setup:rho}
\end{equation}
\end{lemma}

\begin{proof}
The Euler pairing of two sphere branes is minus their intersection
pairing.  On their mirror objects it is minus the Mukai pairing.
Since the test classes span the nondegenerate, $21$-dimensional space
$P_Y$, their Mukai vectors span a space of dimension at least $21$.

The mirror is the geometric generic member of the one-parameter family
obtained by resolving the diagonal-group quotient in Seidel's theorem.
That family has nonconstant period.  Before taking the finite symplectic
quotient, this is visible in the quartic pencil generated by the product
of the four coordinates and their fourth powers: at the Fermat member
the infinitesimal period in the product direction is nonzero, because
$x_0x_1x_2x_3$ survives in the Jacobian ring.  The group fixes the
holomorphic two-form.  Pullback by the quotient correspondence over an
open base therefore transports this variation to the resolved family.
Its geometric generic Picard rank is at most $19$ by period theory.
The HMS parameter is transcendental, including after its formal
reparametrization, and geometric Picard rank is unchanged by
algebraically closed field extension.  Hence $\dim P_B\le21$.

We have equality.  The pairing equality now also proves that the map on
sphere classes is well-defined: a relation among their classes maps to
a vector orthogonal to a spanning set of the nondegenerate space $P_B$,
and thus to zero.  The resulting surjective map preserves the pairings,
which proves \eqref{setup:rho}.
\end{proof}

\subsection{Graph images and full faithfulness on the test systems}

We use Abouzaid's family Floer theorem in its finite-object form
\cite[Theorem~2.10 and Appendix~A]{AbouzaidFaithful}.  It gives a faithful
$A_\infty$ functor to analytic perfect complexes.  For the present test
systems a dimension comparison will give full faithfulness.

The fibration $T_g\to\R^g/\Z^g$ is nonsingular, and both its base and
total space have vanishing second homotopy.  A graph induces
$b\mapsto(b,Mb)$ on fundamental groups, so
$\pi_2(T_g,G_M)=0$.  The graphs and their Hamiltonian copies are
therefore tautologically unobstructed.  The global zero section and
the base frame trivialize the section and spin twists in family Floer
theory.  Let $\mathfrak q$ be the Tate parameter, of positive valuation.
The rigid analytic mirror has multiplicative coordinates
$z=(z_1,\ldots,z_g)$ and periods $z\mapsto\mathfrak q^b z$, and is
$(E^g)^{\mathrm{an}}$.

\begin{lemma}[Graph images]
\label{setup:graph-images}
For each finite test system, the image of $G_M$ is a line bundle up to
shift, with numerical type $\varepsilon M$ for one common sign
$\varepsilon\in\{1,-1\}$.  After fixing the variance convention, all
shifts can be made equal.  The family Floer functor is fully faithful
on the finite systems and their Hamiltonian copies.
\end{lemma}

\begin{proof}
A graph has one transverse intersection with each torus fiber.  Its
local family Floer complex is consequently a rank-one module in one
degree.  We recall the numerical calculation of
\cite[Lemma~9.1]{A8}, including its extension to non-diagonal slopes.
On the real cover of the base lift a perturbed graph as
$p=Mq+\nabla f(q)$, where $f$ is periodic.  Comparison paths along this
lift have zero character offset: the fixed lattice offset of a lifted
graph boundary cancels at its two ends, and a residual fiber loop could
not bound.  Scalar trivializations on the real cover therefore give
the period multiplier
\begin{equation}
 j_M(b,z)=c_bz^{-\varepsilon Mb},\qquad
 b\in\Z^g,\quad c_b\in\Lambda^*.
 \label{setup:graph-multiplier}
\end{equation}
Here $\Lambda$ is the Novikov coefficient field and a change of variance
by perfect duality changes all numerical signs together.

Every integral symmetric matrix is a sum of signed matrices $uu^t$,
with $u\in\Z^g$.  Products of theta lines pulled back under
$E^g\to E$, $(x_j)\mapsto\sum u_jx_j$, have precisely the prescribed
type $\varepsilon M$.  Dividing their multiplier into
\eqref{setup:graph-multiplier} leaves a scalar character, hence a line
in $\operatorname{Pic}^0(E^g)$.  Projective rigid GAGA algebraizes the
image.  For the enhanced comparison one can use the common-refinement
\v Cech models of \cite[Section~9]{A8}; their restriction and cup maps
give the same dg comparison in $g$ variables.

For distinct slopes with definite difference, Floer cohomology has
dimension $|\det(M-M')|$ and is concentrated in one degree.  The quotient
of the corresponding lines is ample or anti-ample, and its total Ext
dimension is the same determinant.  Faithfulness is thus fullness for
these pairs.  For copies of the same slope, self-Floer cohomology has
the exterior-algebra dimensions of a $g$-torus.  A nontrivial
$\operatorname{Pic}^0$ quotient would have zero cohomology, contradicting
faithfulness.  The underlying lines of the copies therefore agree, and
their self-Ext dimensions again give fullness.

Choose the arrow direction in which definite graph comparisons have
degree zero.  The numerical sign must make the corresponding line
quotients ample.  Otherwise every pair among three ordered scalar
slopes would require the same nonzero shift difference $g$, which
contradicts additivity of shift differences.  With the correct sign,
all scalar shifts agree.  Comparison with extreme scalar slopes makes
the shift of every other graph the same.  Grading transport and the
self-cohomology support $[0,g]$ give the same conclusion for copies.
This is the normalization argument of \cite[Lemma~9.2]{A8}.
\end{proof}

\subsection{Finite generation and the order of choices}

\begin{lemma}
\label{setup:graph-generation}
Let $H$ be the principal product theta line on $E^g$.  There are
integers $k_1<\cdots<k_{g+1}$ such that, for every choice
$P_j\in\operatorname{Pic}^0(E^g)$, the lines
$H^{-k_j}\otimes P_j$ split-generate $\Perf(E^g)$.
The assertion persists under algebraically closed field extension.
\end{lemma}

\begin{proof}
Choose a vector-bundle generator $V$ of $\Perf(E^g)$
\cite{BondalVanDenBergh}.  Starting with $K_0=V$, take evaluation
surjections
\[
 0\longrightarrow K_j\longrightarrow
 W_j\otimes(H^{-k_j}\otimes P_j)
 \longrightarrow K_{j-1}\longrightarrow0.
\]
Here $W_j$ is a finite-dimensional vector space.  The exponents can be
chosen successively and uniformly in all preceding $P_j$: perform each
step over the proper parameter space
$(\operatorname{Pic}^0 E^g)^j$, using its Poincar\'e lines and relative
Serre generation and vanishing.  The kernels remain vector bundles.
The length-$g+1$ sequence gives a triangle
\[
 K_{g+1}[g]\longrightarrow C_V\longrightarrow V
 \longrightarrow K_{g+1}[g+1],
\]
where $C_V$ is a complex of sums of the displayed lines.  Since
$\Ext^{g+1}(V,K_{g+1})=0$, the triangle splits and $V$ is a summand of
$C_V$.  All evaluation sequences and base-change properties were
constructed over the parameter spaces, so the argument persists after
field extension.  The original generator remains a generator after
extension, for example by the usual perfect-diagonal criterion.
This is the argument of \cite[Lemma~9.3]{A8} in dimension $g$.
\end{proof}

If the sign convention gives ample rather than anti-ample powers,
dualize the preceding construction.  Thus the scalar graph tests can
always be chosen to include generators.  Further tests may be added
while retaining definite order: the integral symmetric matrices larger
than a fixed matrix in positive-definite order are Zariski dense.
Consequently every finite-dimensional span obtained by polynomial
evaluation on symmetric matrices can be obtained by continuing the
ordered list.  Extreme scalar slopes can then be added for the shift
normalization.

We fix these finite systems and their mirror images before choosing
$a$.  Generation in both factors survives the later field extensions.
Finally, $E$ is non-CM after every complex embedding used here.  Its
$j$-invariant has negative Novikov valuation, so is transcendental over
$\Q$, whereas a CM elliptic curve has an algebraic $j$-invariant.

\section{Cohomological operators and deformation}\label{sec:cohomology}

Proposition~\ref{prop:realization} specifies Euler pairings and bounds an
\(\Ext^2\) space.  We now express both conclusions in linear algebra.
The Euler pairings determine an algebraic component of a Mukai vector;
the bound constrains the kernel of the Hochschild cap action on the
entire vector.  When that kernel has the expected dimension, the
semiregularity map is injective, and the complex formally deforms along
its Hodge locus.  Spreading and specialization then transport the
algebraicity of its Chern character.  This is the mechanism used in
Sections~\ref{sec:exactks} and~\ref{sec:auxiliary}.

\subsection{The linear dictionary}

Two gradings will enter the argument.  To choose the symplectic input
class, we use the Lefschetz action on the cohomology of \(D\).  To study
the resulting complex, we use the actual Hodge grading on the mirror.
We first describe the former as a linear model for the latter; identifying
their transcendental Hodge structures will require the later rank argument.

We use the Mukai form \(q\) on the total even cohomology of a K3
surface: it is the intersection form on \(H^2\), and
\(q(1,[\mathrm{pt}])=-1\).  Mukai vectors are
\(v(\mathcal E)=\ch(\mathcal E)\sqrt{\td}\).
When using the Hodge circle, we suppress Tate twists; invariance means
type \((p,p)\) separately in each cohomological degree.

Suppose first that \(Y\) is a quartic.  Recall the isometry
\(\rho:P_Y\longrightarrow P_B\) supplied by the sphere tests in
Section~\ref{sec:setup}, where
\[
 P_Y=\omega_Y^\perp\subset H^2(Y,\Q),\qquad
 P_B=H^0(B,\Q)\oplus\NS(B)_\Q\oplus H^4(B,\Q).
\]
The orthogonal complement of \(P_Y\) in the Mukai space is
\[
 U=\langle 1,\omega_Y,[\mathrm{pt}]\rangle.
\]
It is a split ternary quadratic space.  Lefschetz theory gives an
\(\mathfrak{sl}_2\)-action with Cartan operator \(2-\deg\) and lowering
operator \(f=\omega_Y\cup-\).  The space \(U\) is its irreducible
representation of highest weight two, while \(P_Y\) is trivial.
Set \(t_+=1\) and \(t_-=-[\mathrm{pt}]\), so
\(q(t_+,t_-)=1\), and choose a weight-zero basis vector \(t_0\) of
\(U\).  We identify \(\Spin(U)\) with \(\SL_2\).
Write \(U_B=T(B)\).  Witt cancellation in the two rational Mukai lattices
identifies \(U\), as a quadratic space, with \(U_B\).
This is a rational quadratic identification only; no assertion about
the transcendental action of mirror symmetry is used.

For the torus write
\[
 W=M_0\otimes H,\qquad \dim M_0=g,\qquad
 H=\Q e_+\oplus\Q e_-,\qquad f(e_+)=e_-.
\]
The vectors \(e_+,e_-\) are the standard \(\SL_2\)-basis, of weights
\(1,-1\).  Write \(a_i,b_i\) for the corresponding basis vectors of
\(W\), \(1\le i\le g\).  The creation--annihilation transform is the
linear isomorphism
\begin{equation}\label{coh:transform}
 \Phi:H^*(T_g,\Q)\longrightarrow\bigwedge W
\end{equation}
that sends \(p_1\cdots p_g\) to \(1\), cup product by \(q_i\) to
exterior multiplication by \(b_i\), and cup product by \(p_i\) to
contraction by \(a_i^*\).  With consistent ordering, the other two
creation--annihilation rules are
\[
 \Phi\,\iota_{q_i^*}\,\Phi^{-1}=\iota_{b_i^*},
 \qquad
 \Phi\,\iota_{p_i^*}\,\Phi^{-1}=a_i\wedge-.
\]
Together these four rules and the chosen vacuum characterize the
transform.  Thus the Cartan
operator on the source is \(g-\deg\), and cup product by
\(\omega_T\) becomes \(f\).  The transform is not a ring isomorphism.

For a symmetric integral slope matrix \(M\), the graph class is, up
to a fixed orientation sign, \(\prod_i(p_i-(Mq)_i)\).  Consequently
\begin{equation}\label{coh:graph-exponential}
 \Phi\bigl(\PD[G_M]\bigr)
 =\pm\exp\left(\sum_{i,j}M_{ij}a_i b_j\right).
\end{equation}
One may verify the identity by the annihilating operators of the graph
and its constant coefficient.  Mark \(H^1(E,\Q)\) by \(H\), and hence
\(H^1(E^g,\Q)\) by \(W\).  The numerical graph calculation in
Section~\ref{sec:setup} makes the right side the Mukai vector of the
mirror line, with the common shift understood.  The marking is oriented
or reverse oriented according to the single sign convention for all
graph lines; it is fixed throughout the construction.

The degree-two \(\SL_2\)-invariants in \(\bigwedge W\) are precisely
the expressions \(\sum M_{ij}a_i b_j\) with \(M\) symmetric.  The
first fundamental theorem for \(\SL_2\) says that they generate its
exterior-algebra invariants \cite{WeylClassicalGroups}.  Expanding the exponentials and polarizing
therefore shows that the classes in
\eqref{coh:graph-exponential} span those invariants.  A sufficiently
large finite graph list has the same span.  Since \(E\) is non-CM,
these are exactly its power's rational Hodge classes.  Their Euler
pairing is nondegenerate and agrees, up to one fixed nonzero scalar,
with the intersection pairing under \(\Phi\): on pairs of graph
exponentials this is the determinant formula, and those classes span.

Combine the torus transform with the sphere isometry to define
\begin{equation}\label{coh:dictionary}
 \Psi=(\rho\oplus\id_U)\otimes\Phi:
 H^*(Y,\Q)\otimes H^*(T_g,\Q)
 \xrightarrow{\sim}(P_B\oplus U)\otimes\bigwedge W.
\end{equation}
Here the target still contains the formal Lefschetz space \(U\),
rather than \(U_B\) with its actual Hodge structure.  When \(Y\) is
a point, set \(\Psi=\Phi\).

We can now state precisely what the Euler tests recover.  In the
quartic case let
\begin{equation}\label{coh:input-space}
 C_*\in (P_B\oplus U)\otimes\bigwedge W
\end{equation}
be rational and invariant under the simultaneous \(\SL_2\), acting
trivially on \(P_B\).  In the torus-only case take
\(C_*\in\bigwedge W\) invariant under \(\SL_2\).
Set \(a=\Psi^{-1}(C_*)\).  This lies in \(H^n(D,\Q)\): total Cartan weight
zero means middle degree, and \(\SL_2\)-invariance gives
\(\omega a=0\).
If the annihilator of \(a\) satisfies
\eqref{setup:wedge-injective}, Proposition~\ref{prop:realization}
gives a perfect complex.  Divide its Mukai vector by the common
nonzero rational scalar in the resulting Euler comparison, including
the fixed pairing normalization of \(\Psi\), and call the result \(C\).
Then
\begin{equation}\label{coh:determined-component}
 C|_{P_B}=C_*|_{P_B};
 \qquad\text{in the torus-only case, }C=C_*.
\end{equation}
Indeed the sphere tests lie in \(P_Y\), and the graph tests span all
the relevant torus Hodge classes.  The \(P_B\)-component of an
algebraic class is itself Hodge.  The nondegenerate pairings therefore
determine that component.  They do not yet determine the component in
\(T(B)\).

\subsection{The annihilator and cap operators}

The degree-two annihilator of \(a\) can be calculated directly on
\(C_*\).  A K3 degree-two class becomes
\(x\in P_B\oplus\langle t_0\rangle\).  Its cup action is
\begin{equation}\label{coh:lowering-k3}
 D_x(t_+)=x,\qquad
 D_x(y)=-q(x,y)t_-\quad
   (y\in P_B\oplus\langle t_0\rangle),\qquad
 D_x(t_-)=0.
\end{equation}
The torus degree-two actions are all operators
\begin{equation}\label{coh:lowering-torus}
 F(\lambda,L,\pi)=\lambda\wedge-+F_L+\iota_\pi,
 \quad
 \lambda\in\bigwedge^2\langle b_i\rangle,
 \quad \pi\in\bigwedge^2\langle a_i^*\rangle,
\end{equation}
where \(F_L\) is the derivation sending the \(a\)-copy of \(M_0\)
to the \(b\)-copy by an arbitrary matrix \(L\), and vanishing on the
\(b\)-copy.  On the tensor product the full operator is
\(D_x\otimes1+1\otimes F(\lambda,L,\pi)\).  These formulas follow
from cup product on a surface and the creation--annihilation rules.

The actual Hochschild cap action on \(C\) has the same operator
formulas, after choosing bases for its Hodge weights.  The two
calculations have different inputs:
\[
\begin{array}{c|c|c}
 &\text{formal calculation}&\text{actual calculation}\\ \hline
 \text{class}&C_*&C\\
 \text{K3 complement}&U&U_B=T(B)\\
 \text{weights}&n-\deg\text{ on }H^*(D)&p-q\text{ on }H^{p,q}\\
 \text{operators}&H^2(D)\text{ by cup product}
   &\HH^2(B\times E^g)\text{ by cap product}
\end{array}
\]
Omit the K3 row and the factor \(B\) in the torus-only case.
Equality of these operator formulas does not yet identify \(C\) with
\(C_*\), or the Hodge circle of \(U_B\) with the formal action on
\(U\).  It will let us constrain the unknown component of \(C\)
using the small \(\Ext^2\) space.  We now justify the actual formulas
directly on the mirror, without a Hochschild comparison under HMS.

Use Todd-modified Hochschild--Kostant--Rosenberg isomorphisms.  On
homology the modification wedges with \(\sqrt{\td}\); on cohomology
use the corresponding contraction modification.  The module theorem
of Calaque--Rossi--Van den Bergh identifies cap product with the
polyvector action on Hodge cohomology
\cite[Theorem~1.4]{CalaqueRossiVanDenBergh}.  It sends the Hochschild
class of a perfect complex to its Mukai vector.  We take normalized HKR,
with the factor \(1/p!\) in form degree \(p\), and totalized form
signs: holomorphic forms have Hochschild homological degree one;
antiholomorphic cochains precede polyvector contractions, with the
Koszul signs for moving each contraction through a cochain.

For this actual calculation, reuse the weight notation by choosing
\(t_+\in H^{2,0}(B)\) and
\(t_-\in H^{0,2}(B)\) with \(q(t_+,t_-)=1\).
Choose \(t_0\) in the remaining Hodge line of \(U_B\).
These are complex Hodge directions, in place of the formal basis
\(1,\omega_Y,-[\mathrm{pt}]\) used above.  The weight-zero space is
now the full \((p,p)\)-part of the Mukai space.  The three summands
\[
 H^2(B,\cO_B),\qquad H^1(B,T_B),\qquad
 H^0(B,\bigwedge^2 T_B)
\]
act as \eqref{coh:lowering-k3}, with \(x\) respectively in
\(H^4\), \(H^{1,1}\), and \(H^0\).  Here \(T_B\) in the displayed
groups is the tangent bundle.  Wedge product gives the first action;
contraction with a tangent-vector-valued one-form gives the middle
derivation, skew for the Mukai form; bivector contraction gives the
last action.  These observations verify the relative signs in
\eqref{coh:lowering-k3}, up to an immaterial rescaling of each acting
coordinate.  The same form calculus on the abelian factor gives
\eqref{coh:lowering-torus}, now with \(a,b\) the Hodge directions of
weights \(1,-1\).

There are no cross terms involving the K3 factor, because its
\(\HH^1\) vanishes.  Consequently
\begin{equation}\label{coh:hh-dimension}
 \dim\HH^2(B\times E^g)=22+\binom{2g}{2}=b_2(Y\times T_g),
\end{equation}
with the term \(22\) omitted in the torus-only case.  Passing from
Dolbeault-normalized to Betti-normalized classes does not alter these
rank calculations.  In an output component of type \((j,j+2)\), an
acting polyvector of degree \(b\in\{0,1,2\}\) uses the Mukai
component of half-degree \(j+b\).  The powers of the comparison unit
\(\pm2\pi i\) thus amount to invertible row and column rescalings.

\begin{lemma}\label{coh:rank-trace}
Let \(\mathcal E\) be supplied by
Proposition~\ref{prop:realization}, and let \(C\) be its normalized
Mukai vector.  The kernel of the operators
\eqref{coh:lowering-k3}--\eqref{coh:lowering-torus} acting on \(C\)
has dimension at least \(\dim\mathcal Z\).  If it has dimension
exactly \(\dim\mathcal Z\), the semiregularity trace on
\(\Ext^2(\mathcal E,\mathcal E)\) is injective.
\end{lemma}

\begin{proof}
The cap action on the class of an object factors as
\begin{equation}\label{coh:trace-factorization}
 \HH^2\longrightarrow\Ext^2(\mathcal E,\mathcal E)
       \longrightarrow\HH_{-2}.
\end{equation}
The first map is evaluation of a Hochschild operator on the object;
the second is the boundary--bulk trace.  Under HKR this trace is the
collection of traces against powers of the Atiyah class, followed by
an invertible Todd modification.  These are the diagonal-adjunction
and Atiyah--Chern character identities
\cite{CaldararuMukaiII,BuchweitzFlennerHH}.
They apply to all the components of the trace, not only to
\(H^2(\cO)\).  Rescaling the Mukai vector has no effect on rank.

By \eqref{coh:hh-dimension} and Proposition~\ref{prop:realization},
the composite in \eqref{coh:trace-factorization} has rank at most
\(b_2(D)-\dim\mathcal Z\).  This proves the kernel bound.  If
equality holds, the composite has rank equal to the upper bound on
the dimension of its middle space.  Both ranks must then equal that
dimension, so the trace out of the middle space is injective.
\end{proof}

\subsection{Deforming a complex along its Hodge locus}

We use the following precise form of semiregularity.  It isolates the
only deformation-theoretic input needed after the kernel calculations.
The obstruction principle is due to Buchweitz--Flenner; the version
for all infinitesimal obstructions and varying ambient varieties is
also a consequence of Pridham's theorem
\cite{BuchweitzFlennerSemiregularity,PridhamSemiregularity}.

\begin{proposition}\label{prop:semiregularity}
Let \(\mathscr X\to S\) be a smooth projective family over a smooth
complex algebraic base, let \(s\in S\), and let
\(\mathcal E\in\Perf(\mathscr X_s)\).  Suppose the parallel transport
of \(\ch(\mathcal E)\) remains Hodge on the germ of \(S\) at \(s\),
and that the map
\[
 \Ext^2(\mathcal E,\mathcal E)
 \longrightarrow\bigoplus_{j\ge0}H^{j+2}(\mathscr X_s,\Omega^j),
 \qquad o\longmapsto
  \bigl(\Tr(o\,\operatorname{At}(\mathcal E)^j)/j!\bigr)_j
\]
is injective.  Then \(\mathcal E\) extends to a compatible family of
perfect complexes over the formal germ of \(S\) at \(s\).

If \(S\) is connected and the Chern character is a global parallel
section remaining Hodge on \(S\), that section is algebraic on every
fiber.  The same conclusions hold for the Mukai vector in a family
whose Todd classes are parallel algebraic classes.
\end{proposition}

\begin{proof}
We first explain the obstruction calculation.  Resolve the initial
perfect complex by a bounded complex of vector bundles, and use the
Dolbeault resolution of its graded smooth bundle.  A deformation is
represented by a total-degree-one \((0,*)\)-superconnection squaring
to zero.  Relative smooth forms may be used over the Artin
neighborhoods of \(s\): in a smooth local trivialization take Taylor
coefficients in base variables and set antiholomorphic base variables
to zero.  Local holomorphic coordinates identify these complexes with
relative Dolbeault and de Rham resolutions.  Smooth projectivity gives
Hodge degeneration and cohomology base change.

Consider a small extension of Artin neighborhoods, with kernel \(I\)
annihilated by the maximal ideal.  Lift an already constructed
superconnection arbitrarily, retaining the relative Dolbeault Leibniz
rule.  Its square \(o\) is an \(I\)-valued cocycle in the central
endomorphism complex and represents the obstruction.  Add a smooth
\((1,0)\)-connection and denote the curvature of the resulting full
superconnection by \(R\).  The closed forms \(\str(R^m)\) represent
the horizontal Chern character, up to the usual nonzero normalization;
this follows from the superconnection Chern--Weil transgression
formula.

Sort the curvature by holomorphic form degree.  Its degree-zero part
is \(o\), and the other parts have degrees one and two.  Since
\(I^2=0\), every term with two copies of \(o\) vanishes.  Hence
\(\str(R^m)\) has no component below holomorphic degree \(m-1\),
and its component of degree \(m-1\) is
\begin{equation}\label{coh:obstruction-trace}
 m\,\str\bigl(o(R^1_0)^{m-1}\bigr),
\end{equation}
where \(R^1_0\) is the degree-one curvature at the central fiber.
It represents the Atiyah class.  The horizontal Chern character lies
in filtration \(F^m\) by hypothesis.  Hodge degeneration therefore
makes \eqref{coh:obstruction-trace} zero in Dolbeault cohomology.
This remains true with coefficients in \(I\) on the central fiber,
by freeness and cohomology base change.  All semiregularity traces of
the obstruction vanish.  Injectivity makes \(o\) exact, so a
correction of the lifted superconnection makes its square zero.
Iterating over the powers of the maximal ideal gives compatible lifts.

These lifts are perfect complexes, not merely formal differential
operators.  Locally on a coordinate disk, assume inductively that the
operator modulo \(I\) has been gauged to ordinary Dolbeault plus a
holomorphic differential.  Lift that gauge and an arbitrary
holomorphic matrix lift of the differential.  The difference has
coefficients in \(I\).  Its highest positive form-degree component is
Dolbeault closed by the square-zero equation.  A local primitive gives
a gauge transformation removing it.  Continue downward in form
degree.  The remaining square-zero operator is ordinary Dolbeault
plus a holomorphic differential.  Thus it is locally the Dolbeault
resolution of a perfect complex.  The relative forms are flat over
the Artin base, so the chosen reductions compute derived restriction.
GAGA identifies these analytic lifts with algebraic perfect complexes.

The fixed bounded graded bundle in this local straightening gives
a uniform finite Tor amplitude for the compatible lifts.
Grothendieck existence algebraizes the resulting pseudocoherent tower
over the completed local base
\cite[Theorem~38.29.6, Tag~0DIG]{PerfectGrothendieckExistence}.
The family is proper, flat, and finitely presented, and the restriction
to the closed fiber is perfect.  The relative-perfectness criterion
\cite[Lemma~37.71.3, Tag~0DJZ]{PerfectGrothendieckExistence} therefore
makes the algebraization perfect relative to the completed base.
That base is regular and the family is smooth, so the total space is
regular; the resulting bounded coherent complex is thus perfect.
Spread its finite presentation to a finitely generated subalgebra of
the completion over an affine neighborhood of \(s\), retaining the
closed point and its original complex.  The resulting finite-type
parameter is dominant over that neighborhood: the base is smooth at
\(s\), and its local ring maps injectively to its completion.
On the relevant irreducible component, topological constancy of Chern
characters identifies the classes of the spreading complex with the
pullback of the prescribed parallel section.  Its image contains a
dense open subset of \(S\).

For completeness, the last passage to every fiber uses proper cycle
parameters, not constancy of a chosen complex.  Express algebraicity
of the fixed rational parallel class by finite sums and differences
of cycles with a common positive denominator.  Relative projective
Hilbert schemes give countably many proper parameter spaces for such
expressions; restrict to connected components on which the cycle
class is the prescribed one.  Their images in \(S\) are closed.
A dense open subset of an irreducible complex variety cannot be
covered by countably many proper closed subvarieties.  One such
proper image is therefore all of \(S\).  This proves algebraicity on
every fiber.  A connected smooth base is irreducible componentwise
with no intersecting components, so the stated connected case follows.
Finally, multiplication by \(\sqrt{\td}\) is an invertible parallel
algebraic operation in the indicated families and changes none of
the conclusions.
\end{proof}

The remaining task is now finite-dimensional: choose \(C_*\) so that
the annihilator has the required wedge-injectivity, and show that the
unknown Mukai component is forced by the cap-kernel bound.  Once the
bound is sharp, Proposition~\ref{prop:semiregularity} transports the
resulting algebraic tensor through the desired period family.

\section{Exact Kuga--Satake correspondences}
\label{sec:exactks}

We now use the realization and deformation results to construct the
prescribed Kuga--Satake map.  The distinction between that map and an
unspecified nonzero embedding matters: the full Clifford action will
subsequently control tensor classes and the actions on different K3
factors.

For a projective K3 surface $S$, put
\[
 T(S)=\NS(S)_{\Q}^{\perp}\subset H^2(S,\Q),
 \qquad W_S=C^+(T(S),q),
\]
where $q$ is the cup form and our Clifford convention is
$v^2=q(v,v)$.  Give $W_S$ its standard left Kuga--Satake Hodge
structure.  Thus, for an oriented orthonormal basis $u_1,u_2$ of
the positive period plane, its complex structure is left
multiplication by $u_1u_2$.  This polarizable weight-one structure
is $H^1(A_S,\Q)$ for an abelian variety, up to the choice of lattice
and isogeny \cite{KugaSatake,DeligneK3,VanGeemenKS}.  If
$w\in T(S)$ has nonzero square, the standard tensor is
\begin{equation}
 \label{ks:standard-map}
 j_w:T(S)(1)\longrightarrow\End(W_S),
 \qquad j_w(v)(x)=vxw.
\end{equation}
A rational polarization identifies $W_S^\vee$ with $W_S(1)$ and
therefore converts \eqref{ks:standard-map} into a Hodge map
$\iota_w:T(S)\to W_S\otimes W_S\subset H^2(A_S^2,\Q)$.

\begin{theorem}
 \label{thm:exact-ks}
Let $S$ be a smooth projective complex K3 surface.  For every
nonisotropic $w\in T(S)$, every full even-Clifford realization
$H^1(A_S,\Q)=C^+(T(S),q)$, and every rational polarization used
to define $\iota_w$, there is a cycle
\[
 \Gamma\in\CH^2(S\times A_S^2)_{\Q}
 \quad\text{such that}\quad
 \Gamma_*|_{T(S)}=\iota_w.
\]
\end{theorem}

The proof first constructs this exact map at the mirror quartic
$B$.  Auxiliary terms in its Mukai vector make the cap-action
kernel rigid enough to determine the unknown transcendental
component.  Equality in the Ext bound then permits deformation
through a full polarized K3 family.  A rational comparison of
polarization lines and specialization reach every projective K3
surface.

\subsection{A spin-invariant tensor at the mirror quartic}
\label{ks:construction}

Retain the spaces $P_B$ and $U$ from the cohomological dictionary.
Choose a rational ample class $h\in\NS(B)_{\Q}$ and set
\[
 N=h^\perp\cap\NS(B)_{\Q},\qquad V=N\perp U.
\]
Then $\dim N=18$, while $V$ has dimension $21$ and signature
$(2,19)$.  The square class of $q(h,h)$ can be prescribed as
$2d$ for any positive integer $d$; we explain this choice when
comparing polarization degrees in Section~\ref{ks:deformation}.

For this $V$, write
\[
 W=C^+(V).
\]
The left spin action on $W$ admits a nondegenerate invariant
alternating form that polarizes the Kuga--Satake structure on
the chosen period component.  One may use
\begin{equation}
 \label{ks:trace-form}
 E(x,y)=\operatorname{Tr}(\alpha x^*y),
\end{equation}
where $*$ reverses Clifford products and $\alpha^*=-\alpha$.
For example, take $\alpha$, up to sign, to be the product of
two orthogonal rational vectors spanning a positive plane.
At that plane, let $I$ be their normalized product.  Evaluating
$\operatorname{Tr}(I x^* I x)$ on the even Clifford monomial
basis gives a definite form: commuting the two positive
generators past a monomial compensates the parity of its
negative generators.  Choose the sign giving positivity.
Spin invariance then gives positivity on the entire component.

Let $\theta\in\bigwedge^2W$ be the nondegenerate alternating
tensor corresponding to this form, and choose a rational volume
element $z$ in the full Clifford algebra of $V$.  Define
\begin{equation}
 \label{ks:J}
 J:V\longrightarrow\bigwedge^2W,
 \qquad J(v)=(L_{vz}\otimes1)\theta.
\end{equation}
Here $L_a$ denotes left multiplication by $a$.  Since $\dim V=21$,
the volume element is central and $vz$ is reversal-even.
Consequently $L_{vz}$ is self-adjoint for \eqref{ks:trace-form},
so \eqref{ks:J} is alternating.  It is also Spin-equivariant.

Restriction to the split group $\Spin(U)=\mathrm{SL}_2$ gives
\[
 W=M_0\otimes H,
 \qquad H=\Q e_+\oplus\Q e_-.
\]
Thus this is precisely a torus space in the preceding dictionary,
with $g=\dim M_0=2^{19}$.  To describe the operators on this large
multiplicity space, work over $\C$.  The even Clifford algebra in
dimension $21$ is the endomorphism algebra of its spin module
$\Sigma$, of dimension $2^{10}$.  Its regular representation is
therefore a sum of $2^{10}$ copies of $\Sigma$.  On restriction to
$N\perp U$, the module $\Sigma$ is $M_1\otimes H$, where $M_1$
is the $2^9$-dimensional irreducible module for the full Clifford
algebra of $N_{\C}$.  Consequently
\begin{equation}
 \label{ks:multiplicity}
 W_{\C}=\Sigma\otimes R_0=(M_1\otimes H)\otimes R_0,
 \qquad M_{0,\C}=M_1\otimes R_0,
\end{equation}
where $R_0$ is the regular-representation multiplicity space.
The module $\Sigma$ has an invariant alternating form.  The tensor $\theta$ is the
product of a symmetric nondegenerate tensor $Q$ on $M_0$ with
the alternating tensor on $H$.  Moreover, $Q$ factors as the
form on $M_1$ times a symmetric form on $R_0$.

For $y\in N$, the operator $L_{yz}$ is, up to a common nonzero
scale, the Clifford generator $\Gamma_y$ on $M_1$ tensored with
the identity on $R_0\otimes H$.  These generators are
self-adjoint for $Q$.  Let $\Delta$ be their chirality operator;
it anticommutes with every $\Gamma_y$ and is skew-adjoint,
since reversal in dimension $18$ has negative sign.  For
$u\in U$, $L_{uz}$ is, up to a common scale, $\Delta$ times
the vector action on $H$.  It follows that $J(y)$ has symmetric
matrix proportional to $\Gamma_yQ$ and alternating $H$ factor,
whereas $J(u)$ has matrix proportional to $\Delta Q$ and
belongs to the $\operatorname{Sym}^2H$ part.

The Clifford generators on $M_1$ generate its full endomorphism
algebra.  An operator commuting with all of them can still act
arbitrarily on $R_0$.  We remove this multiplicity freedom by choosing
two further rational Spin-invariant nondegenerate
alternating tensors $\theta_1,\theta_2$ whose relative
operators to $\theta$ have scalar common commutant on $R_0$.
Such choices exist: over $\C$ two general symmetric operators
relative to the multiplicity form have this property, which is
a nonempty Zariski-open condition in the rational space of
invariant tensors.

We next place these tensors in separate exterior degrees.  This will
let an annihilation equation distinguish the wedge, derivation, and
contraction parts of a torus operator.  Choose a rational basis
$e_1,e_2,e_3$ of $N^\perp\cap P_B$,
each vector lying in one of $\Q h$, $H^0(B,\Q)$, and
$H^4(B,\Q)$.  For dual bases $(v_j)$ and $(v_j^\vee)$ of
$V$, put
\begin{equation}
 \label{ks:test-tensor}
 \begin{split}
 R_1&=1,\qquad
 R_2=\theta^4+\theta_1^7+\theta_2^{10},\qquad
 R_3=\theta^{13},\\
 C_*&=\sum_jv_j\otimes J(v_j^\vee)
            +\sum_{s=1}^3e_s\otimes R_s.
 \end{split}
\end{equation}
The first summand contains the desired Clifford map $J$.
The three $e_s$ components impose the additional equations just
described; after those equations remove the unwanted operators,
the remaining kernel will consist of the $19$ simultaneous spin
lowering directions.  We prove this next.
The full tensor is invariant under the simultaneous $\mathrm{SL}_2$
action.  Its inverse image under the dictionary is therefore
a rational middle-degree class $a$ killed by the symplectic
form.

\subsection{The formal kernel has dimension nineteen}
\label{ks:formal-kernel-section}

For a tensor in $U\otimes\bigwedge W$, write its three weight
components as
\begin{equation}
 \label{ks:weight-components}
 t_+\otimes U_-+t_0\otimes U_0+t_-\otimes U_+.
\end{equation}
Here $U_-,U_0,U_+\in\bigwedge W$.  A superscript $*$ denotes
the components of $C_*$; we will later use the same notation for
the component in $U_B\otimes\bigwedge W$.
In particular, $U_-^*$ is nonzero and has exterior degree two.
The annihilator of $a$ by degree-two classes is computed by
the operators $D_x+F(\lambda,L,\pi)$ in
\eqref{coh:lowering-k3} and \eqref{coh:lowering-torus}.

\begin{lemma}
 \label{ks:formal-kernel}
For the class $a$ defined by \eqref{ks:test-tensor}, its
annihilator $\mathcal Z\subset H^2(Y\times T_g,\Q)$ has
dimension $19$ and projects injectively to $H^2(Y,\Q)$.
In particular, it satisfies \eqref{setup:wedge-injective}.
\end{lemma}

\begin{proof}
Decompose
$x=x_N+x_e+x_0t_0$, where $x_e=\sum_sx_{e,s}e_s$.
The $e_s$ components of the annihilation equation are
\begin{equation}
 \label{ks:e-equations}
 F(R_s)+x_{e,s}U_-^*=0.
\end{equation}
For $s=2$, the second term has exterior degree two.  On a
power $\theta_i^k$, the contraction, derivation, and wedge
terms of $F$ have degrees $2k-2,2k,2k+2$, respectively.
For $k=4,7,10$ these nine degrees are distinct and exceed
two.  Hence $x_{e,2}=0$, and each of the nine terms vanishes
separately.

Symplectic Lefschetz injectivity now gives $\lambda=0$ from
$\lambda\wedge\theta^4=0$.  The derivation terms give
$F_L(\theta_i)=0$ for $i=0,1,2$, where $\theta_0=\theta$.
For the contraction term, contracting two $a$ directions into
$\theta^4$ gives, up to a nonzero factor, the two-$b$ tensor
obtained from $\pi$ by $Q$, wedged with $\theta^2$.
The same injectivity gives $\pi=0$.  All these wedge maps
are in the injective range for $g=2^{19}$.  The remaining
equations \eqref{ks:e-equations} give $x_e=0$.

The equation $F_L(\theta)=0$ says that $L$ is self-adjoint
for $Q$.  Preservation of $\theta_1$ and $\theta_2$ says
in addition that $L$ commutes with their relative operators.
If $x=0$, the $N$ equations become
$F_L(J(y))=0$ for every $y\in N$.  By the matrix description
above, these are exactly
$[L,\Gamma_y]=0$.  Irreducibility on $M_1$, followed by
the common-commutant condition on $R_0$, makes $L$ scalar.
The equation $F_L(U_0^*)=0$ kills this scalar, because
the lowering action on the highest-weight-two representation
is nonzero.  The kernel therefore projects injectively to
$N\oplus\Q t_0$, a space of dimension $19$.

Conversely, each $x\in N\oplus\Q t_0$ occurs in the kernel.
The operator $D_x$ is the lowering-degree element of
$\mathfrak{so}(V)$ displayed in \eqref{coh:lowering-k3}.
Its simultaneous spin action preserves \eqref{ks:test-tensor}.
Since $W$ has only Cartan weights $1$ and $-1$, its lowering
action is of the form $F_L$.  This proves the dimension
assertion.  Finally, projection of $H^3(Y\times T_g,\Q)$
to $H^2(Y,\Q)\otimes H^1(T_g,\Q)$ identifies the map in
\eqref{setup:wedge-injective} with the tensor product of
the injective projection of $\mathcal Z$ and the identity.
\end{proof}

\subsection{Determining the transcendental component}
\label{ks:exactness-section}

Proposition~\ref{prop:realization} now supplies a perfect
complex on $B\times E^g$.  Divide its Mukai vector by the
common nonzero rational scalar in the Euler comparison and
call the result $C$.  The dictionary determines
$C|_{P_B}=C_*|_{P_B}$, but does not initially determine
its component on $U_B=T(B)$.  The Ext bound does.  In the
following statement, give the formal space $U$ the Hodge
structure induced by the actual elliptic structure on $H$
through its highest-weight-two representation.

\begin{lemma}
 \label{ks:exactness}
There is a rational Hodge isometry $U_B\simeq U$, for the
actual Hodge structure on $H=H^1(E,\Q)$, under which
$C=C_*$.  The semiregularity trace of the perfect complex
is injective.
\end{lemma}

\begin{proof}
\emph{Excluding a zero transcendental component.}
Choose normalized actual Hodge directions $t_+,t_0,t_-$
on $U_B$, as in the cap calculation of
Section~\ref{sec:cohomology}, and write its unknown component
in the form \eqref{ks:weight-components}.
By Lemma~\ref{ks:formal-kernel} and the Ext bound, the actual
cap kernel has dimension at least $19$.  Imposing the three
linear conditions $x_e=0$ leaves a subspace of dimension at
least $16$.  On this subspace the
equations $F(R_s)=0$ again imply $\lambda=\pi=0$, and
make $L$ self-adjoint and commuting with the two relative
operators.  The $N$ equations are
\begin{equation}
 \label{ks:N-equation}
 F_L(J(y))+q(x_N,y)U_-=0\qquad(y\in N).
\end{equation}
The $t_-$ component, using \eqref{coh:lowering-k3}, is
\begin{equation}
 \label{ks:t-minus-equation}
 F_L(U_+)=J(x_N)+q(x_0t_0,t_0)U_0.
\end{equation}
If the entire component \eqref{ks:weight-components} were
zero, Equation~\eqref{ks:N-equation} would make $L$ scalar,
and \eqref{ks:t-minus-equation} would give $J(x_N)=0$,
hence $x_N=0$.
Only the scalar $L$ and $x_0$ would remain, contrary to
the dimension lower bound $16$.  Thus the component is
nonzero.

\emph{Identifying the Hodge representation.}
This nonzero component realizes the simple, three-dimensional
K3 Hodge structure $U_B$ in the exterior algebra of copies of
$H$, with the indicated Tate twists.  Since $E$ is non-CM,
its irreducible tensor representations are the irreducible
$\mathrm{SL}_2$ representations.  The only possible one here
has highest weight two.  The resulting rational Hodge
identification $U_B\simeq U$ is a similarity: the invariant
forms on this irreducible representation are proportional.
The underlying quadratic spaces are already rationally
isometric by the dictionary.  In odd dimension their similarity
multiplier is a rational square, as follows by taking
determinants; the signatures give the positive sign.  Rescaling
gives a Hodge isometry.  We may now use it to compare the two
sets of weight directions.  The component
\eqref{ks:weight-components} is equivariant for the simultaneous
$\mathrm{SL}_2$; in particular $U_-\ne0$.

\emph{Determining the component up to one scalar.}
Still on the subspace $x_e=0$, the kernel of projection
to $x_N$ has dimension at most two: setting $x_N=0$
in \eqref{ks:N-equation} again makes $L$ scalar, leaving
at most that scalar and $x_0$.  The image in $N_{\C}$
therefore has dimension at least $14$.  A totally isotropic
subspace of the $18$-dimensional nondegenerate space
$N_{\C}$ has dimension at most nine.  Choose a kernel
element for which $q(x_N,x_N)\ne0$.

For every $y\perp x_N$, Equation~\eqref{ks:N-equation}
says that $L$ commutes with $\Gamma_y$.  On the irreducible
Clifford module $M_1$, the commutant of these $17$
directions is spanned by $1$ and $\Delta\Gamma_{x_N}$;
this is the restriction of the even-dimensional Clifford
module to an odd-dimensional perpendicular.  The relative
operators again remove the multiplicity commutant, so
\[
 L=l\,1+m\,\Delta\Gamma_{x_N}.
\]
The Clifford relation gives
\[
 [\Delta\Gamma_{x_N},\Gamma_y]
   =\kappa\,q(x_N,y)\Delta
\]
for a fixed nonzero normalization constant $\kappa$.
Equation~\eqref{ks:N-equation} therefore makes $U_-$ a
nonzero multiple $cU_-^*$: the two-$b$ tensor with matrix
$\Delta Q$ is exactly the lowest-weight line of $J(U)$.
Equivariance gives this same scalar on all three
components in \eqref{ks:weight-components}.

\emph{Determining the exact normalization.}
Let $L^*$ be the spin lowering lift for
$x=x_N$ in the tensor $C_*$.  Comparing its version of
\eqref{ks:N-equation} with the actual equation gives
\[
 L=cL^*+\widetilde l\,1.
\]
Now use Equation~\eqref{ks:t-minus-equation} modulo the line
of $U_0$.  The scalar part $\widetilde l\,1$ of $L$
lowers $U_+$ into that line.  The equation for $C_*$,
with $x=x_N$, gives $F_{L^*}(U_+^*)=J(x_N)$.
Since $U_+=cU_+^*$, we obtain
\[
 J(x_N)\equiv F_L(U_+)
   \equiv c^2 F_{L^*}(U_+^*)=c^2J(x_N)
       \pmod{\C U_0}.
\]
One factor $c$ comes from $L$ and the other from $U_+$.
The tensor
$J(x_N)$ lies outside that line, since $J(N)$ and $J(U)$
have the distinct Clifford descriptions above.  Hence
$c^2=1$.  Reversing the sign of the transcendental Hodge
isometry if necessary gives $C=C_*$.

The circle on the resulting $V=h^\perp$ has its unique
lift to the spin group, and its action on $W$ is the
actual circle on $H^1(E^g,\Q)$.  Thus this actual
weight-one structure is the stipulated Kuga--Satake
structure.  The formal and actual cap calculations now
agree, so the cap kernel has dimension exactly $19$.
Lemma~\ref{coh:rank-trace} makes the
semiregularity trace injective.
\end{proof}

\subsection{Deformation and all polarization degrees}
\label{ks:deformation}

We have constructed an exact tensor whose complex can deform
whenever its Mukai vector stays Hodge.  We first use this to
move away from $B$ while preserving its polarization degree.
We then compare that degree with any prescribed degree by a
rational isometry, and specialize the exact tensor to fibers
with larger Picard rank.

\paragraph{Deformation from the mirror quartic.}
Consider the full family of $h$-polarized K3 surfaces, coupled
with their Kuga--Satake abelian varieties.  These are algebraic families
after a suitable level cover.  More precisely, the primitive
period domain and its spin cover, together with the invariant
rational alternating form, give a holomorphic equivariant
map to Siegel space.  Arithmetic subgroups under the
isogeny from Spin to $\SO$ have arithmetic image.  Intersect
with the stabilizers of the chosen lattices and pass to a
neat level; the K3 period theorem and algebraicity of the
arithmetic quotients and their morphisms then give a smooth
algebraic base with both universal families
\cite{DeligneK3,BailyBorel,BorelAlgebraicity}.

The lattice at the central point may be chosen to give
exactly the abelian variety $E^g$: preserve that lattice
by further level restriction and multiply the rational
polarization by an integer, with the sign chosen above.
On this spin level cover the tensors in
\eqref{ks:test-tensor} are global parallel sections:
the vectors $h$, $H^0$, and $H^4$ remain fixed and algebraic, the
sum on $V$ is Spin-invariant, and every $\theta_i$ is
Spin-invariant.  Hence the full Mukai vector $C_*$ stays
Hodge throughout the family; so do its Todd factors.
Lemma~\ref{ks:exactness} and
Proposition~\ref{prop:semiregularity} give algebraicity
of this class on every fiber of the polarized component.

On a fiber $S\times A$, take the component in $H^2(S)\otimes H^2(A)$
and project the K3 factor to
$V=h^\perp$.  The additional terms $e_s\otimes R_s$ have
abelian degrees $0$, $8$, $14$, $20$, or $26$, so these
algebraic projections extract $J$.  To regard $J$ as a correspondence
with two $H^1$ slots, use the surface and abelian
K\"unneth projectors, the polarization and its inverse,
and the inclusion of an exterior square into two ordered
slots.  All are algebraic: the inverse polarization is
provided by the Lefschetz standard conjecture for abelian
varieties, and the slot conversion can also be realized
using the addition map and pullback.

This already determines the usual standard tensor, not
only a tensor with the volume element on the right.
Since $z$ is central, $L_{vz}=L_vR_z$ on the even space.
For any nonisotropic $w\in V$, replace $R_z$ by $R_w$
using right multiplication by the invertible even
element $z^{-1}w$.  Every right even multiplication is
a weight-one Hodge endomorphism of the abelian variety,
and is therefore algebraic.  Thus the map
$x\mapsto vxw$ is algebraic in this full primitive
Clifford construction.

\paragraph{Comparison of polarization degrees.}
Fix a prescribed degree $2d$.  The initial rational ample
class $h$ on $B$ can be chosen with square $2d$.  Indeed,
Meyer's theorem gives an isotropic rational vector in the
indefinite rank-$19$ space $\NS(B)_{\Q}$, hence a rational
hyperbolic plane.  The quadric $q(x,x)=2d$ consequently has
rational points, and rational approximation on this quadric
gives a point in the open ample cone.  Scaling to a primitive
integral class changes its square by a rational square.
Thus the primitive integral class chosen on $B$ has square $2d$
times a rational square.  Witt's extension theorem
therefore gives a rational isometry of the K3 quadratic
spaces carrying its polarization line to the line of
a primitive vector of square $2d$.  The isometry
identifies the primitive period domains.  At very general
periods on both sides the surfaces have Picard rank one:
the excluded wall hyperplanes, including their translates
under the rational isometry, are rational hyperplanes
and can all be avoided.  The K3 period theorem supplies
both polarized surfaces.  Buskin's theorem makes their
rational Hodge isometry algebraic \cite{Buskin}.
If $f$ denotes its restriction to the primitive spaces,
the induced Clifford Hodge isomorphism satisfies
\[
 C^+(f)\,j_w(v)\,C^+(f)^{-1}
       =j_{f(w)}(f(v)).
\]
This weight-one isomorphism is induced by an abelian
isogeny up to a rational multiple.  Composing it with
the K3 correspondence, and converting polarizations
as above, therefore transports the exact tensor.
We obtain the exact standard tensor very generally in
degree $2d$.  One may make this comparison on marked
components or on commensurable level quotients of the
two domains.

\paragraph{Specialization and transcendental restriction.}
Specialization gives the full primitive tensor at
every polarized fiber.  Here the class being specialized
is the prescribed parallel tensor in the full
Kuga--Satake family, so Picard jumps cause no loss of
information.  Equivalently, apply the specialization
theorem and its primitive-map corollary in
\cite[Theorem~1.2 and Corollary~5.5]{CKS}: the very-general nonzero
algebraic input required there has just been supplied
in its stronger exact form.  This specialization is
the proper-cycle-parameter argument of
Proposition~\ref{prop:semiregularity}; it does not
replace actual algebraicity by the Hodge condition.

\begin{proof}[Proof of Theorem~\ref{thm:exact-ks}]
Every projective K3 surface has a primitive ample class,
so the preceding construction gives its full primitive
Kuga--Satake tensor.  The projection to $T(S)$ is algebraic,
using the divisor intersection pairing and the degree-zero
and top-degree projectors.  For the prescribed
$w\in T(S)$, restrict the even Clifford space to
\[
 C^+(T(S))\subset C^+(h^\perp).
\]
This is a rational weight-one sub-Hodge structure, and
$x\mapsto vxw$ preserves it for $v\in T(S)$.  It can
therefore be split off by algebraic abelian
correspondences.  Changing its lattice, isogenous
realization, or rational polarization also uses algebraic
abelian correspondences.  The resulting cycle has the
exact action \eqref{ks:standard-map} with the chosen
polarization conversion, proving the theorem.
\end{proof}

\begin{corollary}
 \label{cor:self-powers}
For every projective complex K3 surface $S$, the rational
Hodge conjecture holds in every degree on every self-power
$S^m$.
\end{corollary}

\begin{proof}
Theorem~\ref{thm:exact-ks} supplies the exact standard
tensor on the full even Clifford space, for one
nonisotropic $w$ and one polarization.  These are the
hypotheses of the general self-power criterion
\cite[Theorem~1.1]{AKS}.  That criterion is independent
of the special quadratic locus used elsewhere in the
same reference to construct correspondences.
\end{proof}

\section{Degree-four tensors in spin families}
\label{sec:auxiliary}

The exact Kuga--Satake correspondences control each K3
surface separately.  To compare several factors, we will
also need a quadratic contraction tensor on certain
abelian varieties.  The construction below permits
independent period planes at every real place of a
totally real field.  Its four weight-one slots will
later compare actions on different spin constituents.

Let $F'/\Q$ be a totally real Galois extension and let
$(V',q')$ be a quadratic space over $F'$ of dimension $d'$.
Assume that
\begin{equation}
 \label{aux:quadratic-assumptions}
 d'\ge5,\quad d'\equiv5\pmod8,\qquad
 \operatorname{sig}(V'_\sigma)=(2,d'-2)
       \quad(\sigma:F'\hookrightarrow\R),\qquad
 \operatorname{Witt}(V')\ge2.
\end{equation}
The dimension can be increased as needed in the
application.  Choose an
oriented positive period plane at each real place.  The
corresponding spin circles give the weight-one Hodge
structure
\begin{equation}
 \label{aux:weight-one}
 W'=\operatorname{Res}_{F'/\Q}C^+_{F'}(V').
\end{equation}
It is polarizable: choose an invariant alternating form
with the required positivity at every real embedding,
using approximation over $F'$ in the construction
\eqref{ks:trace-form}, and take its trace to $\Q$.
Thus $W'$ is the rational first cohomology of an abelian
variety $A'$, with a chosen lattice.  On each component
the period construction gives an algebraic family of
abelian varieties.  Indeed, a neat arithmetic subgroup
of the restriction-of-scalars spin group can be chosen
to preserve the lattice and polarization.  The product
of the orthogonal period domains then maps
holomorphically and equivariantly to Siegel space, and
the algebraicity theorems used in
Section~\ref{ks:deformation} apply.

Choose a volume element $z'$ over $F'$.  Define a
rational tensor
\begin{equation}
 \label{aux:P}
 \mathcal P\in(\End W')^{\otimes2}
\end{equation}
by prescribing its extension to $\C$.  For each embedding
$\sigma:F'\hookrightarrow\C$, take dual bases
$(v_j)$ and $(v_j^\vee)$ of $V'_\sigma$ and form
\[
 \sum_jL_{v_jz'}\otimes L_{v_j^\vee z'}.
\]
Each operator is supported on the $\sigma$ summand of
$W'_\C$; sum these tensors over $\sigma$.  This is the
trace construction, so it descends to $\Q$.  It is
invariant under $\operatorname{Res}_{F'/\Q}\Spin(V')$.
Using an invariant polarization in each endomorphism
factor, regard $\mathcal P$ as a degree-four Hodge
tensor in four ordered $H^1(A')$ slots.

\begin{proposition}
 \label{prop:auxiliary-spin}
For the data \eqref{aux:quadratic-assumptions}, the tensor
$\mathcal P$ is the class of a cycle in
$\CH^2((A')^4)_{\Q}$, for every choice of the period planes and
every isogenous realization of \eqref{aux:weight-one}.
\end{proposition}

\begin{proof}
We first choose a point of the family whose abelian
variety is a power of the non-CM elliptic curve $E$
in the torus dictionary.  The Witt-index assumption
gives a split rank-three quadratic subspace of $V'$
with signature $(2,1)$ at every real place: take one
hyperbolic plane and a vector positive at all places
in a second one.  Its spin group is $\mathrm{SL}_2$
over $F'$.  Restricting scalars and then taking the
diagonal $\mathrm{SL}_2$ over $\Q$ makes $W'$ a
multiple of its standard representation.

The common circle can be placed in any prescribed
component of the period family.  Under the ternary
identification, conjugate by an element of
$\mathrm{GL}_2(F')$ with the chosen determinant signs
at the real embeddings.  Approximation permits these
signs to be chosen independently.  This also accommodates
either elliptic marking convention in the graph
dictionary; its common line sign was fixed before the
choice of tensor.  Choose the lattice to give the
corresponding elliptic power exactly at this point.

Apply the torus-only realization with
\[
 W=(W')^{\oplus4}=M_0\otimes H,
\]
using bases compatible with the four copies.  Place
$\mathcal P$ in $\bigwedge^4W$, with one vector in
each copy, after the polarization conversion.  These copies
have a second use: to verify the wedge-injectivity hypothesis,
we will project a degree-one class from one copy against an
annihilator's degree-two component on another.  We arrange
that the annihilator has the same component on every copy.

Let
$\theta$ be the direct sum of four identical copies
of a nondegenerate rational Spin-invariant alternating
tensor on $W'$.  Choose two more tensors
$\theta_1,\theta_2$ by tensoring that tensor with
nondegenerate symmetric matrices on the four copy
indices: the first is diagonal with distinct entries,
and the second has every entry nonzero.  Here the relative
operators act on the four copy indices.  Their commutation
conditions will force identical blocks there, while allowing
operators within $W'$; this is the copy condition just required.
As in the earlier construction, separate the powers so that
wedge, derivation, and contraction terms have different degrees.
Put
\begin{equation}
 \label{aux:test-tensor}
 C_* =\mathcal P+\theta^6+\theta_1^9+\theta_2^{12}.
\end{equation}
This is Spin-invariant and hence invariant under the
chosen simultaneous $\mathrm{SL}_2$.  The torus
dictionary produces a middle-degree primitive class
$a$.  Here
\[
 g=\tfrac12\dim_{\Q}W=[F':\Q]2^{d'}\ge32,
\]
so $g$ is even, the realization dimension condition is
satisfied, and the symplectic wedge maps used below
are in their injective range.

We verify the wedge-injectivity hypothesis of
Proposition~\ref{prop:realization}.  In the formal
annihilation equation $F(\lambda,L,\pi)C_*=0$, the
three operator types on the displayed powers have
degrees $(10,12,14)$, $(16,18,20)$, and $(22,24,26)$.
They are all distinct and exceed the degrees $2,4,6$
obtained by acting on $\mathcal P$.  The same symplectic
wedge and contraction calculation as in
Lemma~\ref{ks:formal-kernel} therefore gives
$\lambda=\pi=0$ and
\[
 F_L(\theta)=F_L(\theta_1)=F_L(\theta_2)=0.
\]
The first condition makes $L$ self-adjoint.  Commutation
with the diagonal relative operator makes it block
diagonal on the four copies.  Commutation with the
second relative operator, whose off-diagonal entries
are nonzero, then makes all four blocks equal.

Under the inverse dictionary, the annihilator
$\mathcal Z$ consequently consists of degree-two forms
obtained by repeating the same form on each of the
four torus copies.  Its projection to the degree-two
cohomology of any one copy is injective.  For an input
in $H^1$ of copy $i$, choose $j\ne i$ and project its
wedge product to multidegree one on copy $i$ and
degree two on copy $j$.  This is the tensor product
of the identity on $H^1$ of copy $i$ with that
injective projection of $\mathcal Z$.  Inputs from
the other copies have different multidegrees.  Taking
these projections for all $i$ proves
\eqref{setup:wedge-injective}.

In the torus-only case, the Euler tests determine
the entire normalized Mukai vector, so the resulting
perfect complex has $C=C_*$.  At this elliptic point,
the actual cap kernel and the formal annihilator
have the same dimension: the two sets of weights
on this invariant tensor are conjugate over $\C$.
The Ext bound is therefore an equality at the cap
rank, and the semiregularity trace is injective.

Throughout the indicated algebraic period family,
all terms of \eqref{aux:test-tensor} remain parallel
Hodge classes by Spin invariance.  Apply
Proposition~\ref{prop:semiregularity} and extract
cohomological degree four.  The three higher powers
have degrees $12$, $18$, and $24$, so the extracted
class is exactly $\mathcal P$.  Its one-slot-per-copy
component gives an algebraic class on $(A')^4$.
The elliptic starting point was available in every
component.  Finally, transport under abelian
isogenies and change of rational polarization use
algebraic correspondences, proving the assertion
for all the stipulated realizations.
\end{proof}

\section{Hodge classes on mixed products}\label{sec:mixed}

Fix projective K3 surfaces $S_1,\ldots,S_m$, and put
\[
 V_i=T(S_i)(1),\qquad W_i=C^+(T(S_i)).
\]
Here $V_i$ has weight zero, whereas $W_i$ carries the weight-one
Kuga--Satake structure.  Theorem~\ref{thm:exact-ks} supplies the exact
Clifford correspondences, and Corollary~\ref{cor:self-powers} supplies all
Hodge classes on each individual self-power.  We must still treat tensors
involving several different $V_i$.

The individual assertions determine the projections of the joint Hodge
group, but not the relations between those projections.  We first use
weight-one Hodge morphisms between the Kuga--Satake varieties to compare
the actions of shared simple factors projectively.  The CM theorem then
controls the torus action.  After removing these actions, only certain
central signs remain, arising from the two half-spin representations in
dimension four.  We encode their compatibility by a finite graph and use
the degree-four tensors of Proposition~\ref{prop:auxiliary-spin} to prove
the required relations.  All auxiliary abelian varieties will be chosen
before choosing the tensor automorphism to which this comparison is
applied.

\subsection{Tensor groups and the individual factors}

For a finite collection of the surfaces and abelian varieties used here,
let $G$ be its Hodge group, with Tate objects fixed.  Let $C$ be the group
of tensor automorphisms of its cohomology realizations that preserve
algebraic correspondences, again fixing Tate objects.  Then $G\subset C$.
We use these groups over $\Q$ and extend their coefficients to $\C$ when
decomposing representations.

We recall why $C$ has the required invariant interpretation.  The
Lefschetz standard conjecture holds for surfaces and abelian varieties
and is stable under products; it gives algebraic K\"unneth and Lefschetz
operators in the category under consideration
\cite{KleimanStandard}.  Hodge--Riemann then identifies numerical and
homological equivalence there.  Numerical semisimplicity
\cite{JannsenMotives}, with the symmetry modified to give the ordinary
flip on cohomology, gives a neutral Tannakian category of pure motives
with its cohomology fiber functor \cite{DeligneMilneTannakian}.
Its tensor automorphism group is the group just described.

In weight zero, the $C$-invariants are precisely algebraic tensors.
Indeed, before fixing Tate objects, invariants are morphisms from the
unit motive.  The central weight torus acts by cohomological weight and
maps onto the group acting on Tate objects up to isogeny.  It acts
trivially in weight zero, so fixing Tate objects does not change those
invariants.  Consequently it will suffice to prove
\begin{equation}\label{eq:mixed-group-goal}
 C\big|_{\bigoplus_i V_i}=G\big|_{\bigoplus_i V_i}
\end{equation}
for one finite collection containing the original surfaces and the
auxiliary abelian varieties constructed below.

\begin{lemma}[Individual actions]\label{lem:mixed-individual}
For any such collection and any $c\in C(\C)$, the action of $c$ on $V_i$
belongs to $\Hg(V_i)(\C)$.  Its simultaneous action on $V_i,W_i$ is
induced by a lift to the spin group of the quadratic space $V_i$.
\end{lemma}

\begin{proof}
Corollary~\ref{cor:self-powers}, equivalently the exact-tensor criterion
of \cite[Theorem~1.1]{AKS} applied after
Theorem~\ref{thm:exact-ks}, makes every Hodge tensor on powers of $V_i$
algebraic.  Algebraic degree projectors and the quadratic pairing allow
the usual tensor criterion for the Hodge group.  Hence $c|_{V_i}$ lies
in $\Hg(V_i)$.

Choose a spin lift of this orthogonal action.  The operator it induces
on $W_i$ preserves the exact KS tensor, all right even Clifford
multiplications, and an invariant polarization.  These tensors are
also preserved by $c$: the KS tensor is algebraic, while the right
multiplications are weight-one Hodge endomorphisms of an abelian
variety and therefore algebraic.  After removing the chosen spin lift,
the remaining operator on $W_i$ commutes with all right multiplications,
so it is left multiplication by an even Clifford element $a$.
Preservation of the KS map
\[
 v\longmapsto L_vR_w,\qquad w^2\ne0,
\]
then says that $a$ commutes with every $v\in V_i$.  The even part of
the center of the full Clifford algebra consists of scalars.  Thus the
remaining operator is scalar, and preservation of the polarization
makes it $1$ or $-1$.  Both scalars are supplied by the spin cover.
\end{proof}

Let
\[
 G_0=\Hg\left(\bigoplus_i V_i\oplus\bigoplus_i W_i\right).
\]
The projection from $G_0$ to the Hodge group of the $V_i$ has finite
kernel: the Hodge circle acts on each $W_i$ through the spin lift of
its action on $V_i$.  Thus the $W_i$ do not add a Lie algebra factor.

For each $i$, write $E_i=\End_{\rm Hdg}(V_i)$.  Zarhin's theorem says
that $E_i$ is either totally real or CM \cite{Zarhin}.  We use the
following notation:
\begin{itemize}
\item in the orthogonal case, $F_i=E_i$ and
  $d_i=\dim_{E_i}V_i\ge3$;
\item in the unitary case, $F_i=E_i^+$ is the maximal totally real
  subfield and $d_i=\dim_{E_i}V_i\ge1$.
\end{itemize}
Over $\C$, the corresponding groups and representations are
\begin{equation}\label{eq:mixed-single-groups}
 \begin{array}{c|c|c}
 &\Hg(V_i)_{\C}&(V_i)_{\C}\\ \hline
 \text{orthogonal}&\displaystyle\prod_{p:F_i\hookrightarrow\C}
       \mathrm{SO}_{d_i}
   &\text{one standard space at each }p\\[3pt]
 \text{unitary}&\displaystyle\prod_{p:F_i\hookrightarrow\C}
       \mathrm{GL}_{d_i}
   &\text{a standard space and its dual at each }p .
 \end{array}
\end{equation}
The quadratic form in the orthogonal case, and the underlying quadratic
form of the Hermitian form in the unitary case, has positive index two
at the distinguished real embedding and is negative definite at the
other real embeddings.  We call the distinguished embedding the
\emph{active place}.  The infinitesimal Hodge circle has weights
$2,-2$ once each on $V_i$, and weight zero elsewhere.

\subsection{Shared simple ideals and spin constituents}

The derived Lie algebra of $G_0$ is a subdirect product of the simple
ideals occurring in \eqref{eq:mixed-single-groups}.  It is therefore
a product of simple ideals, each mapping diagonally onto some of the
factor ideals.  We call two factor ideals \emph{shared} if they are
images of the same ideal of $\operatorname{Lie}(G_0)^{\rm der}$.

\begin{lemma}[Possible identifications]\label{lem:mixed-shared}
Shared simple ideals have one of the following forms:
\begin{enumerate}
\item orthogonal factors of the same dimension, with their standard
  modules matched;
\item unitary factors of the same dimension, with their standard
  modules matched up to duality;
\item ideals isomorphic to $\mathfrak{sl}_2$, shared among orthogonal
  factors of dimensions three or four and unitary factors of
  dimension two.
\end{enumerate}
Whenever sharing occurs, the fields $F_i$, viewed in $\R$ through
their active embeddings, agree.  The associated place identifications
are then the identity on that common field.  In the unitary case with
$d_i\ge3$, matching the standard modules as representations of the
shared simple ideal also matches their characters on the connected
center of $G_0$.
\end{lemma}

\begin{proof}
An identification of shared ideals must identify their infinitesimal
Hodge-circle elements.  Their eigenvalues first restrict the
possible representations; the Galois translates of the nonzero
components will then identify their fields and places.
First transport a shared ideal by Galois to an active place.  In an
orthogonal standard module the circle element rotates one plane.  For
$\mathfrak{so}_4=\mathfrak{sl}_2\oplus\mathfrak{sl}_2$, both simple
ideals are active, with weights $1,-1$ on their standard spin modules.
The spin standard module for $\mathfrak{so}_3$ has the same weights.
In an active unitary standard module there is exactly one nonzero
weight, equal to $2$ or $-2$; the simple ideal sees its traceless part.

These elements distinguish the possible identifications of the simple
Lie algebras.  In the accidental isomorphism
$\mathfrak{so}_6\simeq\mathfrak{sl}_4$, the orthogonal half-spin
module has weights $1,1,-1,-1$, whereas a unitary standard module
has the traceless pattern of a single nonzero weight.  They cannot
match.  In type $D_4$, a triality automorphism carrying the standard
module to a half-spin module is likewise excluded: the latter has
only weights $1,-1$, unlike the orthogonal standard module.  The
remaining coincidences are precisely those in the statement.

For completeness, the same argument identifies the fields and the
place labels.  Work with simple ideals over $\overline\Q$, testing
activity after the fixed embedding into $\C$.  For an ideal at a
place $p$ of $F_i$, the Galois elements carrying it to an active
ideal are exactly those carrying $p$ to the active place of $F_i$.
This holds also in dimension four, since both ideals there are active.
The sets agree for shared ideals.  After transporting one shared ideal
to the active place, equality of these sets is equality of the
stabilizers of the active field embeddings.  Thus the embedded fields
agree, and all place matchings are the matching embeddings of that
field.  Sharing is transported to every such place; for an orthogonal
dimension-four block this may initially involve only one of its two
simple ideals.

Finally consider a shared unitary ideal of dimension $d\ge3$ and
choose the standard or its dual so that the simple modules agree.
The exceptional eigenvalue in the traceless circle element determines
the sign of its scalar part.  Hence the central characters of the two
matched modules have the same differential on the central component
of the circle.  This is true after every Galois transport.  A Galois
stable character condition that vanishes on the circle vanishes on
the rational torus generated by its projection.  Applying this to
the difference of the two characters proves equality on the connected
center.  One may pass to powers of characters to descend to the
central torus quotient; its character lattice is torsion free, so
equality of these powers implies equality of the original characters.
For $d=2$, the traceless pattern does not distinguish the sign, which
is why that case remains in the third alternative.
\end{proof}

We will use the constituents of $W_i$ as representations of the joint
group.  Restriction of the full even Clifford representation to an
orthogonal sum gives all choices of spin constituents.  In the
orthogonal case, these are tensor products over the places of $F_i$:
there is one spin choice for an odd-dimensional block and either
half-spin choice for an even-dimensional block.  Multiplicities will
not matter.  In dimension three the choice is the standard module of
one $\mathrm{SL}_2$; in dimension four the two choices are the
standard modules of the two $\mathrm{SL}_2$ factors.

In the unitary case the choices at a place are
\begin{equation}\label{eq:mixed-unitary-spin}
 \bigwedge^r\mathrm{Std}\otimes(\det)^{-1/2},
       \qquad 0\le r\le d_i,
\end{equation}
with independent choices at the different places.  The half-determinant
notation denotes a character on the relevant spin cover.  When $d_i=2$,
these choices are the standard module of $\mathrm{SL}_2$, with zero
connected central weight, and two opposite torus characters.  These
descriptions follow from the exterior-algebra model of spin and the
spin decomposition of an orthogonal sum.

\begin{lemma}[Weight-one comparison]\label{lem:mixed-constituents}
If a simple ideal is shared between factors $i$ and $j$, then $W_i$
and $W_j$ contain isomorphic representations of the joint Hodge group
on which that ideal acts nontrivially.  The required intertwiners and
the projections onto individual copies of these representations are
algebraic correspondences with complex coefficients.
\end{lemma}

\begin{proof}
For a shared $\mathfrak{sl}_2$, choose its standard module at the
given place and a transported shared standard module at every other
place.  Lemma~\ref{lem:mixed-shared} provides the common field and
these transported identifications.  In the other orthogonal types,
choose matching spin constituents at all places.  In the unitary
types, choose matching constituents in
\eqref{eq:mixed-unitary-spin}; their central characters agree by
Lemma~\ref{lem:mixed-shared}.  Thus the tensor constituents have the
same joint Lie algebra representation.  The Hodge group is connected,
so they have the same group representation.

Morphisms between weight-one Hodge structures of complex abelian
varieties are induced by rational homomorphisms of abelian varieties.
After extending coefficients, every intertwiner just used is therefore
algebraic.  The same applies to the weight-one idempotents selecting
the constituents and individual multiplicity copies.
\end{proof}

\subsection{Removing projective and torus actions}

We prove the reduction before constructing the auxiliary varieties.
For this subsection, allow any finite collection containing the original
surfaces and their $W_i$ whose Hodge group $G$ projects onto $G_0$ with
finite kernel.  Write $C$ for its algebraic tensor group as above, and
identify the Lie algebras of $G$ and $G_0$.  The auxiliary varieties
constructed below will satisfy this hypothesis.  Thus the reduction
will apply to the final collection once all of its objects have been
fixed.

\begin{lemma}[Projective compatibility]\label{lem:mixed-projective}
For $c\in C(\C)$, its actions at all occurrences of a shared simple
ideal agree in the corresponding adjoint group.
\end{lemma}

\begin{proof}
By Lemma~\ref{lem:mixed-individual}, $c$ acts on each $W_i$ by a
spin lift of its action on $V_i$.  This has the tensor-product
description of its block spin constituents, up to an overall sign.
Lemma~\ref{lem:mixed-constituents} provides matching constituents
on which a specified shared ideal acts nontrivially.  Their
intertwiner is algebraic, so $c$ has the same action in these two
realizations.  Equality of invertible tensor-product operators
forces equality of their individual factors up to scalars.  A
nontrivial representation of a simple Lie algebra is projectively
faithful on its adjoint group.  Hence the adjoint actions at the
shared ideal agree.
\end{proof}

Let $W_{\rm CM}$ be the sum, in all the weight-one realizations of
the collection under consideration, of the constituents on which
$\operatorname{Lie}(G)^{\rm der}$ acts trivially.  This is a
rational sub-Hodge structure.  It is the realization of a CM abelian
factor, and its projector is algebraic.  The CM Hodge conjecture
\cite[Theorem~1.1]{CM}, applied jointly to these factors and their
powers, therefore gives
\begin{equation}\label{eq:mixed-cm-image}
 C\big|_{W_{\rm CM}}=G\big|_{W_{\rm CM}}.
\end{equation}

\begin{lemma}[Reduction to central actions]\label{lem:mixed-central}
After multiplying $c\in C(\C)$ by an element of $G(\C)$, one may
assume that its compatible adjoint actions through the original
factors are trivial at every simple ideal and that it acts as the
identity on $W_{\rm CM}$.  Under these assumptions all
actions on the $V_i$ outside the dimension-four orthogonal blocks
and the dimension-two unitary blocks can be removed by centers of
the shared simply connected derived factors.
\end{lemma}

\begin{proof}
Lemma~\ref{lem:mixed-projective} gives a single prescribed action
in the product of the adjoint groups, and
\eqref{eq:mixed-cm-image} gives the prescribed action on the CM
part.  The derived group covers the first product, and the connected
center covers the torus image on $W_{\rm CM}$.  These choices can
be made independently.  Removing their product proves the first
assertion.  Continue to write $c$ for the resulting element.

On an odd-dimensional orthogonal block with trivial adjoint action,
the standard action is already the identity.  On an even-dimensional
orthogonal block it is a sign.  In a spin lift this sign is the ratio
of the scalar actions on the two half-spin choices.  For simple
types other than $\mathfrak{sl}_2$, these ratios agree at shared
ideals: apply the weight-one intertwiners to matching tensor words
with either half-spin choice at the specified place.  Interchanging
the two choices in one description only inverts the ratio and does
not change a sign.  The center of the shared simply connected
factor supplies that standard-module sign, so it can be removed.

On a unitary block at $(i,p)$, choose the standard or its dual so
that the simple modules match on each shared ideal, as in
Lemma~\ref{lem:mixed-shared}.  Write its scalar action as $r_{i,p}$;
the opposite choice has action $r_{i,p}^{-1}$.  In
\eqref{eq:mixed-unitary-spin}, compare choices $0$ and $d_i$ at
that place, using only these two choices at every other place.
Both resulting words belong to $W_{\rm CM}$, and their ratio is
$r_{i,p}^{d_i}$.  Hence
\begin{equation}\label{eq:mixed-unitary-root}
 r_{i,p}^{d_i}=1.
\end{equation}
This removes the case $d_i=1$.  For $d_i\ge3$, compare also the
choice $1$ with the choice $0$.  The ratio is $r_{i,p}$, and the
matched weight-one constituents show that these ratios agree on
shared simple ideals.  They are therefore supplied by the center
of the shared $\mathrm{SL}_{d_i}$ and can be removed there.
Such derived central corrections do not affect $W_{\rm CM}$.
The only remaining actions on the $V_i$ are therefore signs on
orthogonal dimension-four and unitary dimension-two blocks.
\end{proof}

\subsection{The graph and its auxiliary abelian varieties}

We now formulate the compatibility needed for these residual signs and
construct algebraic tensors that enforce it.  The graph and all its
auxiliary varieties depend only on the original Hodge representations,
before any element of the final group $C$ is chosen.

Form a finite graph $\Gamma$ as follows.  Its ordinary vertices are
the simple ideals of $\operatorname{Lie}(G_0)^{\rm der}_{\C}$
isomorphic to $\mathfrak{sl}_2$.  For every orthogonal factor with
$d_i=4$ and every place $p$ of $F_i$, put an edge $e=(i,p)$ between
the two ideals acting on that block.  For every unitary factor with
$d_i=2$ and every place $p$, put an edge between its simple ideal and
a new vertex $t_{i,p}$, used only for this edge.  We call these new
vertices \emph{torus vertices}.  Orthogonal dimension-three ideals
are included as ordinary vertices, even when no edge meets them.

An edge sign $\varepsilon_e$ will denote the scalar action on its
quadratic block after the preceding reduction.  To remove all these
signs using the shared derived group, we seek signs $s_v$ at the ordinary
vertices, and sign $1$ at each torus vertex, such that
\begin{equation}\label{eq:mixed-sign-target}
 \varepsilon_e=s_us_v
 \qquad(e\text{ with endpoints }u,v).
\end{equation}
Indeed, a central sign in $\mathrm{SL}_2$ acts on its standard module
by that sign.  The orthogonal dimension-four standard module is the
tensor product of the two endpoint modules, while the unitary
dimension-two standard module has just the ordinary endpoint's derived
action.  The same center acts trivially on the orthogonal
dimension-three standard module.

The necessary compatibility is that the edge signs multiply to one
around every cycle and along every path between torus vertices.
These relations are also sufficient: path products from a base vertex
give the signs $s_v$, as proved in Lemma~\ref{lem:mixed-node-signs}.
The role of the auxiliary tensors is to establish precisely these
relations for every normalized tensor automorphism.

We therefore use every simple cycle of $\Gamma$,
including a cycle formed by two parallel edges, and every simple path
with both endpoints torus vertices.  There are finitely many.  Fix one
such path or cycle $\gamma$.  All its edges have the same real field $F$ and
the same place $p$ under Lemma~\ref{lem:mixed-shared}.  For each
embedding $\tau:F\hookrightarrow\C$, choose a Galois transport
$\gamma_\tau$ of the path or cycle from $p$ to $\tau$.  It uses the
same list of original factor indices, although its simple ideals have
been transported.  At the original place choose the identity transport,
so that $\gamma_p=\gamma$.

Viewing each edge as an interval, define the \emph{grid} to be the
Cartesian product of these graphs,
\[
 \mathscr G_\gamma=
       \prod_{\tau:F\hookrightarrow\C}\gamma_\tau.
\]
Its vertices are tuples of vertices, one at each place; a grid step
changes one coordinate along an edge.  A \emph{box} is a product
\(\prod_\tau e_\tau\), where $e_\tau=(I(\tau),\tau)$ is one
chosen edge of $\gamma_\tau$.  Its vertices choose one endpoint of
each edge.  A square face of a box varies the endpoints at two distinct
places and freezes all other vertices.  At any torus endpoint we may
select either of its two character lines; include all such choices in
the finite construction.

\begin{lemma}[Realization of a box]\label{lem:mixed-box}
For each box there is an auxiliary abelian variety with weight-one
realization $W_{\Box}$ and an algebraic tensor
\[
 \mathcal P_{\Box}\in(\End W_{\Box})^{\otimes2}
\]
as in Proposition~\ref{prop:auxiliary-spin}.  At a chosen complex
embedding, $W_{\Box}$ contains the tensor representation attached
to every vertex of the box: the standard two-dimensional representation
at an ordinary vertex and the selected character line at a torus
vertex, tensored over the places of $F$.

Adjoining all these finitely many $W_{\Box}$ to the original system
does not change its Hodge Lie algebra.  Vertex representations with
the same choices, whether occurring in different boxes or in an
original $W_i$, are identified by algebraic weight-one correspondences.
\end{lemma}

\begin{proof}
Choose a totally real Galois closure $F'\subset\R$ of the embedded
field $F$.  Every embedding $q:F\hookrightarrow\C$ then takes values
in $F'$.  For each such $q$, take the four-dimensional quadratic block
of factor $I(q)$ and extend scalars to $F'$ through $q$.
In the unitary case this is the underlying
quadratic space of its two-dimensional Hermitian block.  Let $V'$ be
their orthogonal sum together with sufficiently many negative Tate
directions.

At every real embedding of $F'$, exactly one of the chosen blocks
has a period plane: composition with that embedding permutes the
places $q$, and their active places agree under the common-field
identification.  All other blocks are negative definite.  Thus $V'$
has signature $(2,d'-2)$ at every real embedding.  Choose the number
of negative directions so that $d'$ is sufficiently large and
$d'\equiv5\pmod8$.  The Witt index over $F'$ is at least two.
Indeed the nonarchimedean anisotropic dimension bound and the
Hasse principle first give a hyperbolic plane; the residual form
still has positive and negative directions at every real place and,
in sufficiently large dimension, gives a second hyperbolic plane.

Proposition~\ref{prop:auxiliary-spin} applies to
\[
 W_{\Box}=\operatorname{Res}_{F'/\Q}C^+_{F'}(V')
\]
and supplies its tensor $\mathcal P_{\Box}$.  The spin decomposition
at the chosen embedding gives every endpoint choice in the
four-dimensional blocks.  The extra negative Tate directions supply
only multiplicities, since the Hodge group acts trivially on them.
This proves the assertion about vertex representations.

Here is the rational group action underlying this construction.
Write $Q_i$ for the four-dimensional quadratic space over $F$ of
the edge's factor, and put $Q_{i,q}=Q_i\otimes_{F,q}F'$.  Extension
of scalars through $q$ gives a rational group homomorphism
\[
 \operatorname{Res}_{F/\Q}\mathrm{SO}(Q_i)
 \longrightarrow
 \operatorname{Res}_{F'/\Q}\mathrm{SO}(Q_{i,q}).
\]
Compose these maps with the factor actions of $G_0$, take their
orthogonal sum over $q$, and act trivially on the added Tate
directions.  This defines
\[
 \rho_{\Box}:G_0\longrightarrow
       \operatorname{Res}_{F'/\Q}\mathrm{SO}(V').
\]
The fiber product of $G_0$ with the corresponding spin covers,
using all the maps $\rho_{\Box}$, is a rational group mapping
onto $G_0$ with finite kernel.  The original Hodge circle together
with its auxiliary KS actions lies in this fiber product.
Consequently the enlarged Hodge group projects onto $G_0$ with
finite kernel and has the same Lie algebra.  Matching vertex
choices therefore have the same Lie representation in all their
occurrences.  Connectedness gives the same group representation,
and the weight-one correspondence argument of
Lemma~\ref{lem:mixed-constituents} makes their identifications and
multiplicity projections algebraic.
\end{proof}

Fix all the varieties in Lemma~\ref{lem:mixed-box}, and from now on
let $G,C$ denote the Hodge and algebraic tensor groups of this entire
finite collection.  We identify the Lie algebra of $G$ with that
of $G_0$.  In particular this collection satisfies the hypothesis
of the preceding normalization argument.  The preceding lemmas apply
to this final collection.
For a grid vertex $u$, let $R_u$ be the abstract tensor product of
its endpoint representations.  Identify it with a multiplicity copy
in every occurrence by weight-one Hodge embeddings.  Every
intertwiner between such copies is algebraic by
Lemma~\ref{lem:mixed-box}.  Hence any $c\in C(\C)$ induces the
same operator $c_u$ on $R_u$ through all these identifications.
We may therefore choose multiplicity copies in a box using its
Clifford tensor decomposition without changing $c_u$.

\subsection{The signs of a normalized tensor automorphism}

Now choose $c\in C(\C)$ and normalize it by
Lemma~\ref{lem:mixed-central}, using the CM part of this fixed final
collection.  For an edge $e=(i,p)$ of $\Gamma$,
let
\[
 \varepsilon_e\in\{1,-1\}
\]
be its scalar action on the associated quadratic block of $V_i$.
For a unitary edge this is a sign by
\eqref{eq:mixed-unitary-root} with $d_i=2$, and standard and dual
have the same sign.

At each four-dimensional block, a lift to
$\mathrm{Spin}_4=\mathrm{SL}_2\times\mathrm{SL}_2$ acts on the
two endpoint modules.  The normalized trivial adjoint action makes
each semisimple $\mathrm{SL}_2$ lift central, so its endpoint scalar
is a sign.  In the orthogonal case the two signs multiply to
$\varepsilon_e$.  In the unitary case one endpoint is a torus:
if $a_e\in\{1,-1\}$ is the ordinary endpoint sign and $\lambda_e$
is the selected torus character scalar, then
$a_e\lambda_e=\varepsilon_e$.  Thus $\lambda_e$ and its inverse,
the opposite character scalar, are signs as well.
These endpoint signs may at present
depend on the edge.  On a word in an original $W_i$, the action of
$c$ is the product of the endpoint signs selected at its places.
An overall sign of the spin lift can be absorbed into one block.

These are the edge signs in \eqref{eq:mixed-sign-target}.  We will
prove their cycle and torus-endpoint relations by transporting a sign
from one grid step to every parallel step, independently of the frozen
vertices.

\subsection{Transporting a sign across the grid}

We first isolate a linear-algebra observation that will justify the
subspaces used in the geometric tensor calculation.

\begin{lemma}[Supports of an invariant pairing]\label{lem:mixed-support}
Let $A,B$ be finite-dimensional complex vector spaces, let
$U\subset A$ and $V\subset B$, and let $t\in U\otimes V$ induce
an isomorphism $U^\vee\to V$.  If automorphisms $a$ of $A$ and
$b$ of $B$ satisfy $(a\otimes b)t=t$, then $a(U)=U$ and $b(V)=V$.
If moreover $a|_U=\varepsilon\,\mathrm{id}$ for
$\varepsilon\in\{1,-1\}$, then
$b|_V=\varepsilon\,\mathrm{id}$.
\end{lemma}

\begin{proof}
The images of the contraction maps $B^\vee\to A$ and
$A^\vee\to B$ defined by $t$ are exactly $U$ and $V$.
Thus preservation of $t$ preserves both subspaces.  The last
assertion follows from the perfect pairing, since
$\varepsilon^{-1}=\varepsilon$.
\end{proof}

Fix one of the paths or cycles used to construct a grid.  For a step
along its edge $e=(i,p)$ at place $p$, freeze the vertices at every
other place.  If the local endpoint spaces are $X,Y$ and the tensor
product of the frozen spaces is $T$, the two grid representations are
$X\otimes T$ and $Y\otimes T$.  Inside the Hom space between them,
consider
\begin{equation}\label{eq:mixed-local-hom}
 \Hom(X,Y)\otimes\mathrm{id}_T.
\end{equation}
The reverse direction uses $\Hom(Y,X)\otimes\mathrm{id}_T$.
Neither subspace is being assumed algebraic on its own.

\begin{lemma}[Propagation across a box]\label{lem:mixed-propagation}
For the normalized element $c$, the action on
\eqref{eq:mixed-local-hom}, in either direction, is the scalar
$\varepsilon_e$, independently of the frozen vertices at the
other places of the grid.
\end{lemma}

\begin{proof}
First choose every frozen vertex at an endpoint of the transported
edge of factor $i$.  The corresponding tensor words occur in the
original $W_i$.  Its endpoint-sign description shows that the action
on \eqref{eq:mixed-local-hom} is $\varepsilon_e$ in both directions:
the frozen signs cancel, and the ratio of the two endpoint signs is
their product.

It remains to move one frozen vertex across an edge at a different
place $q$.  Choose a box containing these two steps, with the other
places fixed.  Denote its four vertex representations by
\[
 R_{00}=X\otimes Z\otimes T_0,\quad
 R_{10}=Y\otimes Z\otimes T_0,\quad
 R_{01}=X\otimes Z'\otimes T_0,\quad
 R_{11}=Y\otimes Z'\otimes T_0.
\]
Project $\mathcal P_{\Box}$ to
\begin{equation}\label{eq:mixed-cross-pairing}
 \Hom(R_{00},R_{10})\otimes\Hom(R_{11},R_{01}).
\end{equation}
The second $p$-direction has been reversed.  This projection is
algebraic, since its vertex projections are weight-one Hodge
endomorphisms.  Figure~\ref{fig:mixed-grid} records the two directions.

\begin{figure}[htbp]
\centering
\setlength{\unitlength}{1pt}
\begin{picture}(310,110)
 \put(30,85){\makebox(0,0){$R_{00}$}}
 \put(260,85){\makebox(0,0){$R_{10}$}}
 \put(30,20){\makebox(0,0){$R_{01}$}}
 \put(260,20){\makebox(0,0){$R_{11}$}}
 \put(52,85){\vector(1,0){186}}
 \put(238,20){\vector(-1,0){186}}
 \put(30,66){\vector(0,-1){28}}
 \put(260,66){\vector(0,-1){28}}
 \put(145,99){\makebox(0,0){$p$-step: $\varepsilon_e$}}
 \put(145,4){\makebox(0,0){reverse $p$-step: $\varepsilon_e$}}
 \put(16,52){\makebox(0,0)[r]{$q$}}
 \put(274,52){\makebox(0,0)[l]{$q$}}
 \put(145,52){\makebox(0,0){$\mathcal P_{\Box}$ pairs the local Hom subspaces}}
\end{picture}
\caption{A square face of a box in the grid of transported paths or
cycles.  Only the distinct place coordinates $p$ and $q$ vary.
The projected tensor pairs the local Hom subspace represented by the
upper $p$-arrow with the one represented by the reversed lower arrow.
Preservation of this perfect pairing transports the scalar sign when
the frozen $q$-vertex changes.  The vertical arrows indicate that change
of vertex, not a chosen linear map.}
\label{fig:mixed-grid}
\end{figure}

We verify that the tensor in \eqref{eq:mixed-cross-pairing} is a
perfect pairing between the two local Hom subspaces.  Only the
$p$-block of $V'$ contributes.  Clifford multiplication by a vector
in that four-dimensional block interchanges its two half-spin
spaces and acts as the identity on spectator slots, up to fixed
grading signs.  In a graded Clifford tensor model these signs are
the chirality operators on preceding even-dimensional slots; after
choosing the spectators they are scalars.  The volume element in
the total odd-dimensional Clifford module likewise acts as a
scalar.  Choose multiplicity copies in this Clifford tensor
decomposition and transport the projected tensor to the fixed
tensor models $R_u$.  Its spectator factors are then identities;
the transported actions are the already fixed operators $c_u$.

For the two two-dimensional half-spin spaces $X,Y$, the vector
representation is $X\otimes Y$, and Clifford multiplication gives
all of $\Hom(X,Y)$ and all of $\Hom(Y,X)$.  More explicitly, with
the skew-form identifications $j_X:X\to X^\vee$ and
$j_Y:Y\to Y^\vee$, the two maps belonging to $x\otimes y$ are,
up to nonzero scalars,
\[
 y\otimes j_X(x),\qquad x\otimes j_Y(y).
\]
The inverse quadratic metric therefore gives a nondegenerate tensor
between the two Hom spaces, with identities on the frozen factors.
If an endpoint is a torus character line in $Y$, project to that
line.  The first Hom space consists of maps from $X$ into the selected
line, and the second of maps from that same line to $X$.  In the
displayed formulas these use respectively the selected character
component of $y$ and its opposite under $j_Y$.  They remain perfectly
paired by the metric.  Thus the projected tensor has
exactly the local Hom subspaces as its minimal tensor supports.

Since $c$ preserves $\mathcal P_{\Box}$ and all the vertex
projections, Lemma~\ref{lem:mixed-support} first makes these Hom
subspaces invariant and then transfers the sign $\varepsilon_e$
from the upper one to the lower one.  Reversing the roles of the
two directions gives the reverse propagation as well.  Moving the
frozen places successively along their paths proves the assertion
for every vertex choice.
\end{proof}

\begin{lemma}[Cycle and endpoint relations]\label{lem:mixed-cycle}
The edge signs of $\Gamma$ satisfy
\begin{equation}\label{eq:mixed-cycle-signs}
 \prod_{e\in\gamma}\varepsilon_e=1
\end{equation}
for every cycle $\gamma$ and for every path $\gamma$ whose two
endpoints are torus vertices.
\end{lemma}

\begin{proof}
It suffices to use simple cycles and simple paths, since cancellation
of repeated edges reduces the general case to them.  Fix one such
path or cycle and a place $p$.  Freeze all other places at the
starts of their transported paths.  At the initial grid vertex,
the action of $c$ is scalar: all these endpoint choices occur in
the original representation belonging to the first edge.

Follow the path at place $p$.  Suppose $c$ is scalar $a$ on its
current vertex representation.  By
Lemma~\ref{lem:mixed-propagation}, conjugation on the local Hom
space for the next edge $e$ is $\varepsilon_e$.  Hence
\[
 c_{\rm next}\,f=\varepsilon_e a f
 \qquad
 \text{for every }f\in\Hom(X,Y)\otimes\mathrm{id}_T.
\]
The images of these maps span the next vertex representation.
Thus its scalar is $\varepsilon_e a$.  Around a cycle one returns
to the same representation and obtains
\eqref{eq:mixed-cycle-signs}.  For a torus-to-torus path, both
endpoint words consist only of character slots, including at the
frozen places.  They belong to $W_{\rm CM}$, so their scalars are
both one.  The same calculation proves the endpoint relation.
\end{proof}

\begin{lemma}[Vertex signs]\label{lem:mixed-node-signs}
There are signs $s_v$ at the vertices of $\Gamma$, with $s_v=1$
at every torus vertex, such that
\[
 \varepsilon_e=s_us_v
 \quad\text{for every edge }e\text{ with endpoints }u,v.
\]
\end{lemma}

\begin{proof}
In each connected component choose a base vertex, choosing a torus
vertex when one is present, and give it sign one.  Define the sign
of any other vertex as the product of the edge signs along a path
from the base.  The cycle relations in
Lemma~\ref{lem:mixed-cycle} make this independent of the path.
If another torus vertex occurs, the endpoint-path relation gives
it sign one.  The edge formula follows immediately.  Isolated
vertices may also be assigned sign one.
\end{proof}

\subsection{Completion of the mixed-product argument}

\begin{proof}[Proof of Theorem~\ref{thm:main}]
Take $c\in C(\C)$ for the final collection fixed after
Lemma~\ref{lem:mixed-box}.  Lemma~\ref{lem:mixed-central} removes
an element of $G(\C)$ so that the only remaining actions on the
$V_i$ are the edge signs of $\Gamma$.  By
Lemma~\ref{lem:mixed-node-signs}, these are products of signs
assigned to the shared $\mathfrak{sl}_2$ ideals, with sign one at
the torus vertices.  Use those signs in the centers of the
corresponding simply connected derived factors.  On an orthogonal
dimension-four block their action is the product of the two
endpoint signs; on a unitary dimension-two standard block it is
the ordinary-vertex sign.  Their action on orthogonal
dimension-three blocks is trivial.  They therefore remove every
remaining action on the $V_i$.  This proves
\eqref{eq:mixed-group-goal} over $\C$, and hence over $\Q$.

Every rational Hodge tensor in the $V_i$ is now fixed by $C$ and
is therefore algebraic.  Apply the algebraic K\"unneth projectors
to each factor of $S_1\times\cdots\times S_m$.  Besides the
transcendental summands in degree two, the pieces are generated by
the unit, divisors, and the point class, all algebraic.  Thus every
Hodge class in a K\"unneth summand is a sum of the algebraic
transcendental tensors just obtained tensored with these elementary
classes.  Products and the projectors realize them in the required
codimension.  This proves surjectivity of the rational cycle-class
map for every ordered list, including repetitions, and every
codimension in the statement.
\end{proof}


\part{The realization theorem}
\label{part:realization}

The preceding argument reduces Theorem~\ref{thm:main} to
Proposition~\ref{prop:realization}.  We now prove that proposition.
Its two conclusions, prescribed Euler pairings and a bound on
\(\Ext^2\), will come from one Lagrangian immersion: its fundamental
class controls the pairings, while its second cohomology controls the
extension space.  The construction must also remove the disk
obstruction and identify the category generated by the embedded mirror
tests.  These are the successive tasks of this part.

\section{A Lagrangian immersion with controlled cohomology}
\label{sec:topology}

We begin the proof of Proposition~\ref{prop:realization} by constructing
a geometric representative of its input class $a$.  Two features of that representative will control its
Floer theory: ambient classes detect its degree-two cohomology, and its
self-intersections contribute only odd-degree generators.  We establish
these features together.  All homology and cohomology in this section have
rational coefficients unless otherwise indicated.

\begin{proposition}\label{prop:immersion}
Let $D=Y\times T_g$, $\omega$, and $n$ be as in Section~\ref{sec:setup}, with
$g$ even and $n\geq12$.  Suppose that $a\in H^n(D,\Q)$ satisfies
$\omega a=0$, and that its annihilator
\[
 \mathcal Z=\{x\in H^2(D,\Q):xa=0\}
\]
satisfies the injectivity hypothesis \eqref{setup:wedge-injective}.
There exist a closed connected spin $n$-manifold $L$, a positive integer
$\nu$, and a Maslov-zero Lagrangian immersion $i:L\to D$ such that
\begin{equation}\label{topo:class}
 \PD\,i_*[L]=\nu a,
\end{equation}
\begin{equation}\label{topo:h2}
 i^*:H^2(D,\Q)/\mathcal Z\xrightarrow{\ \sim\ }H^2(L,\Q),
\end{equation}
and
\begin{equation}\label{topo:push}
 i_*:H_{n-2}(L,\Q)\longrightarrow H_{n-2}(D,\Q)
 \quad\text{is injective}.
\end{equation}
The immersion is surjective on fundamental groups and has only transverse
double points, with no triple points.  With an integer grading whose
orientation agrees with the orientation of $L$, both ordered branch
indices at every double point are odd.
\end{proposition}

We first realize the class with the required tangential structure and low
cohomology.  The immersion theorem then supplies a graded Lagrangian
representative.  To control the double-point signs, we place all of its
crossings inside a contractible exact domain and apply the flexibility
results of Eliashberg--Murphy and Ekholm--Eliashberg--Murphy--Smith.  This
adapts the local-domain construction of \cite[Section~3]{A8}; the choice
of filling classes below also treats the nonzero $\pi_2$ of the K3 factor.

\subsection{Tangential realization and low-dimensional surgery}

We seek a representative whose degree-two cohomology is exactly
$H^2(D,\Q)/\mathcal Z$ and whose complexified tangent bundle is
identified with the ambient tangent bundle.  We impose these two
conditions by lifting the representing map: a torus bundle kills
the classes in $\mathcal Z$, and a spin pullback records the required
stable tangent identification.  The following calculation shows that
the lift introduces no further degree-two classes; surgery will then
transfer this degree-two description to the representing manifold.

Choose independent integral classes $z_1,\ldots,z_r$ spanning
$\mathcal Z$ over $\Q$, and let
$\pi_0:\mathscr T_0\to D$ be the principal $r$-torus bundle with these
Chern classes.  Choose the natural stable $\SU$ structure on $TD$, and
form the homotopy pullback
\[
 \mathscr T=\mathscr T_0\times^{h}_{B\SU}B\Spin,
 \qquad \pi:\mathscr T\longrightarrow D,
\]
where $B\Spin\to B\SU$ is complexification.  The tautological stable
real bundle $\xi$ on $\mathscr T$ has virtual rank $n$, is spin, and
comes with a stable identification
\begin{equation}\label{topo:tautological}
 \xi\otimes\C\simeq\pi^*TD
\end{equation}
that retains the determinant trivializations.

\begin{lemma}\label{topo:lift}
Pullback induces
\[
 H^2(\mathscr T,\Q)\simeq H^2(D,\Q)/\mathcal Z.
\]
Moreover, $\PD^{-1}a\in H_n(D,\Q)$ is in the image of
$\pi_*:H_n(\mathscr T,\Q)\to H_n(D,\Q)$.
\end{lemma}

\begin{proof}
The compact K\"ahler manifold $D$ is formal over $\R$
\cite{DeligneGriffithsMorganSullivan}.  We first compute over $\R$.
A real model for its torus bundle is obtained from $H^*(D,\R)$
by adjoining degree-one
generators $u_1,\ldots,u_r$ with $du_j=z_j$.  In degree two, the
horizontal classes give $H^2(D,\R)/(\mathcal Z\otimes\R)$.
The differential on
$H^1(D,\R)\otimes\langle u_1,\ldots,u_r\rangle$ is the scalar
extension of the map in \eqref{setup:wedge-injective}, and is
injective.  The differential
on $\bigwedge^2\langle u_1,\ldots,u_r\rangle$ is also injective:
\[
 d(u_j u_l)=z_j u_l-z_l u_j,
\]
and the $z_j$ are independent.  The images of these two differentials
have different vertical degrees, so cannot cancel.

The fiber of $B\Spin\to B\SU$ is two-connected.  Over $\R$,
complexification takes each even-index Chern generator to a nonzero
multiple of a free Pontryagin generator, and every odd-index Chern
generator to zero.  In the relative Koszul model for the second
pullback, eliminate the even-index pairs.  What remains is the torus
bundle model with generators of degrees $4j+1$, $j\geq1$, whose
differentials are the odd Chern classes of $TD$.  These classes vanish:
$TD$ is the sum of the tangent bundle of a K3 surface, when present,
and a trivial bundle.  In particular the second pullback changes
neither the preceding degree-two calculation nor its pullback map.

Define a functional on the degree-$n$ part of this model by pairing
a horizontal class $v\in H^n(D,\R)$ with $a$, and setting the
functional equal to zero on terms with positive vertical degree.
It is a cycle in the dual complex.  Indeed, a differential with a
horizontal term is a sum of multiples of the $z_j$, and
$a z_j=0$; the remaining new generators have zero differential.
Its induced real homology class maps to $\PD^{-1}a$.
The cohomology pullback and homology pushforward are defined over
$\Q$.  Faithfulness of scalar extension from $\Q$ to $\R$ therefore
descends the asserted kernel and isomorphism, as well as membership
of the rational class $\PD^{-1}a$ in the image of pushforward.
\end{proof}

\begin{lemma}\label{topo:source}
There are a closed connected spin $n$-manifold $L_0$, a map
$f:L_0\to D$, and a positive integer $\nu$ satisfying
\eqref{topo:class}--\eqref{topo:push} with $f$ in place of $i$.
They admit a stable identification
\begin{equation}\label{topo:stable}
 TL_0\otimes\C\simeq f^*TD
\end{equation}
with the specified determinant homotopy.  The map is surjective on
fundamental groups, and $L_0$ has two additional free circle generators
mapping trivially to $D$.  After fixing $\nu$, its Euler characteristic
can be made arbitrarily large while preserving these assertions.
\end{lemma}

\begin{proof}
Use the spin bundle $\xi$ on $\mathscr T$, including
\eqref{topo:tautological}, as a stable tangential structure.
With our virtual-rank-$n$ convention, the group of $n$-dimensional
bordism classes carrying this structure is
$\pi_0\operatorname{Th}(-\xi)$ by the Pontryagin--Thom construction;
see \cite{Thom1954} and \cite[Section~2.3]{RathelFournier2024}.
Rational stable homotopy and the Thom isomorphism give
\[
 \pi_0\operatorname{Th}(-\xi)\otimes\Q
 \simeq H_0(\operatorname{Th}(-\xi);\Q)
 \simeq H_n(\mathscr T;\Q).
\]
Thus a positive integral multiple of the lift in Lemma~\ref{topo:lift}
is represented
by a closed manifold mapping to $\mathscr T$, with its stable tangent
bundle identified with the pullback of $\xi$.  Composing with $\pi$
gives \eqref{topo:stable} and the required fundamental class.  We now
modify the representative by bordisms over $\mathscr T$.

Put $\Gamma=\pi_1(\mathscr T)$.  The torus-bundle homotopy sequence
exhibits its fundamental group as an extension of $\pi_1(D)$ by a
quotient of $\Z^r$; the second pullback has two-connected fiber.
Consequently $\Gamma$ is polycyclic, and $\Gamma\to\pi_1(D)$ is
surjective.  The space $\mathscr T$ has a CW model with finite skeleta
in the degrees needed here: the base and torus bundle are compact,
and the characteristic-class pullback uses classifying spaces and
loop spaces of finite type.  The group ring $\Z\Gamma$ is
Noetherian \cite{Hall1954}.

Join the components of the representative by one-handles and add
loops mapping to generators of $\Gamma$.  Its fundamental-group map
is then surjective.  Its kernel is normally finitely generated,
since both groups are finitely presented.  Surgery on disjoint
embedded representatives of normal generators makes that map an
isomorphism.  Each surgery uses a null homotopy over $\mathscr T$;
a twist of the normal framing on a circle makes the stable tangential
data extend.

Next make the map on $\pi_2$ surjective by adding
$S^2\times S^{n-2}$ summands mapped through $S^2$, for a finite set
of module generators of the missing classes.  Such a pulled-back
spin bundle is stably trivial over $S^2$, so these additions carry
the prescribed tangential data.  The kernel on $\pi_2$ is finitely
generated as a $\Z\Gamma$-module: this follows from the finite
cellular chain modules of the universal covers and Noetherianity.
Represent its generators by disjoint embedded two-spheres and
perform surgery on them, using their null homotopies over
$\mathscr T$.  Their normal bundles are trivial: the spin condition
makes their stable normal bundles trivial, and the rank is in the
stable range.  The stable framing extends over surgery because
$\pi_2\SO=0$.  Since $n\geq12$, the embedding dimensions and the
reverse-handle indices preserve the already established
fundamental-group assertion.  These operations make $\pi_2$ an
isomorphism as well.  The resulting map to $\mathscr T$ consequently
induces an isomorphism on $H^2$.

Now take connected sums with two copies of $S^1\times S^{n-1}$,
and with as many copies of $S^4\times S^{n-4}$ as desired, mapping
these summands constantly to $\mathscr T$.  These products are
stably parallelizable and preserve the tangential structure and
pushed fundamental class.  They do not change $H^2$.  The first
two summands provide the extra null-image circles; each subsequent
summand raises the Euler characteristic by two.  Lemma~\ref{topo:lift}
therefore proves \eqref{topo:h2}.

To verify \eqref{topo:push}, write any $u\in H_{n-2}(L_0,\Q)$ as
$f^*x\cap[L_0]$ using Poincar\'e duality and the surjectivity in
\eqref{topo:h2}.  The projection formula gives
\[
 f_*u=x\cap f_*[L_0]=\nu\,\PD^{-1}(xa).
\]
If this is zero, then $x\in\mathcal Z$ and $f^*x=0$, so $u=0$.
\end{proof}

Since $\omega\in\mathcal Z$, the map just constructed satisfies
$f^*[\omega]=0$.  Destabilize \eqref{topo:stable} in the unitary
stable range over the $n$-dimensional manifold $L_0$, retaining the
determinant homotopy.  Restriction to the real tangent bundle gives
a formal Lagrangian immersion of Maslov class zero.  The
Gromov--Lees theorem realizes it by a Lagrangian immersion homotopic
to $f$ \cite{GromovLees,RathelFournier2024}.  We continue to call
this immersion $f$, and choose its integer grading with the
orientation of $L_0$.

\subsection{Chosen fillings of zero area}

The remaining task is to alter the double points without changing
the cohomological properties.  To do this inside a contractible
exact domain, the loops formed by the two branches of each crossing
must admit chosen fillings of zero symplectic area.  No uniqueness
of a filling area is asserted or required.

\begin{lemma}\label{topo:zero-actions}
After Lagrangian perturbations, $f$ is self-transverse without triple
points, and there are pairwise disjoint embedded arcs in $L_0$
joining the two preimages of each double point such that every
image loop has a chosen filling in $D$ of symplectic area zero.
The perturbations preserve the homotopy class and integer grading.
\end{lemma}

\begin{proof}
Choose fillings for the two extra null-image circles from
Lemma~\ref{topo:source}.  In an immersed cotangent neighborhood,
replacing the zero section by a small closed one-form changes each
chosen filling area by the corresponding period, up to a fixed
sign.  The two periods can be adjusted independently.  Choose
the resulting areas to generate a dense subgroup of $\R$.
Concatenations and reversals of these circles are given the
concatenated and reversed fillings, so their chosen areas add.
Small exact perturbations now make the immersion self-transverse
and remove triple points, without changing these periods.

For a double point $f(a_j)=f(b_j)$, begin with a path from $a_j$
to $b_j$.  Surjectivity on fundamental groups allows insertion of
a source loop that makes its image contractible.  Choose a filling.
By inserting combinations of the two null-image circles with their
chosen fillings, make its area arbitrarily small.  Choose disjoint
source patches at the $a_j$, avoiding every other double-point
preimage.  An exact perturbation by $dF_j$, where $F_j$ is constant
near $a_j$ and zero off its patch, changes this path's filling area
by $F_j(b_j)-F_j(a_j)$, up to the same sign convention.  It fixes
the sheet near the crossing.  Thus it makes the area exactly zero.
The endpoint formula remains valid if the path visits the patch
several times.  With the patches fixed, a uniform sufficiently
small bound on these perturbations preserves all crossings and
creates no new ones; take the preceding residual areas within that
bound.

Finally, homotopies of the source paths relative to their endpoints
preserve the chosen areas, because $f^*\omega=0$ as a differential
form.  General position makes the finitely many paths embedded,
mutually disjoint, and disjoint from the other double-point
preimages.  Grading and spin data transport through the
Lagrangian perturbations.
\end{proof}

\subsection{A contractible exact neighborhood}

\begin{lemma}\label{topo:exact-domain}
There is a compact contractible Weinstein domain
$(P,\lambda)\subset(D,\omega)$, with $d\lambda=\omega|_P$, such that
$W=f^{-1}(P)$ is a disk
containing all double-point preimages.  Its primitive satisfies
$f^*\lambda|_W=0$, and its Liouville field is tangent to the
immersed sheets.  The immersion is cylindrical along $\partial W$.
At one interior point, $W$ has a single linear sheet in a radial
Darboux ball.
\end{lemma}

\begin{proof}
The source and target neighborhoods have different topology.  A tree
joining the double-point preimages has a disk as its source thickening,
but its image acquires a loop at each double point.  We first thicken
that image to a Weinstein zero- and one-handlebody $P_1$.  We then
attach two-handles along its free generators, with cores disjoint from
the immersion.  The target becomes contractible while its preimage
remains a disk.  We give the primitive, collar, and filling details
needed for this construction, adapting \cite[Lemma~3.5]{A8}.


\emph{Thickening the tree and its image.}
Extend the arcs of Lemma~\ref{topo:zero-actions} to an embedded tree
$\Theta\subset L_0$ by connector arcs avoiding all other sheet
intersections.  Its image $G=f(\Theta)$ has precisely the
identifications at the double-point pairs.  Each chosen source arc
becomes one of the free graph cycles, with its previously chosen
zero-area filling.  Subdivide an ordinary edge to obtain a vertex
with only one preimage.
Figure~\ref{fig:topo-spines} illustrates the resulting spines.

\begin{figure}[htbp]
\centering
\setlength{\unitlength}{1pt}
\begin{picture}(420,105)
 \put(10,53){\line(1,0){90}}
 \put(10,53){\circle*{4}}
 \put(35,53){\circle*{4}}
 \put(75,53){\circle*{4}}
 \put(100,53){\circle*{4}}
 \put(10,67){\makebox(0,0){$a_1$}}
 \put(35,67){\makebox(0,0){$b_1$}}
 \put(75,67){\makebox(0,0){$a_2$}}
 \put(100,67){\makebox(0,0){$b_2$}}
 \put(55,19){\makebox(0,0){source tree $\Theta$}}
 \put(55,5){\makebox(0,0){thickening: $W$}}
 \put(119,53){\vector(1,0){31}}
 \put(134,68){\makebox(0,0){$f$}}
 \put(177,53){\circle{32}}
 \put(241,53){\circle{32}}
 \put(193,53){\circle*{4}}
 \put(225,53){\circle*{4}}
 \put(193,53){\line(1,0){32}}
 \put(209,19){\makebox(0,0){image graph $G$}}
 \put(209,5){\makebox(0,0){thickening: $P_1$}}
 \put(271,53){\vector(1,0){31}}
 \put(321,53){\color{black!15}\scalebox{2.1333}{\circle*{15}}}
 \put(385,53){\color{black!15}\scalebox{2.1333}{\circle*{15}}}
 \put(321,53){\circle{32}}
 \put(385,53){\circle{32}}
 \put(337,53){\circle*{4}}
 \put(369,53){\circle*{4}}
 \put(337,53){\line(1,0){32}}
 \put(321,53){\makebox(0,0){$e^2_1$}}
 \put(385,53){\makebox(0,0){$e^2_2$}}
 \put(353,19){\makebox(0,0){filled spine}}
 \put(353,5){\makebox(0,0){thickening: $P$}}
\end{picture}
\caption{Schematic spines for two double points.  Identifying
$a_j$ with $b_j$ turns a contractible source tree into a graph with
loops.  Two-handles kill these loops in the target; their filling
interiors miss the immersion, so the preimage remains the disk $W$.
The picture represents the abstract spines, not an embedding in $D$.}
\label{fig:topo-spines}
\end{figure}

At each vertex take a small radial Darboux ball.  The single sheet,
or the two transverse sheets at a double point, can be straightened
simultaneously.  For two sheets, first make one horizontal; the other
is a gradient graph over the vertical plane, and a symplectic shear
straightens it while fixing the first.  Arrange the incident tree
edges as distinct radial rays.  Thin Weinstein one-handles along the
remaining image arcs join these balls and carry strips in the
corresponding source sheets.

Here are compatible models for these attachments and the two-handles
used below.  In symplectic coordinates $(u,v,z)$, with
$u,v\in\R^j$ and $z\in\R^{2(n-j)}$, use the Liouville field
\begin{equation}\label{topo:handle-field}
 Z_j=(-u,2v,z/2),\qquad j=1,2.
\end{equation}
A box $|u|\leq1$, $|(v,z)|\leq\varepsilon$ has negative face near
$|u|=1$, isotropic core $v=z=0$, and a quadratic Lyapunov function
with one critical point of index $j$.  For $j=1$, the subspace
$v=0$ with $z$ in a fixed Lagrangian plane is an invariant
Lagrangian strip.

The attaching contact identifications can be chosen strict for the
specified contact forms.  Indeed, near an isotropic attaching locus
choose a slice tangent to the contact plane along that locus.
The Reeb flow writes the contact form as $dt+\beta$ on a
neighborhood of the slice.  Match the slices by the symplectic
isotropic neighborhood theorem with the prescribed tangent and
symplectic normal data.  The resulting difference of primitives is
a closed form that restricts to zero on the attaching locus.
It is $dh$, with $h=0$ there, since the neighborhood retracts onto
the locus.  Adjusting the Reeb coordinate by $h$ makes the contact
identification strict.  The same argument applies to the Legendrian
patches in the sheets.  Extending by Liouville flow gives collars
preserving the primitives.

For the one-handles, choose transverse strip framings along the
source arcs and extend the prescribed end germs over each strip.
The relative Lagrangian neighborhood theorem extends these maps
symplectically with the collar germs fixed: after fixing tangent
data, the remaining linear lifts are contractible, and the tubular
construction and Moser argument are relative to the collars.
Shrinking the handles keeps them embedded, mutually disjoint away
from their attachments, and disjoint from every other sheet.
The model primitives glue through the strict collars.  After
rounding the positive corners, we obtain a Weinstein zero- and
one-handlebody $P_1$ with spine $G$.  Its preimage is a thickening
of the source tree, hence a disk, and its primitive vanishes on the
sheets.  The singled-out radial ball can be retained inside $P_1$.

\emph{Isotropic fillings outside the one-handlebody.}
The loops of $P_1$ now have to be killed without meeting any source
sheet.  The fillings must retain their chosen zero-area classes:
contractibility of the attaching circles alone would not control area
in the presence of $\pi_2(D)$.  Choose a free basis of $\pi_1(P_1)$
represented by disjoint isotropic circles in $\partial P_1$, avoiding the Legendrian boundary of the
source disk.  These representatives exist by the relative subcritical
contact isotropic embedding theorem
\cite[Theorem~12.4.1]{EliashbergMishachev2002}; formal tangent lines
can be moved into the contact distribution, and all the avoidance
conditions are in the general-position range.

Each attaching circle inherits a zero-area filling in $D$: join it
in $P_1$ to its graph cycle by an annulus and append the chosen
filling.  Both boundary primitive integrals vanish, so the annulus
has area zero.  We may represent this same filling class outside $P_1$.
In fact, recovering a $2n$-dimensional zero- and one-handlebody from
its boundary uses only handles of indices $2n$ and $2n-1$.
The inclusion of the exterior into $D$ therefore has no relative
homotopy in the dimensions of this disk and its deformation.
This preserves the chosen relative disk class, not just the
contractibility of the attaching circle.  Consequently its
symplectic area remains zero even when $\pi_2(D)\neq0$.
Prescribe the initial collar to be the outward Liouville cylinder.
General position relative to these collars gives embedded filling
disks with disjoint interiors, all missing $f(L_0)$.

We next make the disks isotropic while retaining those collars.
Use a prequantization circle bundle for the integral form $\omega$.
The collars lift horizontally.  In a trivialization over a filling
disk, the zero area gives winding number zero for its boundary
lift, so this lift extends over the disk.  The space of isotropic
ordered two-frames in the contact distribution retracts to
$\mathrm U(n)/\mathrm U(n-2)$ and is simply connected.  Thus the
prescribed isotropic derivative on the collar extends as a formal
isotropic monomorphism over the disk.  The ordinary real two-frame
space in rank $2n+1$ is two-connected, so the derivative of the
lifted disk is homotopic to that monomorphism relative to the
collar.  The relative, $C^0$-dense subcritical contact embedding
$h$-principle realizes these data, fixing a slightly smaller collar.

Projection gives isotropic immersions of the disks into $D$.
They can be chosen embedded and mutually disjoint: contact
Hamiltonian variations near distinct lift points give arbitrary
independent projected variations, and $2+2<2n$.  The same argument
preserves avoidance of the immersed source because $2+n<2n$.
Smallness protects avoidance away from the fixed collars; the
collars themselves protect the attachments.

\emph{Attaching the two-handles.}
We now have embedded isotropic disks with disjoint interiors outside
$P_1$, prescribed outward collars, and no intersections with the
immersion.  Choose symplectic normal framings over these disks.  They
restrict to the contact symplectic normal framings needed at the attaching
circles.  Parameterize the attaching circles by parameterizations
extending over the disks.  The relative isotropic neighborhood
theorem and relative Moser argument now embed thin index-two
handles with the prescribed strict collars and model
\eqref{topo:handle-field}.  These handles miss the immersion.
Their primitives glue, and rounding corners produces $P$.

Its spine is the graph $G$ with a two-cell attached along each
free generator.  After collapsing a maximal tree, these attaching
maps form a free basis of a bouquet of circles; the resulting
complex is contractible.  Thus $P$ is contractible.  The two-handles
have not changed the preimage disk, its primitive-zero sheets,
or its outer cylindrical collar.  The radial single-sheet ball
also remains intact.
\end{proof}

We have now localized all double points in an exact simply connected
target.  This is what permits their algebraic signed count, rather
than a fundamental-group-valued intersection invariant, to control
the final replacement.

\subsection{Double-point signs and the integer grading}

\begin{proof}[Proof of Proposition~\ref{prop:immersion}]
Write $n=2k$.  For the immersion supplied by the preceding lemmas,
the self-intersection formula reads
\begin{equation}\label{topo:double-number}
 2I(f)=f_*[L_0]\cdot f_*[L_0]-(-1)^k\chi(L_0).
\end{equation}
Here the normal Euler number is $(-1)^k\chi(L_0)$, since the
normal bundle of an oriented Lagrangian is its tangent bundle
with this orientation factor.  In Lemma~\ref{topo:source}, choose
the Euler characteristic large enough that $I(f)$ is nonzero
and has sign $(-1)^{k-1}$.  Put $d=|I(f)|$.

Complete the domain $P$ of Lemma~\ref{topo:exact-domain} by its
positive cylindrical end, extend $W$ by its cylinder, and remove
the center of its single-sheet radial ball.  The punctured ball
is a negative Liouville end.  The source is now
$S^{n-1}\times\R$, with the standard real-plane Legendrian
sphere $\phi_0$ at that end.  The target is simply connected:
$P$ was contractible, and removing a point in dimension $2n$
does not change its fundamental group.

The stabilization traces of
\cite[Lemma~3.2 and proof of Lemma~3.4]{EEMS} have signs
\begin{equation}\label{topo:trace-signs}
 \phi\longrightarrow\phi_1:\ (-1)^{k-1},
 \qquad
 \phi_1\longrightarrow\phi:\ (-1)^k,
\end{equation}
where $\phi_1$ is a loose stabilization and the arrows run
outward along the symplectization.  Insert $d$ reverse traces
in separated collars near the puncture, from the $d$-fold
stabilization $\phi_d$ at the negative end to $\phi_0$ outward.
The new negative link is connected and loose, and its total
signed double count is zero by \eqref{topo:double-number}.

Eliashberg--Murphy's exact cancellation theorem
\cite[Theorem~2.3]{EliashbergMurphyCaps} now applies.  The target
is a simply connected Liouville manifold with a negative end;
the immersion is proper, cylindrical there, self-transverse,
and has finitely many double points.  It is exact because
$H^1(S^{n-1}\times\R;\R)=0$.  The loose component and vanishing
signed count were just arranged.  The theorem's extra width
condition concerns source dimension three, so is irrelevant
here.  We obtain a compactly supported Hamiltonian regular
homotopy to an embedding, leaving both cylindrical tails fixed.

In the unchanged negative tail insert $d$ forward traces from
$\phi_0$ to $\phi_d$, and fill the center by the real plane.
The resulting disk immersion has exactly $d$ double points,
all of sign $(-1)^{k-1}$.  Its noncylindrical part is compact.
A sufficiently long negative Liouville flow in the filled
positive completion brings that part inside $P$.  Backward
flow exists for all time; points of $P$ cannot leave it
backward, while the positive end is translated toward $P$.
The flow rescales the symplectic form and hence preserves
Lagrangianity.  The cylindrical set is unchanged, so its
intersection with the original outer collar is unchanged.
Trim at that collar and glue to the immersion outside $W$.

This replaces a source disk by a source disk with the same
collar.  A collar identification extends continuously over
the disks, and their maps into the contractible $P$ are
homotopic relative to the collar.  It follows that the pushed
fundamental class and the maps on degree-two and degree-$(n-2)$
(co)homology are unchanged, under the identifications through
the common complement.  The same complement carries every
first homology class and the fundamental-group map.
The boundary sphere has a unique spin structure, so spin extends
over the replacement disk.  The Maslov class remains zero.
Choose the extended integer grading with the
original orientation parity.

For graded oriented Lagrangian intersections in real source
dimension $2k$, a generator of degree $|x|$ has intersection
sign
\[
 (-1)^k(-1)^{|x|}.
\]
The sign $(-1)^{k-1}$ obtained above therefore makes every
branch degree odd.  Reversing the ordering replaces its degree
by $n-|x|$, which is also odd because $n$ is even.  The disk
replacement has only transverse double points and introduces
no triples or intersections with the unmodified source.
Together with Lemma~\ref{topo:source}, this proves all assertions.
\end{proof}

\section{Removing the disk obstruction}\label{sec:curvature}

Let $i:L\to D$ be the graded spin immersion of
Proposition~\ref{prop:immersion}.  Its transverse self-intersection
generators have odd integer degrees, and
\begin{equation}\label{curv:topological-input}
 H^2(L,\Q)\simeq H^2(D,\Q)/\mathcal Z,
 \qquad
 i_*:H_{n-2}(L,\Q)\longrightarrow H_{n-2}(D,\Q)
 \text{ is injective}.
\end{equation}
We will use the second assertion to remove the curvature of its
Fukaya algebra.  The first will then bound its degree-two endomorphisms.
The geometric ingredient is a Stokes identity for disks with one
interior marked point and no external boundary marked point.  We
construct that identity geometrically, then use filtered transfer to
remove curvature.

\subsection{Branes, gradings, and the filtered construction}

Choose countably many Hamiltonian copies of the product test branes,
ordered by the positive integers, with each test occurring cofinally.
Use product Hamiltonian perturbations, small enough that the torus
factors remain perturbed graphs.  Arrange pairwise transversality and
absence of triple intersections, including intersections with $L$.
When both factors are present, impose these conditions separately on
the factors of all embedded copies.  Hamiltonian variations supply the
required tangent directions, so these are countably many generic
conditions.

Fix a split compatible almost complex structure $J_0$ with no
nonconstant spheres and no nonconstant disks on any embedded test
copy.  The sphere and disk exclusion on the quartic factor is the
genericity assertion used in the test-brane construction; on the torus
factor it follows from the vanishing of the relevant second relative
homotopy groups.  No disk exclusion is imposed on $L$.  All category
constructions in this section involve a finite subcollection of the
copies together with $L$.

Write $\Lambda_0$ for the universal Novikov ring with complex
coefficients and $\Lambda$ for its fraction field; its monomials are
written $T^\lambda$, with valuation $\lambda$.  We use the immersed
cyclic de Rham construction of Abouzaid--Fukaya--Oh--Ohta--Ono
\cite[Theorems~9.3 and~9.6]{AFOOO2026}.  It gives a curved cyclic
filtered $A_\infty$ category with homotopy units for finite collections
of closed spin immersions in a compact symplectic manifold.  Here the
relative spin background and bulk class are zero.  The morphism spaces
are forms on the source diagonals and the generators at transverse
intersections.  We denote the operations by $m_k$, their cyclic pairing
by $\langle\ ,\ \rangle$, and the curvature by $m_0$.

These operations have integer degree $2-k$.  Indeed, the brane
gradings specify integer indices for the cap paths at every transverse
end, including self-intersection ends.  After capping, the boundary
grading has Maslov index zero.  There is no ambiguity from spheres
because the ambient first Chern class vanishes.  The same index
calculation applies to parameterized operations.  Energy monomials
have degree zero, and all contractions, transfers, and comparison
maps below are chosen homogeneous.  Thus the grading is retained by
the construction itself, including the shifted degrees or auxiliary
variables used in a particular model.

We use Kuranishi structures with normalized boundaries and corners,
outer collars where required, and continuous-family perturbations
(CF perturbations).  The local chart theorem and its evaluation
submersion criterion are
\cite[Theorem~7.1, Condition~7.11, and Lemma~7.12]{FOOOChartsI};
compatible systems arise from disk-component-wise obstruction data
\cite[Theorems~5.3 and~5.4]{FOOOChartsII}.  We use the corrected chart
conventions of \cite{FOOOChartsICorrigendum}.
The analytic gluing input is the exponential control of the gluing
maps and their length derivatives
\cite[Assumption~3.12 and Theorems~3.13, 6.4, 8.16, and~8.17]{FOOOGluing}.
For transverse immersed ends, the fixed-jump charts and spin
determinant cappings are those of
\cite[Sections~4--5]{AkahoJoyce}, as incorporated in the immersed
construction of \cite{AFOOO2026}.

In particular, the nodal chart uses component surjectivity augmented
by matching equations at ordinary boundary nodes and by the transverse
asymptotic conditions at jumps.  Auxiliary markings stabilize the
domains, and the obstruction spaces are transported on compact cores
before gluing.  These are the local inputs needed below when the total
curve has no external boundary mark.  All the additional spaces used
for that purpose have one interior mark and no sphere
components; there is no gluing at an angular interior node.

The construction is first made through finite energy and arity
cutoffs.  Filtered pseudo-isotopies, homotopy limits, and canonical
models then give the complete category
\cite[Sections~6, 9--10, and~12, in particular Propositions~9.55
and~12.10]{AFOOO2026}.  We keep $J_0$ fixed in this passage.  The
following observation gives a common energy filtration for a finite
list of branes.

\begin{lemma}[Energy values]\label{curv:energy}
For a finite list of the immersed and embedded branes above, the
energies of $J_0$-holomorphic polygons form a locally finite subset of
$\R_{\geq0}$.  In particular, the nonconstant polygons have a common
positive energy lower bound, independent of their number of marked
points.
\end{lemma}

\begin{proof}
Let $K\subset D$ be the union of their images.  It has the homotopy
type of a finite complex and has arbitrarily small regular
neighborhoods: away from the finitely many double points these are
ordinary tubular neighborhoods, and at a double point they are the
standard transverse-double neighborhoods.  Relative de Rham theory
therefore represents every class in $H^2(D,K;\R)$ by a smooth closed
form that vanishes near $K$.

The symplectic form has a relative lift that computes the polygon
areas.  To see this directly, construct near $K$ a primitive $\alpha$
of $\omega$ whose pullback to every sheet is zero.  At a double point,
simultaneous straightening of the transverse Lagrangians gives a radial
primitive with this property.  On the remaining pieces use the
Lagrangian neighborhood theorem.  After shrinking the neighborhoods,
their overlaps at the ends of the pieces are collars on disjoint
sheets.  The differences of the primitives are exact there, with
potentials zero on the sheets, by the relative Poincar\'e lemma.
These potentials patch the primitives.  Cutting off $\alpha$ gives a
closed form $\widetilde\omega$ vanishing near $K$ and representing the
relative lift.  Stokes' theorem, with the usual decay at jump ends,
gives
\[
 \int u^*\widetilde\omega=\int u^*\omega
\]
for every polygon $u$.

Choose closed forms vanishing near $K$ that represent a basis of
$H^2(D,K;\R)$.  Their integrals over a $J_0$-holomorphic polygon are
bounded in absolute value by fixed constants times its area: this is
the pointwise comass estimate for the compatible metric.  A polygon
defines an integral class in $H_2(D,K;\Z)$, whose quotient by torsion
is a lattice.  The preceding bounds leave only finitely many lattice
points below a fixed area cutoff.  Pairing with
$[\widetilde\omega]$ leaves only finitely many energies.  Energy
positivity now gives the asserted lower bound whenever a nonconstant
polygon exists.  If none exists, any positive lower bound will do.
\end{proof}

The additive monoid generated by these energies is locally finite as
well: below an energy bound $B$, each summand belongs to a finite set
and a sum has at most $B/\delta$ positive summands, where $\delta>0$
is the gap.  This is the filtration used for the transfer operations.

At energy zero, the unary operation is the de Rham differential on
the source diagonals and zero on transverse-point generators; the
constant strip with no stabilizing mark is unstable.  Transfer at
energy zero consequently gives the graded canonical space
\begin{equation}\label{curv:canonical-space}
 C_L^{\mathrm{can}}
 =H^*(L;\C)\otimes\Lambda_0
   \;\oplus\!\!\bigoplus_{x\text{ a branch generator}}\!\Lambda_0 x.
\end{equation}
The diagonal transfer can be made cyclic using a compatible de Rham
contraction.  Homotopy units will give cohomological units after
curvature is removed.  Since every branch generator has odd degree,
the degree-two part of this canonical space is exactly
$H^2(L;\C)\otimes\Lambda_0$.  Removing curvature will therefore
give the cohomological bound in the following proposition, which is
the output of this section.

\begin{proposition}\label{prop:curvature}
For every finite subcollection of the embedded test copies, $L$ and
those copies define an uncurved proper cohomologically unital cyclic
$A_\infty$ category over $\Lambda$.  Its split-closed triangulated
hull has Serre functor $[n]$, and
\begin{equation}\label{curv:degree-two-bound}
 \dim_\Lambda \Hom^2(L,L)
 \leq b_2(L)=b_2(D)-\dim_\Q\mathcal Z.
\end{equation}
For a sequence of embedded copies with strictly increasing labels,
the operations are the geometric disk operations at $J_0$, up to
filtered pseudo-isotopy through every finite cutoff.
\end{proposition}

The proof occupies the rest of the section.  To remove curvature, we construct scalar operations on disks with
an interior root, relative to the ordinary cyclic disk operations.
Their Stokes identity detects the leading curvature class against
ambient cohomology and remains valid after positive-valuation
boundary insertions.  Filtered transfer then gives the proposition.

\subsection{A Stokes identity with one interior mark}

Fix the collared componentwise Kuranishi data and CF-perturbations
defining the cyclic boundary-rooted disk operations through a finite
energy and arity cutoff, enlarged to contain the factors of all
boundary strata that will occur.  We construct additional scalar
operations relative to these geometric choices.  Their disks carry one distinguished
interior mark and cyclically ordered boundary inputs, but no boundary
output.  The component containing the interior mark roots a broken
disk tree.  Every other component has its ordinary boundary output
at the node pointing toward that root.

Let $\eta$ be a closed smooth $(n-2)$-form on $D$.  Pulling back $\eta$
at the interior mark and the inputs at the boundary marks, and
integrating, gives functionals
\begin{equation}\label{curv:bulk-derivative}
 p_k^\eta(x_1,\ldots,x_k)
 =\sum_\beta T^{\omega(\beta)}p_{k,\beta}^\eta(x_1,\ldots,x_k),
 \qquad k\geq1.
\end{equation}
We include all words of diagonal and transverse branch sectors, with
their capping orientation lines.  The signs use the suspended degrees
$|x|'=|x|-1$.  Thus a cyclic permutation of the inputs has its usual
Koszul sign in these suspended degrees.  On a continuous-lift sector
the virtual dimension of the underlying disk problem is
\begin{equation}\label{curv:unmarked-dimension}
 \operatorname{vdim}\mathcal P_{k,\beta}=n+k-1.
\end{equation}
The capped index gives the corresponding formula on a word containing
branch sectors.  The operation is zero unless its integrand has top
degree.  Spaces with $k=0$ are used for Stokes; the unstable constant
disk with just the interior mark is excluded.

\begin{lemma}[Interior-rooted virtual integration]\label{curv:stokes}
Relative to the fixed cyclic boundary-rooted operations, these spaces
admit oriented virtual integrations through every finite cutoff.
Their normalized boundary products have an interior-rooted parent
and an ordinary boundary-rooted child, matched either over the source
diagonal or at reciprocal branch sectors.  The integrations are
cyclic in their external inputs and obey Stokes and composition.
In particular,
\begin{equation}\label{curv:pairing-identity}
 p_1^\eta(m_0)=0,
 \qquad
 p_{1,0}^\eta(c)=\pm\int_L c\wedge i^*\eta
 \quad(c\in\Omega^2(L)).
\end{equation}
The sign in the second identity is fixed and nonzero.  The ordinary
disk operations retain their chosen values through the cutoff.
\end{lemma}

\begin{proof}
We construct the interior-rooted spaces with all immersed boundary
labels.  Their local boundary problems are the ordinary diagonal and
transverse-jump problems of \cite[Sections~4--5]{AkahoJoyce}; the
relative construction below specifies their organization when no
boundary output is chosen.  It uses the mixed parent--child
organization of \cite[Propositions~16.3 and~16.14]{AFOOO2026}, retaining
reciprocal internal switches even when the external word is empty.
The parent is always additional interior-rooted data.  Its
identification with a cyclic boundary-rooted operation is neither
required nor imposed.

\emph{Compactification and normalization.}
Keep the interior mark in immersed stable-map compactness.  At a
boundary node the two lifts either agree in $L$ or end at distinct
preimages $p,q$ of a double value.  In the latter case the two
normalized flags have the reciprocal labels $(p,q)$ and $(q,p)$.
These flags disappear from the external word on gluing; they remain
part of the normalized stratum.  In particular, an empty external
word does not restrict its compactification to the diagonal sector.

There are no nonconstant spheres.  A connected tree of constant
spheres attached at one interior node and carrying at most the one
interior mark cannot be stable: it has fewer than three external
special points in total.  The degeneration in which the entire
boundary collapses would retain a nonconstant sphere, and is also
absent.  All remaining strata are disk trees.  Constant stable
vertices are retained, including the disk with the interior mark
and one boundary node.  The constant disk with no boundary input or
node and only the interior mark is unstable and gives no stratum.

Let $v_*$ be the vertex carrying the interior mark.  Orient every
edge toward $v_*$.  A vertex other than $v_*$ then has exactly one
outgoing flag, which is its ordinary disk output; every remaining
flag is an input.  A normalized corner is the product of the moduli
problems on these vertices, with the diagonal matching equations
and reciprocal-label matches along its edges.  For any collection
of edges, cutting them in any order produces this same oriented
tree and the same factors.  The root of each separated component
is either $v_*$ or its unique edge toward $v_*$.  This rule also
describes intermediate faces, so the identifications on iterated
normalized corners agree.

Stabilization respects these identifications.  Forgetting auxiliary
marks contracts only constant unstable components.  Contracting a
non-root such component carries the unique path toward $v_*$ to the
unique path toward $v_*$; hence the output flag on every surviving
ordinary component is unchanged.  A constant component carrying the
interior mark and an attaching boundary node is stable and cannot
be contracted.  If it has no such node, it is the excluded constant
zero-input root.  An unstable constant bivalent disk has either two
diagonal nodes on the same sheet or two reciprocal switch nodes:
its two constant boundary arcs rule out exactly one switch.  Its
contraction preserves the labels of the surviving flags.  The
two-switch constant strip is the identity under the dual cap-line
identification fixed by determinant gluing; it is not a new
orientation choice \cite[Sections~4.3 and~5]{AkahoJoyce}.
Thus a contraction performs the same pairing
of the ordered cap lines whichever order of normalization is used.
For a chain of these components, the neck lengths concatenate as
$T=T_1+\cdots+T_j$.  This is associative.  Forgetting is interpreted
in the admissible category of \cite[Section~14.3]{AFOOO2026}:
its pullbacks of admissible obstruction and CF data are admissible
and smooth in the gluing profiles used for integration.  The
forgetting map itself need not be smooth in ordinary corner
coordinates.  This admissible compatibility gives compatibility
with stabilization as well as with cutting a fixed tree.

These constructions are equivariant under isomorphisms of decorated
trees.  Such an isomorphism fixes the interior mark and preserves
the directed paths to it, so it never changes an ordinary child's
output.  Relabeling external inputs permutes their positions without
rerooting any ordinary child.  The same statement holds for an
automorphism permuting several isomorphic subtrees.  Normalized
corners retain their ordering of the cut edges; changing that
ordering permutes the outward normal lines with their determinant
sign.  Their subsequent finite quotients therefore have the required
corner identifications and automorphism weights.

\emph{Relative charts and smoothing.}
At a smooth nonconstant root choose auxiliary marked points on
compact cores, in free orbits of its finite automorphism group,
and impose transverse incidence conditions.  Choose obstruction
sections supported away from the distinguished mark, the nodes,
and the strip ends.  Their span can make the augmented
Cauchy--Riemann operator surjective: a cokernel vector annihilating
all such sections would vanish on an open core and then vanish by
unique continuation.  Taking the span of the stabilizer translates
makes this an invariant finite-dimensional obstruction space.
If effective charts are used, include vectors $D_u\xi$ for map
variations $\xi$ supported on disjoint free stabilizer-orbit balls
in the core.  These variations lie in the augmented kernel and
carry a faithful stabilizer representation modulo infinitesimal
reparametrizations; they make the thickened chart effective.
No submersivity is requested at the interior evaluation, which is
used only to pull back $\eta$.

At a broken root use the direct sum of the obstruction data on its
normalized components.  Ordinary vertices carry the fixed ordinary
data.  At a diagonal edge the child output evaluation is submersive,
so component surjectivity remains surjectivity after imposing the
matching equation in $L$.  All matching equations can be solved
recursively from the root toward the leaves: having chosen the
parent variation, choose the child's variation with the required
output value and continue with its children.  Joint submersivity
of all boundary evaluations on one component is not needed.
At a reciprocal switch the matching
space is a point; the operator has precisely the transverse
fixed-jump asymptotic condition of the immersed disk charts.
The linear and nonlinear gluing problems are therefore those of
the ordinary immersed construction, with the interior-mark
coordinate added on the root component.  A constant root
stabilized by that mark is kept as a marked component in this
description.

Transport obstruction sections from the compact cores into the
plumbed domains.  Exponential gluing estimates give charts through
the smoothing parameters, with the corrected smooth gluing
profiles cited above.  Coordinate changes are obtained from
componentwise obstruction data.  We specify their relative choice,
so that compatibility is more than a choice of unrelated local
thickenings.  Use the quasi-component construction of
\cite[Conditions~8.12--8.15, Lemmas~8.19--8.20, and
Definition~8.21]{FOOOChartsII}.  A quasi-component records a
transported compact-core obstruction choice and its auxiliary
stabilization data.  On a boundary tree take the componentwise
union of the already chosen indices: those of its interior root
and the fixed indices of its ordinary vertices.  The directed-tree
identifications show that this prescription agrees on every
iterated boundary and stabilizer quotient.

Extend the corresponding open family of indices inside a proper
family over a sufficiently small boundary neighborhood, as in
\cite[Proposition~8.18]{FOOOChartsII}.  Complete its finite cover
by new compact-core choices on the root away from a smaller
neighborhood, and saturate the choices under the finite
automorphism groups.  Choose the new compact-core spaces generically
transverse to the finitely many inherited spaces and their
stabilizer translates, as in the quasi-component construction.
Surjectivity and these direct-sum conditions are then open;
shrinking the neighborhoods preserves them.
Properness gives the semicontinuity inclusion of obstruction
spaces under changes of the central map, as in
\cite[Proposition~8.17]{FOOOChartsII}.  These inclusions and
forgetting the auxiliary marks define the coordinate changes;
they satisfy the cocycle conditions because the index inclusions
compose and the admissible forgetting operations commute with
normalization.  The original boundary families are unchanged on
the smaller collar.  This is the componentwise proof of
\cite[Theorems~5.3 and~5.4]{FOOOChartsII}, with its additional root
data specified relatively.  The local diagonal and fixed-jump
operators and the augmented matching surjectivity verified above
are exactly what that construction uses.  The fixed ordinary
factors are retained.

\emph{Finite extension of perturbations.}
Let $\delta>0$ be a lower bound for positive component energies,
and choose $0<\varepsilon<\delta$.  Assign cost
$E+\varepsilon(k-1)$ to an ordinary vertex of energy $E$ with $k$
inputs and cost $E+\varepsilon k$ to an interior-rooted vertex.
These costs add under boundary splitting.  Every ordinary stable
vertex has positive cost: the zero-energy unary strip is unstable
and represents the separately fixed de Rham differential, and a
zero-energy stable ordinary vertex has at least two inputs.  The
constant one-input interior root has cost $\varepsilon$; the
constant zero-input root is excluded.  Each nontrivial breaking
therefore has lower-cost factors.  The locally finite energy set,
the arity cutoff, and stability give a finite collection of
problems and normalized corner types, after adjoining all required
faces.  This includes constant switching polygons; their number of
inputs is controlled by arity, not by their zero energy.

Induct on this finite order.  On every normalized boundary prescribe
the product of the already constructed interior-rooted parent
perturbation and the fixed ordinary child perturbation.  The
prescriptions agree on higher corners by the directed-tree and
stabilization identifications.  Diagonal fiber products are defined
because the child output is strongly submersive for its fixed
CF-perturbation; reciprocal switches are products over points.
First extend these compatible normalized-boundary products to a
corner neighborhood by outer collaring, as in the collar step of
\cite[Proposition~14.8 and Remark~14.7]{AFOOO2026}.  The required
agreement on iterated normalized boundaries is precisely the
agreement verified above, including their stabilizer quotients.
Apply \cite[Proposition~14.5]{AFOOO2026} with the map to a point
and with the closed smaller corner neighborhood as its prescribed
set.  It supplies a compatible thickening and extends the
CF-perturbation over the new root space while retaining it near
that set.  Composition on the prescribed collar still uses the
original ordinary child perturbations.  There is no requirement
to forget a boundary input on the new root space and no new
submersion condition at its interior mark.

For cyclic relabeling and finite stabilizers, take invariant
obstruction spaces and include the finite-group representations in
the auxiliary parameter spaces of the CF-perturbations.  The
prescribed corner products are equivariant, so the same relative
extension is equivariant.  This uses equivariant families, not
single-valued equivariant transverse sections.  Finite quotients
are integrated with their ordinary stabilizer weights.  Stokes and
composition now apply to every prescribed product.  The identical
construction over a parameter interval, relative to its endpoints,
supplies finite-cut comparisons whenever the ordinary choices vary.

\emph{Orientations and the zero-input identity.}
Use the spin determinant orientations and the graded cap lines at
all transverse flags.  At a reciprocal node the two cap lines are
paired in their ordered duality; at a diagonal node they are
replaced by the matching determinant of $TL$.  Determinant gluing
is associative.  Permuting normalized smoothing normals or cap
lines gives exactly their Koszul signs, including at an
automorphism quotient.  Stabilizing marks together with their
incidence constraints contribute canceling orientation factors.
These are the orientation conventions for ordinary immersed disk
composition; adding the even-dimensional interior-mark coordinate
does not alter its cap-line rule.

For $k=0$, pull back the closed $(n-2)$-form $\eta$ to the virtual
$(n-1)$-space.  Each boundary product has a single ordinary
curvature child and an interior-rooted parent with one input.
Composition first identifies its contribution with
$p_1^\eta(m_0)$, with the common Stokes sign.  This includes every
reciprocal-switch product.  Only after composition do we use the
grading: a curvature output has degree two, while each branch
generator has odd degree, so those child outputs vanish.
There is no sphere or boundary-collapse term.  Stokes gives the
first identity in \eqref{curv:pairing-identity}.

At energy zero the one-input diagonal root is the regular constant
disk space $L$.  Its one boundary mark and one interior mark fix
the domain; its evaluations are the identity and $i$.  Its integral
is the second identity in \eqref{curv:pairing-identity}.  The
boundary-rooted operations were never changed in this construction.
\end{proof}

\subsection{Leading obstructions and canonical models}

For a degree-one cochain $b$ of positive valuation, put
\[
 m_0^b=\sum_{k\geq0}m_k(b,\ldots,b),
 \qquad
 R_\eta^b(c)=\sum_{r\geq0}p_{r+1}^\eta(c,b,\ldots,b).
\]
These sums converge in the energy filtration.  The cochain $b$ may
have both diagonal and degree-one branch components.  Through a
fixed cutoff only boundedly many insertions occur.

\begin{lemma}[Boundary deformations]\label{curv:deformation}
For the fixed ordinary operations and the relative integrations of
Lemma~\ref{curv:stokes},
\begin{equation}\label{curv:deformed-pairing}
 R_\eta^b(m_0^b)=0.
\end{equation}
If $m_0^b$ has a first nonzero energy coefficient, that coefficient
is a closed exact two-form on $L$.
\end{lemma}

\begin{proof}
We first establish the identity using the complete virtual boundary
products.  Apply Stokes to the interior-rooted spaces with $k$
external inputs all equal to $b$.  A boundary child carrying $s$
consecutive inputs has output $m_s(b,\ldots,b)$, and the root carries
the remaining $r=k-s$ inputs and this output.  Differentiating an
input form supplies the zero-energy differential term in $m_1$.
There is no differential of $\eta$ and no boundary output
differential.

The orientation signs are the suspended insertion signs.  At a
diagonal matching node this is the ordinary disk Stokes calculation,
with the even interior mark in place of the root's boundary-output
coordinate.  External branch cap lines are moved with their total
shifted degrees in the same calculation.  Since every input $b$
has suspended degree zero, all insertion signs are $+1$; rotating
the degree-two curvature output past these inputs also has sign
$+1$.  This applies to degree-one point sectors as well as to
one-forms.  Reciprocal child outputs are retained in the product
formula, and then vanish because they would have degree two in an
odd-degree point sector.

For $k>0$ there are $k$ cyclic positions for each contribution.
For $0<s<k$ they are the starts of the child block; for $s=0$ they
are the $k$ gaps; and for $s=k$ they are the $k$ gaps at the child's
attaching node, or equivalently its rooted linear orders.  The
input differentials also have $k$ positions.  These are counts on
normalized labeled strata, with stabilizer weights retained; we
do not quotient the spaces by the external cyclic action.
Dividing Stokes by $k$ thus gives
\[
 \sum_{r+s=k}
 p_{r+1}^\eta\bigl(m_s(b,\ldots,b),b,\ldots,b\bigr)=0.
\]
The separate $k=0$ identity is $p_1^\eta(m_0)=0$, proved above.
Summing over $k$ proves \eqref{curv:deformed-pairing}.  Through any
cutoff the sum is finite, so all these operations and divisions are
legitimate.

Let $c_E$ be the first nonzero coefficient of $m_0^b$.  Grading
places it in $\Omega^2(L)$, since no branch sector has degree two.
The leading curved $A_\infty$ equation is $dc_E=0$.  Every term
of $R_\eta^b$ involving a $b$ insertion or positive disk energy has
positive valuation.  Its remaining term is
$c\mapsto\pm\int_L c\wedge i^*\eta$.  The coefficient of energy
$E$ in \eqref{curv:deformed-pairing} is therefore
\begin{equation}\label{curv:ambient-annihilation}
 \int_L c_E\wedge i^*\eta=0.
\end{equation}
Choose closed representatives of a basis of $H^{n-2}(D;\C)$.
Equation~\eqref{curv:ambient-annihilation} says that
$i_*([c_E]\cap[L])=0$.  The injection in
\eqref{curv:topological-input} remains injective over $\C$.
Poincar\'e duality on $L$ now gives $[c_E]=0$.
\end{proof}

\begin{proof}[Proof of Proposition~\ref{prop:curvature}]
First consider the cyclic canonical model on
\eqref{curv:canonical-space}.  Suppose its curvature has a first
nonzero coefficient $c_E^{\mathrm{can}}$.  The local finiteness of
the energy filtration ensures that a first coefficient exists if the
curvature is nonzero.  Filtered transfer from this model to the
geometric forms model includes an arity-zero component $F_0$ of
positive valuation.  Set $b=F_0$.  The zero-input functor equation is
\begin{equation}\label{curv:transfer-equation}
 m_0^{\mathrm{forms},b}
       =F_1(m_0^{\mathrm{can}}).
\end{equation}
Its leading term is
$F_{1,0}(c_E^{\mathrm{can}})$, where $F_{1,0}$ induces the ordinary
cohomology identification.  Since the canonical degree-two summand
is $H^2(L;\C)$, this is a nonzero cohomology class.  But
Lemma~\ref{curv:deformation} says that the first geometric curvature
coefficient is exact, a contradiction.

The same argument applies to a canonical structure obtained by
promotion from finite cutoffs.  Choose a geometric cutoff strictly above $E$, and compose transfer
with the cutoff pseudo-isotopy.  Apply Lemma~\ref{curv:stokes}
relative to the ordinary operations on this geometric representative.  Its arity-zero
translation still has positive valuation, and its leading unary map
still induces an isomorphism on ordinary cohomology.  Thus
\eqref{curv:transfer-equation} through that cutoff gives the same
contradiction.  All these comparisons use the energy/class filtration
of the fixed $J_0$ disk system.  Lemma~\ref{curv:energy} supplies the common gap for the ordinary
operations and their comparisons.  The interior functionals only
need to be constructed on the chosen geometric representative; no
identification of them across the pseudo-isotopy is used.  Terms containing positive-valuation cochains require
only bounded arity below a fixed energy.  Equivalently, one may carry
out the argument in filtered $A_{n',k'}$ models, with both cutoffs
larger than those needed to detect $c_E^{\mathrm{can}}$.
The canonical curvature is therefore zero.

For an embedded test copy there are no positive-energy self-disks
at $J_0$, so its geometric curvature is zero.  Keep $J_0$ fixed in
the pseudo-isotopies and choose transfers with zero arity-zero
translation on these objects.  The self-disk contributions, including
the arity-zero homotopy terms, are empty, so promotion preserves this
choice.  On a strictly increasing word all morphisms are isolated
intersection generators.  Energy-zero transfer is the identity on
these point complexes, and its restriction needs no operations with
repeated labels.  The operations on this sector therefore remain the
geometric disk operations up to the stated cutoff pseudo-isotopies.

Homotopy units give cohomological units on the resulting uncurved
objects.  For $L$ this can also be seen by transferring the zero
canonical bounding cochain to the forms model: its image is the
positive cochain $F_0$, which is a bounding cochain by
\eqref{curv:transfer-equation}.  Deform the homotopy units there and
use the deformed filtered quasi-isomorphism.  No strict chain unit
is required for this step.  All operations are graded, so the
cohomological unit has degree zero.

After extending to $\Lambda$, the canonical morphism complexes are
finite-dimensional.  The self complex of $L$ has degree-two chain
space exactly $H^2(L;\C)\otimes\Lambda$, because all branch
generators have odd degree.  Its degree-two cohomology is a
subquotient of that space, proving \eqref{curv:degree-two-bound}.
Finally, the nondegenerate cyclic pairing gives duality with shift
$n$ on morphism cohomology.  Its extension to twisted complexes and
duality on retracts gives the same pairing on the split-closed
triangulated hull, whose Serre functor is consequently $[n]$.
\end{proof}

\section{Comparison with the ordinary factor categories}
\label{sec:comparison}

Proposition~\ref{prop:curvature} places the immersion and the embedded
test branes in a proper uncurved cyclic category.  To use the ordinary
mirror equivalences, we must identify the category generated by the
embedded tests inside that category.  The needed comparison concerns
only words of distinct, ordered copies.  This restriction gives isolated
asymptotic intersections at every end, including every internal end of a
broken polygon, and removes both self-intersection sectors and unit
conventions from the analytic comparison.

Fix one of the finite full cyclic categories constructed in
Proposition~\ref{prop:curvature}, and select its embedded copy labels.
Choose any total ordering of these labels and write the copies in that
order as
\[
 K_a=\Sigma_a\times G_a,\qquad 1\leq a\leq N,
\]
where each is a copy of a test from Proposition~\ref{prop:realization}.
Omit $\Sigma_a$ when $Y$ is a point.  Here $\Sigma_a$ and $G_a$ include the Hamiltonian
perturbations and transported brane structures of the copies.  As in
the preceding construction, all pairs are transverse and there are no
triple intersections, also in each factor separately.  The fixed split
almost complex structure $J_0$ has no nonconstant spheres, and none of
these embedded copies bounds a nonconstant $J_0$-holomorphic disk.
A \emph{strict word} means a sequence of labels
$a_0<\cdots<a_k$.  Its inputs lie in the intersection complexes for
$(K_{a_{i-1}},K_{a_i})$, and its output lies in the complex for
$(K_{a_0},K_{a_k})$.  We initially omit identities and self-morphism
spaces.  Thus the directed category in this section means precisely
these complexes and their operations on strict words.

\begin{proposition}[Directed comparison]
\label{prop:directed-comparison}\label{prop:comparison}
For each fixed finite full cyclic category above and every total
ordering of its embedded copy labels, its strict directed sector is
quasi-isomorphic, after
an extension of its Novikov field if necessary, to the corresponding
ordinary directed Fukaya category of the embedded product copies.
When $Y$ is present, this ordinary directed category is
quasi-isomorphic to the directed restriction of the $A_\infty$ tensor
product of the two ordinary factor categories.  These comparisons use
the specified brane labels and their usual graded intersection
complexes.  Consequently the ordinary mirror functors in the two
factors give a comparison with the directed category on the specified
objects $Q_\ell\boxtimes N_M$.
All the finitely many ordering comparisons for this same full cyclic
category may be made over one field extension.
\end{proposition}

The two comparisons change different parts of the construction.
For the first, we keep $J_0$ and the intersection generators fixed
and change disk integration into regular Floer counts by introducing
Hamiltonian terms.  Section~\ref{comp:virtual} verifies the relative
chart and gluing hypotheses for that change.  It then removes the
energy cutoffs algebraically, while retaining the directed sector
of the originally fixed full cyclic category.

For the second comparison, we retain the regular counts and change
their domain parameters.  The product category uses the same polygon
in both factors; the tensor product uses an associahedral diagonal.
Sections~\ref{comp:domain-chains}--\ref{comp:orientations} construct
and count a homotopy between these choices, including the boundary
and sign identities that make the counts an operad map.
Section~\ref{comp:operadic-completion} then identifies the result with
the tensor product of the ordinary factor categories.  All these
comparisons concern embedded copies, and ordinary mirror symmetry
is used in the individual factors.  The freedom to order the same
copies in every way will allow us to test compositions through three
distinct copies and prove that the copy arrows are invertible in
Lemma~\ref{realization:copies}.

\subsection{Regular Hamiltonian models with fixed intersections}
\label{comp:regular-models}

For each ordered pair choose strip data of the form $H(t,x)\,dt$,
$0\leq t\leq1$, with $H$ small and identically zero on fixed small
balls around the intersections of that pair.  Take $H$ sufficiently
small on the complement as well.  Its chords are then exactly the
original intersections.  On a stable polygon, with at least two
inputs, use a domain-dependent Hamiltonian one-form whose tangential
restriction to every Lagrangian boundary arc is zero.  The almost
complex structure remains $J_0$.  Product data are the sums of the
corresponding data in the two factors.

We require agreement with the pair data on the terminal parts of the
external ends and on long necks.  On polygon cores the data are smooth
up to domain faces and restrict there to the data on the component
polygons.  More precisely, use the usual consistent strip-like ends
and simultaneous plumbing collars on the compactified domain spaces.
If $T$ is a neck length, smoothness of the core coefficients is taken
in a coordinate $\exp(-cT)$, for a fixed $c>0$.  Thus their difference
from the limiting component coefficients, together with every
required length derivative, is exponentially small.

These conditions leave enough freedom for regularity.  To see their
compatibility, first choose data agreeing with their component data
on smaller collars and carrying the prescribed strip model near each
node.  Retain that strip model on terminal end and neck
neighborhoods; variations occur on the complementary core patches.
The restrictions to intersecting domain faces agree, so extend their
coefficients by the ordinary extension theorem for functions on a
manifold with corners, using inclusion--exclusion and a partition of
unity.  The zero tangential boundary condition is linear and is
preserved by this extension.  At later arities one may shrink the
fixed neighborhoods, since the lower-arity data already equal the
strip model there.  This construction gives smooth extensions with
equality on faces; it does not impose equality with component data on
an entire domain collar.  In particular, at any smooth domain, even
one lying in a plumbing collar, there remain open interior core
patches on which the Hamiltonian data may vary.  The radial
plumbing description and strip height used here are also those of
\cite[Proposition~7.5]{A8}.

\begin{lemma}
\label{comp:regularity}
For the finite list above, the Hamiltonian data can be chosen small,
with the stated face and strip conditions, so that the ordinary
strict-word operations are regular.  They can also be chosen regular
over any prescribed countable collection of smooth parameter chains
in products of domain spaces, including all their relative face
interiors and all component configurations at those faces.
Compactness below an area bound has only polygon and strip
degenerations.  All nonconstant strips and all contributing stable
polygons have positive area, with a common positive lower bound for
this finite directed system.
\end{lemma}

\begin{proof}
The energy identity compares analytic energy with ordinary symplectic
area and a curvature term supported on the cores.  The integral of
that curvature is uniformly bounded for each word, because the pure
strip model has zero curvature on every neck.  There is no boundary
term along a labeled arc, and the chords have not moved.  On strips,
analytic energy equals area.  Areas are additive at broken
configurations and are constant in a relative class during the
allowed homotopies of data.

At a small-scale bubble the Hamiltonian terms disappear, leaving a
$J_0$-holomorphic sphere or a disk on one of the embedded copies.
Neither exists.  Compactness therefore reduces to ordinary Floer
compactness with isolated transverse asymptotes.  Below a fixed area
bound, only finitely many relative classes are realized, even before
regularity: near a limiting broken configuration the total relative
class is fixed.  For regular data this gives finite zero-dimensional
counts and compactifications of the one-dimensional moduli spaces.

Nonconstant strips have positive area and a positive energy gap.
Stable polygons also have a uniform positive area gap once the data
are sufficiently small.  Otherwise a sequence of data tending to
zero and polygons with areas tending to a nonpositive value would
have analytic energies tending to zero.  Compactness, including on
long necks, would make every component constant.  A stable vertex in
a strict word has at least three distinct boundary labels, so such
a constant component would give a triple intersection.  This is
excluded.  There are only finitely many strict words, which gives a
single small-data neighborhood and a common gap.

For strips, regularity follows by varying $H(t,x)$.  The regular-point
argument based on unique continuation gives dense injective points
of $(t,u(s,t))$ on every nonconstant trajectory, away from its
asymptotes; it applies with Lagrangian boundary just as in the Floer
strip argument of \cite{FloerHoferSalamon}, with the corrected
linear-algebra step in \cite{FloerHoferSalamonErratum}.
The required variations
may be supported away from the intersection balls.  Indeed a
nonconstant strip cannot remain entirely in those balls: the local
sheets are exact there, so its area would be zero.  A constant strip
has an invertible transverse operator.

For polygons, domain-dependent Hamiltonian variations on interior
core patches give transversality without a simplicity assumption on
the map.  The same argument applies over any fixed smooth parameter
chain.  On a domain face, the vertices in one factor have different
tuples of boundary labels, by strictness.  Their variations can
therefore be chosen independently.  A variation on the relative
interior of a face extends to the full system by the face-extension
construction above; the restrictions on higher families do not
constrain variations on these component core patches.  Thus the
universal linearized problem is surjective on every relevant
stratum.  Sard--Smale, with countable exhaustions by energy, class,
and parameter stratum, supplies simultaneous regularity.  Applied to
the diagonal chains it gives regular product polygon operations.
Product strip regularity is the direct-sum statement for the factor
strip operators.
\end{proof}

\subsection{Relation to the ordinary Fukaya categories}
\label{comp:ordinary}

The preceding data compute the ordinary factor categories used in
mirror symmetry.  For the graphs this is ordinary weakly exact
Floer theory.  For the quartic spheres, take their disk-excluding
almost complex structure to be $J_0$, as permitted in Seidel's
construction \cite[Section~8]{SeidelQuartic}.  It is also sphere-free,
so regularity at these fixed data needs no additional bubble-incidence
condition.  Extend the strict-word data to all ordinary words on the
same copies; repeated labels use nondegenerate chord perturbations.
The extension is again inductive over domain faces, with the same
terminal strip models and smooth core coefficients.

We explain why the slightly more flexible collar condition used
above gives the same ordinary category.  One can put two choices on
separately labeled copies and extend their mixed-word data.  Energy
and compactness for the comparison follow from the same no-bubbling
argument.  At a broken curve the core coefficients differ from their
limits exponentially in the neck lengths, so ordinary gluing with
isolated asymptotes applies.  Continuation, binary products, and
self-Floer units then identify the two labels of each object.  This
is the ordinary comparison argument; the collar proof in
\cite[Proposition~7.5]{A8} uses exactly these hypotheses and even
allows decaying errors in strip-end coefficients.

When comparing with an ordinary choice that has varying almost
complex data, the quartic factor requires the usual low-dimensional
sphere-avoidance argument.  The prescribed disk-excluding structure
at each boundary label is retained, which excludes boundary bubbles.
In real dimension four with $c_1=0$, a simple unmarked sphere has
expected dimension $-2$.  Adding the two incidence marks, one on
the sphere and one on the principal component, and imposing the
four-dimensional matching condition leaves a loss of two dimensions.
Removing a sphere tree does not change the principal component
index.  Hence sphere incidence can be avoided in the zero- and
one-dimensional systems used for operations and comparisons.  For
a multiply covered sphere, its image lies on a simple sphere, so the
same incidence exclusion suffices.

The required transversality uses almost complex variations for the
sphere and Hamiltonian variations for the principal component,
including evaluation at the incidence point.  Constant or contracted
components do not alter this test.  In a relative construction,
previously fixed faces already satisfy the exclusion, and the new
variations are available on the relative interiors.  The terminal
strip and neck data may be chosen sphere-free; strip data depending
on strip height form only a one-parameter family.  Continuation data
vary on compact cores, where the same argument applies.  These are
the sphere-exclusion and comparison mechanisms of
\cite[Section~8]{SeidelQuartic}.  Thus the mixed-copy construction
identifies our regular systems with conventional ordinary systems,
including cohomological units.  The actual Hamiltonian sphere copies
represent the original objects by
\cite[Lemma~8.14]{SeidelQuartic}.

If one uses capped-generator Novikov exponents in the rational-brane
normalization, rather than ordinary polygon areas, the coefficients
differ by fixed input and output scalar factors.  Rescaling the
generator lines after the permitted scalar extension makes these
descriptions agree.  This completes the identification of the
regular factor models.  It remains to compare them with the disk
integration operations that were used to include the immersion.

\subsection{The relative comparison from disk integration}
\label{comp:virtual}

Join the zero Hamiltonian data to the regular data by a small path,
constant near its endpoints.  Keep $J_0$ and every intersection chord
fixed.  We use the disk Kuranishi and CF-perturbation construction
over this path.  Here a CF perturbation is the compatible finite
dimensional perturbation used to integrate along a Kuranishi
correspondence.  The comparison will use its relative extension and
Stokes properties; it does not require cyclic perturbations along
the path.

\begin{lemma}
\label{comp:pseudoisotopy}
For every finite energy and arity cutoff, the strict-word disk
integration operations at the zero end and the ordinary regular
counts at the other end are related by a filtered pseudo-isotopy.
The linear comparison has the identity reduction on the
intersection generators.  The data at either endpoint may be
recovered up to the usual finite-cutoff pseudo-isotopy.
\end{lemma}

\begin{proof}
We verify that the relative disk construction applies to the
Hamiltonian terms just introduced.  The analytic point is that the
unchanged transverse asymptotes and the controlled core coefficients
give uniform gluing estimates at every corner.  With these estimates,
the cited chart and relative perturbation constructions apply with
the path parameter included.  The inputs are the disk charts
from obstruction bundle data, their compatible products on normalized
corners, and the relative CF extension and Stokes theorems used in
the cyclic construction
\cite[Sections~9--10 and~12]{AFOOO2026}.  For the chart and gluing
theorems, we use \cite[Theorem~7.1, Condition~7.11, and
Lemma~7.12]{FOOOChartsI}, the disk-component-wise compatibility
construction of \cite{FOOOChartsII}, and the exponential gluing
estimates in \cite[Assumption~3.12 and Theorems~3.13, 6.4,
8.16--8.17]{FOOOGluing}.  The following verifications explain their
application to this relative problem.

\paragraph{Charts and the path parameter.}
Use stabilized strip and polygon domains, with the path coordinate
included among the parameters.  The Hamiltonian terms on a core are
smooth lower-order terms in the Cauchy--Riemann equation.  Compact-core
obstruction sections, transported and summed as in the disk-chart
construction, make the linearized equation surjective at a fixed
domain and fixed path parameter.  Interior test sections suffice:
the adjoint unique-continuation argument is unchanged by these
smooth lower-order terms.  This also makes the path projection
submersive.  Strip translation is removed by the usual local slice.
When stabilization uses auxiliary domain points constrained to
target slices, suitable immersive points exist: the maps under
consideration have strictly positive area.  Strict words give no
map isotropy, and the no-bubbling argument requires no sphere chart.

\paragraph{Coefficients near broken configurations.}
At a gluing seam the ends lie in the balls where the Hamiltonian
vanishes.  The seam equation therefore coincides with the
unperturbed transverse-end equation.  A nonconstant strip core down
a neck may carry a Hamiltonian term, but that term is translation
invariant and transports smoothly with the core.

There is a coordinate issue when a domain neck contains several
Floer breaking seams.  Its domain length is the total length of the
intervening pieces.  Consequently its exponential plumbing
coordinate is a product of the exponential seam coordinates, times
a smooth positive factor.  The core coefficients retain their
smoothness and exponential derivative estimates after this
contraction.  Equivalently, in stabilized charts keep the auxiliary
marks that locate all the retained cores.

For completeness, use radial variables
\[
 r_l=e^{-\pi s},\qquad r_r=e^{-\pi(T-s)}
\]
on a boundary neck, together with the strip height.  Transport to a
stabilized strip core multiplies each one-sided radial variable by
the intervening radial factors and smooth core factors.  Changes of
boundary-node coordinate have a positive nonzero leading real
coefficient.  An origin shift of order $r_lr_r$, after division by
one radial variable, still contains the opposite radial factor.
The strip coframe thus tends smoothly to the limiting strip
coframe in these coordinates.  These are the radial collar
calculations in \cite[proof of Proposition~7.5]{A8}.
For coordinate changes between map charts one retains the marks
needed for local stability, including marks on strip components.
There is no requirement to express the inverse total neck length
smoothly in the inverse lengths of a contracted chain.

\paragraph{Uniform gluing estimates.}
Choose exponential weights below the spectral gaps of the finitely
many transverse asymptotic operators.  On a joined strip of length
$T$, a smoothed version of
$e^{\delta\min(s,T-s)}$ gives a suitable weight.  Augment the
component operators by kernel observables on compact cores and by
the obstruction variables.  Pregluing errors are exponentially
small in a slightly weaker weight: the maps decay at the ends, the
seam equations agree, and the core coefficient errors are
exponentially small.

To build a parametrix, split the forcing near each seam midpoint in
a fixed-width overlap, solve with the augmented component inverses,
and splice over intervals whose lengths are proportional to $T$.
The component one-sided weights dominate the joined weight on their
respective sides, uniformly on that overlap.  The parametrix error
is the sum of an $O(T^{-1})$ cutoff error and an exponentially small
coefficient error.  There is no matching zero mode at a transverse
intersection.  A Neumann-series correction therefore gives a
uniform augmented inverse; the Fredholm indices agree under
gluing.  The quadratic estimate is uniform as well, since all new
lower-order terms occur on fixed compact cores and the seam
equation near the joined values is unchanged.

The differentiated estimates have the same form.  Core length
derivatives of coefficients are exponentially small.  Derivatives
of middle cutoffs and reparametrizations multiply exponentially
decaying map terms.  Passing to weaker weights accommodates every
fixed finite number of derivatives, for several necks and with
some nodes still broken.  Hence the usual disk-chart smoothness in
inverse-length corner coordinates applies.  These are precisely
the exponential estimates required by the cited FOOO gluing
theorems.  Alternatively, the isolated-end case of
\cite[Lemma~4.11 and Proposition~4.12]{A8} uses the same augmented
surjectivity, spectral gaps, core transports, and exponentially
controlled coefficients.  Its strip splice estimates are estimates
for finite-rank systems and apply in the present dimension.  The
converse gluing assertion follows from closeness on the cores and
decay along the joining ends, giving the local chart coverage.

\paragraph{Compatible charts and relative perturbations.}
At a normalized corner, choose the product of the component
charts over their common path parameter, summing their
core-supported obstruction data.  At every iterated corner the
limiting equation data agree by the face construction.  The
disk-component-wise compatibility procedure now uses the same
stabilizations, nested neighborhoods, transported core sections,
openness of augmented surjectivity, and obstruction-data
enlargements as for unforced disks.  Compactness gives a finite
construction at each cutoff.  The preceding estimates supply the
inverse-function and gluing steps, with the path parameter included.
Thus the smooth lower-order core terms introduce no new condition
in that construction.  There are no forgetful-map or cyclic
compatibility requirements on strict words, and all seam matches
are at points.

Use the usual outer collars and relative CF extension, preserving
the component products and the submersive path projection.  Stokes
then gives the pseudo-isotopy equations with the required face
products.  Endpoint chart choices need only be recovered up to a
finite-cutoff pseudo-isotopy, as in the ordinary relative
construction.  At the regular endpoint, a virtual zero-dimensional
moduli space has no boundary strata and is already an oriented
zero-manifold without isotropy.  Its Kuranishi integral is its
ordinary signed count, also if one uses a chart with redundant
obstruction directions.  Thus the endpoint operations are the
ordinary operations from Lemma~\ref{comp:regularity}.

The positive-energy gap makes all nonidentity corrections to the
linear comparison have positive valuation.  Its reduction is
therefore the identity on the unchanged intersection generators.
\end{proof}

All orientations in this comparison come from the same graded spin
determinant lines, with zero background class.  The comparison of
sign conventions is proved in Section~\ref{comp:orientations} below.
If an ordinary convention reverses arrows, use the opposite
category and the corresponding algebraic duality; no reflection of
holomorphic disks is involved.

We next remove the cutoffs without imposing simultaneous geometric
choices at all orders.  Let $\Lambda$ be the Novikov field in use,
with valuation $v$.  There are finitely many strict words in the
finite list, and every intersection complex is finite dimensional.
Thus the equations for an $A_\infty$ isomorphism between its two
directed structures are a finite polynomial system over $\Lambda$.
Include variables for the inverse of its linear part.  For any
valuation bound, Lemma~\ref{comp:pseudoisotopy}, together with the
finite-cutoff comparisons of the chosen promoted cyclic model, provides an
approximate solution.  Its coordinates have nonnegative valuation,
as do the inverse linear-part coordinates, because that linear
part has invertible reduction.  In particular, the source of
these comparisons is the restriction of the fixed full cyclic category.
Changing chart or perturbation choices for a particular ordering is
accounted for by these endpoint pseudo-isotopies; it does not replace
that ambient category.

If this finite system had no solution over an algebraic closure of
$\Lambda$, the Nullstellensatz would give a fixed polynomial
combination of its equations equal to $1$.  Evaluation at an
approximate solution of sufficiently large valuation contradicts
this identity: all polynomial coefficients have a fixed lower
valuation bound, whereas the equation errors have arbitrarily high
valuation.  A solution therefore exists after field extension.
The argument applies to every ordering of the same finite label set.
There are only finitely many orderings, so all their polynomial
solutions may be chosen over one algebraically closed extension of
the original coefficient field.  This proves the first assertion of
Proposition~\ref{prop:directed-comparison}.  We now retain the
regular ordinary models and compare their product operations.

\subsection{Chains of polygon domains}
\label{comp:domain-chains}

Write $\mathcal K_j$ for the compactified space of disk domains with $j$
inputs and one output, $j\geq2$.  It has dimension $j-2$ and its
faces are products of lower polygon spaces.  We use its usual
smooth face structure from strip-like ends and plumbing.  A pair
of factor polygons has domain parameter in $\mathcal K_j\times \mathcal K_j$.
The product Fukaya operation uses the diagonal copy of $\mathcal K_j$ in
this product.  Replacing this diagonal by a cellular operad
diagonal will give the tensor product.

We need a chain model in which this replacement can be counted
geometrically.  A geometric chain in $\mathcal K_j^2$ is a smooth map
\[
 c:P\longrightarrow \mathcal K_j^2
\]
from a compact oriented smooth manifold with corners and finitely
many faces.  Every relative face interior of $P$, including its
interior, must map into a relative face interior of $\mathcal K_j^2$.
Impose also the following condition on the way target faces are
left.  If a face of $P$ has codimension one in a containing face,
the inward derivative of $c$ is strictly inward at every target
face hypersurface that contains the smaller image interior but
does not contain the larger image interior.  Derivatives here are
taken in the exponential plumbing coordinates.  The condition
allows several target nodes to be smoothed by one inward
parameter.

Take linear combinations of these maps, with the usual orientation,
reparametrization, and disjoint-union relations; the boundary is
the sum over oriented codimension-one faces.  A map that ignores
a direct product factor of positive dimension is declared zero.
This is a subcomplex relation: if the ignored factor has dimension
one, its signed endpoints cancel, and in larger dimensions its
boundary factors are still ignored.  Products, face inclusions,
diagonals, and grafting in both target coordinates preserve the
face condition.  With the identity in arity one adjoined, these
chains form a nonsymmetric differential graded operad, denoted
$\mathcal G$; a chain of dimension $d$ has cohomological degree
$-d$.

\begin{lemma}
\label{comp:chain-acyclic}
In every arity $j\geq2$, the chain complex $\mathcal G(j)$ has the
homology of a point, with the ordinary augmentation in degree zero.
\end{lemma}

\begin{proof}
Choose a smooth homotopy of $\mathcal K_j^2$ that immediately moves every
boundary point into the interior, with strictly inward initial
derivative at all its incident hypersurfaces, and then contracts
within the interior.  Such a homotopy is obtained from an inward
vector field followed by a contraction of a compact subset of the
open cell.  Apply it to a chain and use its prism.

At time zero the strict inward condition gives exactly the
required condition on the new prism faces.  At positive time the
image is in the target interior, so no further inward condition
remains.  The prism identity is therefore valid in $\mathcal G$.
At the constant endpoint every positive-dimensional chain ignores
its whole domain as a product factor and is zero.  A
zero-dimensional chain is reduced to its augmentation times the
chosen interior point.  This gives the asserted contraction.
\end{proof}

We can now choose the domain homotopy, before choosing the Floer
data with which it will be counted.  Let
$\mathcal A_\infty=C_{\mathrm{cell}}(\mathcal K_\bullet)$ be the cellular
dg operad of associahedra, with the same cohomological convention.
Choose a cellular diagonal
\begin{equation}
 \Delta:\mathcal A_\infty\longrightarrow
             \mathcal A_\infty\otimes\mathcal A_\infty
 \label{comp:cellular-diagonal}
\end{equation}
lifting multiplication.  Such a diagonal exists by induction on
arity, or by the free resolution description; see
\cite{MarklShnider}.  No coassociative choice is required.

There are two maps from $\mathcal A_\infty$ to $\mathcal G$.
The first takes the fundamental chain of $\mathcal K_j$ to its actual
diagonal in $\mathcal K_j^2$.  The second applies $\Delta$ and then the
cross product of geometric fundamental chains.  Both preserve
the augmentation.  Lemma~\ref{comp:chain-acyclic} and cofibrancy
of the $A_\infty$ operad imply that they are dg-operad homotopic.
One can construct the homotopy arity by arity, using polynomial
differential forms on an interval as a path object: each
obstruction is an augmentation-zero cycle in an acyclic arity.

Fix these two maps and one such homotopy in $\mathcal G$.
Their coefficients involve only countably many geometric chains.
Take the dg suboperad $\mathcal G_0\subset\mathcal G$ generated
by those chains, and include the representatives and chain
relations used in this construction among the choices below.
The countable family remains countable after adjoining products,
faces, and maps with ignored factors removed.  If a relation requires
additional representatives, add finite witnesses for a countable
basis of the relation space and repeat this closure.  This still
uses only countably many maps.  Choose regular
split Hamiltonian data for this family by
Lemma~\ref{comp:regularity}.  We require counts only on
$\mathcal G_0$; no assertion of its acyclicity is needed.

When a parameter lies in a
target face, count the polygons at all its vertices, matching the
internal ends and summing over their discrete intersection
generators.  Weight a configuration by its total area.  The
fixed-domain Fredholm index is
\[
 |x_0|-\sum_{i=1}^j|x_i|.
\]
Consequently a rigid count over a $d$-dimensional parameter chain
acts in degree $-d$.

\subsection{Compactification, gluing, and the chain identities}
\label{comp:chain-counts}

We show that these parametric counts, with the orientation
normalization specified in the next subsection, define a dg
operad map from $\mathcal G_0$ to the endomorphism operations of the directed
intersection complexes.  This is the geometric assertion needed
for the diagonal replacement.  It consists of compatibility with
the ignored-factor relation, boundary, and grafting.

If a chain ignores a positive-dimensional product factor, a
zero-dimensional count on it would have negative expected
dimension after that factor is removed.  Include the maps with
those factors removed among the regularity requirements.  Their
moduli spaces are empty, so the ignored-factor relation is
respected.

Consider next a one-dimensional parametric moduli space below a
fixed area bound.  Compactness in
Lemma~\ref{comp:regularity} permits two phenomena: approach to a
face of $P$ and breaking of a Floer strip.  Passing to a
codimension-$r$ face of $P$ lowers the virtual dimension by $r$.
This remains true when the face creates several nodes in one or
both polygon factors.  Indeed the fixed-domain indices add under
matching at isolated transverse intersections; the domain
parameters are supplied by $P$.  Each additional nonconstant
quotient strip lowers dimension by one more.  Stratum regularity
therefore excludes a strip breaking over a parameter boundary in
a one-dimensional problem, including a breaking down any newly
formed neck.  The same argument excludes boundary strata in a
zero-dimensional problem.  Repeated labels on strip components
cause no additional fiber-product transversality requirement,
since every match is discrete.

The uniform curvature bounds apply while a chain parameter tends
to a face.  Thus this is compactness for the parametric problem,
not merely compactness at each fixed parameter.  Its limiting
configurations at a target face are exactly the component
polygons and the strip breakings just described.

We must also know that each parameter-face configuration is the
end of the one-dimensional moduli space with its prescribed
orientation.  Use a local collar $(t,y)$ on $P$, with $t\geq0$
increasing inward.  For every newly smoothed neck its plumbing
coordinate has the form
\begin{equation}
 \lambda_i(t,y)=t\,c_i(t,y),\qquad c_i(t,y)>0,
 \label{comp:simultaneous-plumbing}
\end{equation}
by the face condition.  The neck lengths tend to infinity, and
their derivatives in $y$ stay bounded.  At $t=0$, regularity of
the nodal problem includes the $y$ variables.  Its augmented
linear operator is therefore the invertible operator needed for
the rigid parametric configuration.

For fixed small $t>0$, simultaneously preglue all these necks and
apply the weighted inverse construction from
Lemma~\ref{comp:pseudoisotopy}, including the $y$ variables in the
linearization.  Core coefficients converge to their limiting
values.  Differentiating in $y$ introduces only bounded
length derivatives multiplied by core-supported terms or by
terms decaying along the necks.  On the component halves, the
length-dependent transport has the same exponential estimates.
The implicit function theorem gives a unique glued solution for
each sufficiently small $t>0$, and ordinary gluing coverage
identifies all solutions near the nodal configuration with this
end.  Transport of the determinant orientation gives the boundary
orientation, with $t$ increasing inward.

This argument uses oriented-end coverage and the inverse estimates
for $t>0$.  It does not assert that the solution map is smooth at
$t=0$ in exponential plumbing coordinates.  The disk charts used
in Lemma~\ref{comp:pseudoisotopy} have the usual inverse-length
smoothness; these two uses of gluing require different conclusions.

At an edge already internal in a target face, a broken strip can
occur as an output breaking of the child polygon or as an input
breaking of the parent.  These are the same matched
configurations with opposite boundary signs, so they cancel.
External strip breakings give the Hom differential.  Grafting of
chains gives composition by the same discrete matching rule.
Finally, on a product of chains in the individual domain spaces,
the count factors: regularity makes a rigid product configuration
a product of rigid configurations in the two factors.  The
remaining point is that all these identities use a single
consistent sign convention.

\subsection{Determinant lines and signs}
\label{comp:orientations}

We use the graded orientation lines of Floer generators and the
spin determinant capping and gluing theorem, with zero background
class; for the immersed conventions see
\cite[Sections~4--5]{AkahoJoyce}.  Our factor Lagrangian dimensions
are even.  This parity removes the extra signs that would come
from moving the oriented rank contribution of a capped disk.

For a generator $x$, choose an input cap of index $|x|$ and an
output cap of reciprocal index.  Their parities agree in even
dimension.  Let $o_x$ be the input-cap line.  Order a fixed map
operator first, its input caps in order next, and its output cap
last.  The spin orientation of the resulting capped disk
identifies the map determinant symbol with the graded Hom line
\[
 \Hom\bigl(o_{x_1}\otimes\cdots\otimes o_{x_j},o_{x_0}\bigr).
\]
Pair reciprocal cap lines so that the invertible transverse
operator on a constant strip gives the identity of $o_x$.

The determinant gluing comparison at a specified end, with its
length held fixed, is local at that end after these caps are
inserted.  Determinant excision detaches the outer framing and
disk contributions; sphere determinants are complex oriented.
Only the two node caps remain in the comparison.  We record the
permutation that recovers the usual Hom sign.  Let a child map
determinant have parity $b$, let the sum of the preceding parent
input-cap parities be $l$, and let the node cap have parity $q$.
The sum of the child input-cap parities is $b+q$.  Starting with
the parent and child determinants, the external input caps, the
parent output cap, and the reciprocal node pair, reorder so that
each determinant is followed by its own caps.  The permutation
exponent is
\[
 bl+q.
\]
For example, if $r$ is the sum of trailing parent input-cap
parities and $o$ is its output-cap parity, the exponent before
simplification is
\[
 b(l+q+r+o)+q(b+q+r+o)+(r+o)(b+q)=bl+q
 \quad\text{in }\Z/2.
\]
The term depending only on the node, including $q$, is absorbed by
the reciprocal contraction already normalized on the constant
strip.  Fixed-operator gluing therefore gives Hom composition
with Koszul exponent $bl$.  This calculation applies to formal
Fredholm capping and gluing, independently of the existence of a
rigid holomorphic curve.

For direct sums choose the product cap to be the direct sum of the
factor caps, with line $o_x\otimes o_y$.  In the order
$o_x\otimes o_y\otimes o_x^\vee\otimes o_y^\vee$, its contraction
first moves $o_y$ past $o_x^\vee$ and then applies the two factor
contractions.  The crossing contributes $(-1)^{|x||y|}$, exactly
the graded tensor-evaluation sign.  The constant product strip is
the direct sum of the constant factor strips, so its determinant
unit is the tensor of their units and fixes this same reciprocal
normalization.  At a node the factors
$(-1)^{q_x}(-1)^{q_y}=(-1)^{q_x+q_y}$ absorb the product node parity;
the intervening mixed permutations remain precisely the tensor
Hom signs.  Determinant excision and product spin frames give
this direct-sum identification, with no additional rank sign in
our even dimensions.

For a parameter chain $c$ of dimension $d$, order its orientation
before the map determinant and write $f(c)$ for the resulting raw
count.  Composing chains of dimensions $p$ and $q$ introduces
$(-1)^{pq}$: the second parameter passes the first map
determinant.  The same sign appears for cross products.  Define
\begin{equation}
 \gamma(c)=(-1)^{d(d-1)/2}f(c).
 \label{comp:count-normalization}
\end{equation}
The identity
\[
 \frac{(-1)^{(p+q)(p+q-1)/2}(-1)^{pq}}
      {(-1)^{p(p-1)/2}(-1)^{q(q-1)/2}}=1
\]
shows that $\gamma$ respects both composition and cross products.

For the boundary identity let $v$ be the raw strip map with a fixed
translation convention.  An output strip places its map
determinant before the polygon determinant, so its translation
direction passes the $d$ parameter directions.  Write $s\in\{1,-1\}$
for its comparison with outward-normal-first orientation before
that permutation.  At an input the strip translation has the
opposite end direction, and its determinant follows the polygon
determinant.  Moving this translation past both the parameter and
the rigid polygon determinant adds no further parity.  Thus the
oriented boundary equation is
\begin{equation}
 f(\partial c)+s(-1)^d v\circ f(c)
       -s\sum_i f(c)\circ_i v=0,
 \label{comp:raw-boundary}
\end{equation}
where the slot compositions already include their graded Hom
signs.  On product generators $v$ has the tensor differential
rule.

The same calculation verifies internal-edge cancellation.  Put
the internal strip determinant in the same position on both
descriptions of the broken configuration.  Every permutation used
to do this is also the permutation in Hom composition, by the
fixed-operator calculation above.  The two translation directions
are opposite, which gives cancellation.  At a parameter face,
determinant gluing is compared at fixed positive smoothing values.
The parameter orientations then compare by their ordinary
boundary orientation, with $t$ inward in
Equation~\eqref{comp:simultaneous-plumbing}.  There is no separate
normal sign for each node smoothed by that same $t$.

Set the Floer differential to be $s v$.  Since $f(c)$ has degree
$-d$, Equation~\eqref{comp:raw-boundary} gives
\[
 d_{\mathrm{End}}f(c)=(-1)^{d+1}f(\partial c).
\]
The quotient of the normalization signs in dimensions $d$ and
$d-1$ is $(-1)^{d-1}$.  Therefore
\begin{equation}
 d_{\mathrm{End}}\gamma(c)=\gamma(\partial c).
 \label{comp:dg-count}
\end{equation}
Together with the compactification and grafting arguments, this
proves the promised dg operad count identities.

We finally compare this convention with conventional Fukaya
signs.  Match the cap lines and the order of geometric inputs,
using the same dg or suspended translation throughout.  The
remaining difference is removed by strict sign changes at
generators, as in \cite[Lemma~7.7]{A8}.  Here is the argument in the
present setting.  Both conventions use the same spin orientation
of the capped disk, the same complex orientation on sphere
factors, and the same reciprocal contractions.  They consequently
have no relative-class-dependent ratio.  At fixed index the
transverse graded-plane data are connected; frame paths with the
fixed spin data and determinant winding give the same conclusion.
The remaining factors arise from order, marking, and suspension
permutations and depend only on the corner indices.

To make the two-vertex argument explicit, use suspended degrees
$a_i=|x_i|-1$ and write the ratio as
$(-1)^{r_k(a_1,\ldots,a_k)}$, with $r_k$ valued in $\Z/2$.
The output suspended degree of a rigid vertex is
$1+\sum_i a_i$.  The determinant boundary identity says that
\begin{equation}
 \begin{split}
 &r_s(a_i,\ldots,a_{i+s-1})\\
 &\quad+r_{k-s+1}\left(a_1,\ldots,a_{i-1},
     1+\sum_{j=i}^{i+s-1}a_j,a_{i+s},\ldots,a_k\right)
 \end{split}
 \label{comp:sign-cocycle}
\end{equation}
is independent of the inserted block and its length.  Both
conventions have the same suspended insertion sign.  Thus
Equation~\eqref{comp:sign-cocycle} compares the same oriented
one-dimensional formal gluing problem.  It applies to arbitrary
integer cap indices, without assuming that the associated
holomorphic moduli space is nonempty.

Choose $h:\Z\to\Z/2$ recursively so that
$h(a+1)+h(a)=r_1(a)$.  Rescaling the line of suspended degree $a$
by $(-1)^{h(a)}$ replaces $r_k$ by
\[
 r'_k=r_k+h\left(1+\sum_i a_i\right)+\sum_i h(a_i).
\]
It makes $r'_1=0$ and preserves
Equation~\eqref{comp:sign-cocycle}.  Comparing a unary child at
input $i$ with a unary parent now gives
$r'_k(\ldots,a_i+1,\ldots)=r'_k(\ldots,a_i,\ldots)$.
Thus $r'_k$ is a constant $b_k$.  Comparing a binary child with a
unary parent gives $b_k=b_2+b_{k-1}$, so $b_k=(k-1)b_2$.
A common line rescaling by $(-1)^{b_2}$ removes these constants.
This proves a strict sign identification for both the ordinary
counts and the directed disk-integration comparison.  The
integer gradings and even factor dimensions used here are exactly
the features needed in this argument.

\subsection{The operadic diagonal and the product comparison}
\label{comp:operadic-completion}

Apply the count map $\gamma$ on $\mathcal G_0$ to the homotopy
chosen in Section~\ref{comp:domain-chains}.  At the first
endpoint, counts on the actual diagonal are the split product
polygon counts.  At the second endpoint, factorization on cross
products gives precisely the tensor-product operations determined
by $\Delta$.  For strips the two endpoints have the same tensor
differential: a rigid quotient product strip consists of a rigid
quotient strip in one factor and a constant strip in the other;
if both factors move, the independent translation gives excess
dimension.

The construction applies with object labels.  The endomorphism
operad simply has a component for each composable ordered word,
and its compositions use the corresponding slots.  Only the
bounded arities allowed by the finite strict ordering occur.  The
operadic path therefore gives isomorphic directed $A_\infty$
structures: integrate its coderivation homotopy successively in
arity over the characteristic-zero field.  The differential is
fixed and there is no arity-one homotopy term, so every step is a
finite integration of polynomial coefficients, with identity
linear term.

To finish, the cellular-diagonal tensor computes the ordinary
tensor through dg models and respects $A_\infty$
quasi-isomorphisms.  One may first forget units.  Bar--cobar
rectification replaces an $A_\infty$ quasi-isomorphism by a
zigzag of strict quasi-isomorphisms of operad algebras
\cite{LodayVallette}.  For fixed-object categories the same
construction uses composable tensor words over the finite
product of ground fields given by the object idempotents.
Units may be adjoined freely when augmented resolutions are
used.  Tensor these strict diagrams with
Equation~\eqref{comp:cellular-diagonal}; over a field their maps
remain quasi-isomorphisms.  At the dg endpoints only the
associative multiplication survives, and the resulting tensor is
the ordinary dg tensor.  This also explains the agreement with
the tensor-product framework in the product theorem of
\cite{AmorimProduct}.

Restricting these comparisons to the product labels and strict
words preserves every operation and functor equation used here.
The ordinary factor mirror equivalences then give the stated
directed comparison with their dg mirror categories.  This proves
the second assertion and completes
Proposition~\ref{prop:directed-comparison}.

\section{Completion of the realization construction}
\label{sec:realization}

We now prove Proposition~\ref{prop:realization}.  The immersion of
Proposition~\ref{prop:immersion} has been placed in a proper uncurved
cyclic category by Proposition~\ref{prop:curvature}, and
Proposition~\ref{prop:directed-comparison}
identifies the directed operations on every finite collection of
embedded copies with those on the corresponding mirror tests.
We first assemble these finite comparisons.  Cofinal copies will then
recover the full mirror category, to which the immersed object admits
an orthogonal projection.

\subsection{Copies and a common coefficient field}

Retain the ordered Hamiltonian copies chosen in
Section~\ref{sec:curvature}.  Each of the finitely many tests
$K_{\ell,M}$ occurs cofinally: after every label there is a later copy
of each test.
Finite sets in the exhaustion will always include the original tests
and at least three copies of each test.  Denote
the immersed object by $L$, with its grading and spin data understood.

\begin{lemma}[Invertibility of copy arrows]
\label{realization:copies}
In each finite immersed category containing at least three copies of
each test, every nonzero degree-zero cohomology class between distinct
copies of the same test is invertible.
\end{lemma}

\begin{proof}
The strict pair comparison gives a one-dimensional degree-zero
cohomology space between copies.  Given three distinct copies, choose
their directed order in any desired way.
Proposition~\ref{prop:directed-comparison}, ordinary HMS, and
Hamiltonian-copy invariance in each factor show that the product of two
nonzero composable degree-zero classes between those copies is nonzero.

Let $u:A\to B$ and $v:B\to A$ be such classes.  Choose a third copy $C$
and a nonzero class $w:A\to C$.  If $vu=0$, associativity would give
$0=w(vu)=(wv)u$, contrary to the nonzero-product assertion on distinct
copies.  Thus $vu\ne0$, and similarly $uv\ne0$.  The self degree-zero
cohomology has dimension at most one, because its canonical chain space
in that degree is $H^0(K_{\ell,M})$.  The cohomological units therefore
make $vu$ and $uv$ nonzero multiples of the identities.  Rescaling one
arrow gives inverse classes.
\end{proof}

The ordinary mirrors of all copies can be fixed once and for all.
Indeed, first choose the finite ordinary factor categories on the
original tests and their specified images $Q_\ell,N_M$.  For each larger
finite system, compare ordinary choices after including these original
labels.  Hamiltonian copies are isomorphic to them in the ordinary
categories.  The finite comparison can consequently use the same
mirror objects, with repeated labels.  No analytic functor on an
infinite collection is needed.

The finite cyclic categories and their comparisons need not be
compatible as the label set increases.  We assemble them by a field
ultraproduct: the canonical complexes bound each Hom dimension, and
every equation involving fixed labels is a finite matrix identity.
These two facts allow all the required structures to pass to one
coefficient field in the next lemma.

\begin{lemma}[Simultaneous finite comparisons]
\label{realization:field}
Over an extension field $\Bbbk$ of characteristic zero and cardinality
at most the continuum, there is a proper uncurved cyclic
$A_\infty$ category containing $L$ and all the test copies.  Its
directed subcategory on the test copies admits an $A_\infty$ functor to fixed mirror
models, inducing isomorphisms on the cohomology of every strict arrow.
The degree-two bound for $L$ is
\begin{equation}
 \dim_{\Bbbk}\Hom^2(L,L)
 \le b_2(D)-\dim\mathcal Z.
 \label{realization:ambient-bound}
\end{equation}
The category has cohomological units and a nondegenerate cyclic pairing
of dimension $n$.
\end{lemma}

\begin{proof}
Exhaust the labels by finite systems.  For each one, take the uncurved
$A_\infty$ category of Proposition~\ref{prop:curvature}, with its bound
\eqref{curv:degree-two-bound}, and the
directed comparison of Proposition~\ref{prop:directed-comparison},
extending its coefficient field when necessary.  Pass to minimal
models.  The immersed models can be chosen cyclic: a finite complex
with a nondegenerate compatible pairing admits a splitting and a
contraction compatible with that pairing.  Transfer the directed
functors to these models and to the fixed mirror models on the repeated
labels.

For any fixed pair of labels, the dimensions and degrees of the minimal
Hom spaces are bounded by the ranks of the corresponding finite
canonical complexes.  Choose a nonprincipal ultrafilter on the
exhaustion.  For each fixed Hom, its graded ranks are constant on a set
in the ultrafilter.  Use bases of those ranks and take the ordinary
field ultraproduct of the coefficient fields.  Every fixed
$A_\infty$ equation and functor equation is a finite matrix identity,
so passes to this field.  The same applies to units, cyclicity,
nondegeneracy of the pairings, invertibility of the strict-arrow
cohomology maps and of the copy arrows, and
\eqref{realization:ambient-bound}.
The fixed mirror models undergo ordinary scalar extension along the
diagonal embedding of their common ground field.

The finite canonical ranks show that the resulting category is proper.
All fields in the finite constructions can be taken to have cardinality
at most the continuum, and a countable ultraproduct has the same bound.
There is therefore an abstract embedding of the resulting field into
$\C$.  The construction has used ordinary finite-object HMS before
taking this ultraproduct; it makes no assertion of analytic mirror
symmetry over an ultraproduct field.
\end{proof}

\subsection{Recovering full morphisms from directed copies}

We explain why the directed comparison determines the enhanced full
subcategory of copies.  The two receiving categories are the full
subcategory of test copies in the immersed category and the full
subcategory of their specified mirror images, with repeated labels.
Replace both receiving categories by
strictly unital models, transport the functors on strict words, and
adjoin formal units to the directed category.  Write $\mathcal D$ for
this unital directed category, and let $\mathcal C$ denote either full
category of copies receiving it.  The receiving functor is a
quasi-isomorphism on strict arrows, and copy arrows become isomorphisms
in cohomology by Lemma~\ref{realization:copies} and its ordinary mirror
counterpart.

For a label $j$, write $h_j^{\mathcal D}=\Hom_{\mathcal D}(-,j)$ for its
right Yoneda module.  Choose successive later copies
$j_0,j_1,\ldots$ of $j$ and degree-zero copy arrows between them.
Restriction of the full representable admits a natural comparison
\begin{equation}
 \mathop{\operatorname{hocolim}}_k h_{j_k}^{\mathcal D}
 \longrightarrow
 \Hom_{\mathcal C}(-,j)|_{\mathcal D}.
 \label{realization:telescope}
\end{equation}
To construct it, first take the same telescope of restricted full
representables.  All its arrows are equivalences, so its homotopy
colimit is the right-hand side.  The Yoneda comparisons from directed
representables to restricted full ones commute with the arrows up to
the natural homotopies and induce \eqref{realization:telescope}.
At a fixed argument label $i$, every sufficiently late $j_k$ satisfies
$i<j_k$, where the comparison on Hom is a quasi-isomorphism.
Thus \eqref{realization:telescope} is a quasi-isomorphism on every
argument.

Taking Hom from this telescope to another restricted full
representable gives the homotopy limit of its values at $j_k$.
The transition maps in that homotopy limit are equivalences, because
the full-category copy arrows are invertible.  It follows that
restriction of Yoneda is fully faithful on the full copy category.
This argument applies to higher $A_\infty$ functors as well, using
their natural Yoneda comparisons.

On both sides the restricted representable of each original test is
therefore the same telescope in the derived module category of
$\mathcal D$.  The enhanced split-closed triangulated hulls of the two
full categories consequently agree.  By
Lemma~\ref{setup:graph-generation} and sphere generation, their mirror
hull is
\begin{equation}
 \mathcal A\simeq\Perf(B\times E^g).
 \label{realization:mirror-hull}
\end{equation}
Here we use ordinary external-product generation for smooth projective
varieties: external products of generators of the two factors generate
the perfect category of their product.

\subsection{Projection and the Euler pairings}

Let $\mathcal F$ be the split-closed triangulated hull of the immersed
category including $L$.  Its nondegenerate cyclic pairing makes it
proper with Serre functor $[n]$; the pairing extends to twisted complexes
and then to retracts.  The full subcategory $\mathcal A$ in
\eqref{realization:mirror-hull} is saturated
\cite{BondalVanDenBergh}.  The module
$\Hom_{\mathcal F}(-,L)|_{\mathcal A}$ has finite-dimensional total
cohomology at every object, by properness.  Saturation therefore
represents it by an object
$\mathcal E\in\mathcal A$.  The representing map fits into a triangle
\[
 \mathcal E\longrightarrow L\longrightarrow R\longrightarrow
 \mathcal E[1],
 \qquad \Hom^*(\mathcal A,R)=0.
\]
Serre duality in $\mathcal F$ gives
$\Hom^*(R,\mathcal A)=0$ as well.  The connecting map vanishes, and the
triangle splits.  In particular $\mathcal E$ is a summand of $L$, so
\[
 \dim\Ext^2(\mathcal E,\mathcal E)
 \le\dim\Hom^2_{\mathcal F}(L,L)
 \le b_2(D)-\dim\mathcal Z.
\]
After the abstract field embedding, this proves
\eqref{setup:ext-bound} on the complex mirror.

For every test object, the orthogonal summand $R$ has zero Euler
pairing.  The Euler pairing with $L$ is computed by its transverse
Floer complex, hence equals the oriented intersection pairing with
the immersed fundamental class, up to the common sign fixed by the
grading convention.  Equation~\eqref{topo:class} says
$\PD\,i_*[L]=\nu a$ with $\nu>0$.  It follows that all the test Euler
pairings equal the pairings with $a$ times one common nonzero rational
scalar.  This is \eqref{setup:euler} and completes the proof of
Proposition~\ref{prop:realization}.


\bibliographystyle{alpha}
\bibliography{references}
\end{document}
